\documentclass[11pt,a4paper]{report}

% ---------------------------------------------------------------
%  Algebraic Topology of Manifolds II: Characteristic Classes
%  Lecture notes -- course by Penka Georgieva (Sorbonne Universite)
% ---------------------------------------------------------------

\usepackage[utf8]{inputenc}
\usepackage[T1]{fontenc}
% (lmodern unavailable in this TeX installation; using default fonts)

% --- Page layout: a narrower measure, for comfortable reading ---
\usepackage[a4paper,left=4cm,right=4cm,top=3cm,bottom=3.2cm]{geometry}
\linespread{1.10}
\setlength{\parskip}{2pt plus 1pt minus 0.5pt}

\usepackage{amsmath,amssymb,amsthm,mathtools}
\usepackage{tikz}
\usepackage{tikz-cd}
\usetikzlibrary{arrows.meta,decorations.markings,calc,patterns}
\usepackage[colorlinks=true,linkcolor=blue!50!black,citecolor=blue!50!black]{hyperref}
\usepackage{enumitem}
\setlist{itemsep=4pt,topsep=5pt,parsep=0pt}

% --- Breathing room around displayed mathematics ----------------
\setlength{\abovedisplayskip}{13pt plus 3pt minus 3pt}
\setlength{\belowdisplayskip}{13pt plus 3pt minus 3pt}
\setlength{\abovedisplayshortskip}{9pt plus 3pt}
\setlength{\belowdisplayshortskip}{9pt plus 3pt}

% --- Breathing room around theorem environments -----------------
\makeatletter
\def\thm@space@setup{%
  \thm@preskip=1.3\baselineskip
  \thm@postskip=1.1\baselineskip
}
\makeatother

\theoremstyle{plain}
\newtheorem{theorem}{Theorem}[chapter]
\newtheorem{lemma}[theorem]{Lemma}
\newtheorem{proposition}[theorem]{Proposition}
\newtheorem{corollary}[theorem]{Corollary}
\newtheorem{claim}[theorem]{Claim}

\theoremstyle{definition}
\newtheorem{definition}[theorem]{Definition}
\newtheorem{example}[theorem]{Example}
\newtheorem{construction}[theorem]{Construction}
\newtheorem{exercise}[theorem]{Exercise}

\theoremstyle{remark}
\newtheorem{remark}[theorem]{Remark}
\newtheorem{notation}[theorem]{Notation}

% --- The guiding question opening a topic -----------------------
\newenvironment{guiding}
  {\par\medskip\noindent\textbf{Guiding question.}\itshape\ }
  {\par\medskip}

% --- Proof layout: a strategy paragraph, then numbered steps ----
\newenvironment{strategy}
  {\par\smallskip\noindent\emph{Strategy.}\ }
  {\par\smallskip}

\newcommand{\step}[1]{\par\medskip\noindent\emph{Step #1.}\ \ignorespaces}

% --- Macros ------------------------------------------------------
\newcommand{\R}{\mathbb{R}}
\newcommand{\Z}{\mathbb{Z}}
\newcommand{\C}{\mathbb{C}}
\newcommand{\N}{\mathbb{N}}
\newcommand{\Q}{\mathbb{Q}}
\newcommand{\RP}{\mathbb{RP}}
\newcommand{\CP}{\mathbb{CP}}
\newcommand{\eps}{\varepsilon}
\newcommand{\iso}{\cong}
\newcommand{\gaman}{\gamma^1_n}
\DeclareMathOperator{\Hom}{Hom}
\DeclareMathOperator{\rk}{rk}
\DeclareMathOperator{\im}{im}
\DeclareMathOperator{\Sq}{Sq}
\DeclareMathOperator{\Maps}{Maps}
\DeclareMathOperator{\Ext}{Ext}
\DeclareMathOperator{\GL}{GL}
\DeclareMathOperator{\Or}{O}
\DeclareMathOperator{\spann}{span}
\DeclareMathOperator{\proj}{proj}

% --- Lecture headings: the chapter is titled only by its number ---
% \lecture starts a new lecture: it advances the chapter counter (so that
% theorem numbering stays "5.3"), prints "Lecture 5" as the heading, and puts
% the same string in the table of contents.
\newcommand{\lecture}{%
  \refstepcounter{chapter}%
  \chapter*{Lecture \thechapter}%
  \phantomsection
  \addcontentsline{toc}{chapter}{Lecture \thechapter}%
  \markboth{Lecture \thechapter}{}%
}

\title{\textbf{Algebraic Topology of Manifolds II:\\ Characteristic Classes}\\[1em]
\large Lecture notes by Luis Zegarra}
\author{Course by Penka Georgieva\\ Sorbonne Universit\'e}
\date{}

\begin{document}

\maketitle

\noindent\textbf{About these notes.}
These notes follow the twelve lectures of the course. They are written to be
\emph{read}: each topic opens with the question it answers, and every proof
begins with a short paragraph stating the strategy before the argument is
carried out step by step. Results imported from outside the course are marked
as such, with a reference, rather than proved. The standard reference is
Milnor--Stasheff, \emph{Characteristic Classes}.

\vspace{1em}

\tableofcontents

\lecture

\begin{guiding}
What is a ``continuously varying family of vector spaces'' over a space, how do
we prove that such a family is genuinely twisted, and which constructions of
linear algebra survive the passage from single vector spaces to families?
\end{guiding}

\section{Definitions and first examples}

We begin by fixing the stage. A \emph{topological manifold} is a Hausdorff space
$X$ in which every point has a neighbourhood homeomorphic to $\R^n$; it is
\emph{smooth} when the transition functions between charts are smooth. For a
smooth manifold the tangent space $T_xX$ is obtained from paths $\gamma$ through
$x$, with $\gamma(0)=x$, by taking velocities
\[
\frac{d}{dt}\Big|_{t=0}\gamma(t).
\]

The total space of all tangent spaces is the first and most important example of
the structure we are about to define.

\begin{proposition}
For a smooth $n$-manifold $X$, the set $TX=\bigsqcup_{x\in X}T_xX$ is a smooth
$2n$-manifold, the projection $\pi\colon TX\to X$ is smooth, each fibre
$\pi^{-1}(x)=T_xX$ is a real vector space, and each coordinate chart
$U\subset X$ gives a homeomorphism $U\times\R^n\to\pi^{-1}(U)$ which is linear
on each fibre.
\end{proposition}

\begin{definition}[Vector bundle]
A \emph{real vector bundle} $\xi$ over a topological space $B$ consists of

\begin{enumerate}
\item a topological space $E=E(\xi)$, the \emph{total space};
\item a continuous surjection $\pi\colon E\to B$, the \emph{projection};
\item a real vector space structure on every fibre $F_b(\xi):=\pi^{-1}(b)$,
\end{enumerate}

subject to the following \emph{local triviality} condition: every $b\in B$ has a
neighbourhood $U$, an integer $n\ge0$, and a homeomorphism
\[
h\colon U\times\R^n\longrightarrow \pi^{-1}(U)
\]
such that $x\mapsto h(b',x)$ is a vector space isomorphism
$\R^n\to F_{b'}(\xi)$ for every $b'\in U$. Such an $h$ is a \emph{local
trivialization}. If one may take $U=B$ the bundle is \emph{trivial}; if $n$ is
the same for all $b$ we speak of an \emph{$\R^n$-bundle}, or a bundle of
\emph{rank} $n$.
\end{definition}

\begin{example}\label{ex:mobius}
Three examples to keep in mind throughout.

\begin{enumerate}
\item The cylinder $S^1\times\R$ is the trivial line bundle over the circle.

\item The \emph{M\"obius bundle}: take $[0,1]\times\R$ and glue
$(0,t)\sim(1,-t)$. Locally it is a cylinder; globally it is twisted.

\item $TS^2$, an $\R^2$-bundle over $S^2$, which we will prove is non-trivial.
\end{enumerate}

\begin{center}
\begin{tikzpicture}[scale=0.95]
\draw[thick] (0,0) rectangle (2.4,1.4);
\draw[-{Stealth},thick] (0,0.5) -- (0,0.9);
\draw[-{Stealth},thick] (2.4,0.5) -- (2.4,0.9);
\node at (1.2,-0.55) {\small $S^1\times\R$: edges glued in the same direction};
\begin{scope}[xshift=6.4cm]
\draw[thick] (0,0) rectangle (2.4,1.4);
\draw[-{Stealth},thick] (0,0.5) -- (0,0.9);
\draw[-{Stealth},thick] (2.4,0.9) -- (2.4,0.5);
\node at (1.2,-0.55) {\small M\"obius: $(0,t)\sim(1,-t)$};
\end{scope}
\end{tikzpicture}
\end{center}
\end{example}

\begin{definition}
Two vector bundles $\xi,\eta$ over $B$ are \emph{isomorphic}, written
$\xi\iso\eta$, if there is a homeomorphism $f\colon E(\xi)\to E(\eta)$ commuting
with the projections and restricting to a linear isomorphism
$F_b(\xi)\to F_b(\eta)$ on every fibre.
\end{definition}

\begin{definition}
For a smooth submanifold $M\subset\R^n$, the \emph{normal bundle} $\nu$ of $M$
is the subspace of $M\times\R^n$ consisting of the pairs $(x,v)$ with
$v\perp T_xM$.
\end{definition}

\begin{definition}
A \emph{(cross-)section} of $\xi$ is a continuous map $s\colon B\to E(\xi)$ with
$s(b)\in F_b(\xi)$ for every $b$. It is \emph{nowhere zero} if $s(b)\neq0$ for
all $b$.
\end{definition}

Sections are our main probe. Note three things: a section of a tangent bundle is
exactly a vector field; every bundle has a zero section; and a trivial line
bundle has a nowhere zero section, obtained by transporting the constant section
$b\mapsto(b,1)$ through a trivialization.

\section{Sections and non-triviality}

\begin{guiding}
How can one \emph{prove} that a bundle is non-trivial? The strategy that will
serve us for the rest of the course: exhibit something that a trivial bundle
must possess, and show the bundle at hand does not possess it.
\end{guiding}

\begin{example}[The canonical line bundle]
Write $\RP^n=\{\{\pm x\}\mid x\in S^n\}$ and let $q\colon S^n\to\RP^n$ be the
quotient map. The \emph{canonical line bundle} $\gaman\to\RP^n$ is the subspace
of $\RP^n\times\R^{n+1}$ consisting of pairs $(\{\pm x\},v)$ with $v$ on the
line through $x$ and $-x$, that is, $v$ a multiple of $x$.
\end{example}

\begin{theorem}\label{thm:gamma-nontrivial}
For every $n\ge1$ the bundle $\gaman$ is not trivial.
\end{theorem}

\begin{proof}
\begin{strategy}
A trivial line bundle has a nowhere zero section, so it is enough to prove that
$\gaman$ has none. Pull a hypothetical section back to the sphere, where the
bundle becomes trivial: the section becomes an honest real-valued function,
antisymmetric under the antipodal map, and the intermediate value theorem
produces a zero.
\end{strategy}

\step{1} \emph{The pull-back of $\gaman$ to $S^n$ is trivial.} Consider
\[
q^*\gaman
=\big\{(x,v)\in S^n\times\R^{n+1}\ \big|\ v\in\R x\big\},
\]
\[
\Phi\colon q^*\gaman\longrightarrow S^n\times\R,
\qquad (x,tx)\longmapsto (x,t).
\]
The map $\Phi$ is a bijection, linear on fibres, and both it and its inverse
$(x,t)\mapsto(x,tx)$ are continuous, the former because $t=x\cdot v$ is a
polynomial in the coordinates.

\step{2} \emph{A section becomes a function.} Suppose $s\colon\RP^n\to E(\gaman)$
is a section. Then $x\mapsto s(q(x))$ is a section of $q^*\gaman$, so by Step 1
there is a continuous $t\colon S^n\to\R$ with
\[
s\big(q(x)\big)=\big(\{\pm x\},\ t(x)\,x\big)
\qquad\text{for all }x\in S^n.
\]

\step{3} \emph{The function is odd.} Since $q(-x)=q(x)$, evaluating the formula
of Step 2 at $-x$ gives $t(-x)\cdot(-x)=t(x)\cdot x$, whence
\[
t(-x)=-t(x).
\]

\step{4} \emph{Conclusion.} Fix $x\in S^n$ and let $\alpha$ be a path in $S^n$
from $x$ to $-x$, which exists because $S^n$ is path connected for $n\ge1$. The
continuous function $t\circ\alpha$ takes the values $t(x)$ and $-t(x)$ at its
endpoints. If $t(x)=0$ the section already vanishes at $q(x)$; otherwise the
two values have opposite signs and the intermediate value theorem gives
$x_0$ on the path with $t(x_0)=0$, so $s(q(x_0))=0$. In either case $s$ is not
nowhere zero, and $\gaman$ is not trivial.
\end{proof}

\begin{remark}
For $n=1$ the total space $E(\gamma^1_1)$ is precisely the M\"obius band of
Example~\ref{ex:mobius}: the theorem recovers the familiar fact that the
M\"obius band is twisted.
\end{remark}

The next lemma is the workhorse that lets us verify isomorphisms fibre by fibre,
without ever checking continuity of an inverse by hand.

\begin{lemma}\label{lem:fibrewise-iso}
Let $\xi,\eta$ be vector bundles over $B$ and let $f\colon E(\xi)\to E(\eta)$ be
continuous, commuting with the projections, and restricting to a linear
isomorphism $F_b(\xi)\to F_b(\eta)$ for every $b$. Then $f$ is a homeomorphism
and $\xi\iso\eta$.
\end{lemma}

\begin{proof}
\begin{strategy}
Bijectivity is immediate; the whole content is continuity of $f^{-1}$, which is
a local question. In local trivializations $f$ becomes multiplication by an
invertible matrix depending continuously on the base point, and Cramer's rule
inverts it continuously.
\end{strategy}

\step{1} \emph{$f$ is a bijection.} It is injective and surjective on each
fibre by hypothesis, and it preserves fibres, so it is a bijection on total
spaces.

\step{2} \emph{Reduction to a local statement.} Continuity of $f^{-1}$ may be
checked on the members of an open cover of $E(\eta)$. Fix $b\in B$ and choose a
neighbourhood $U$ over which both $\xi$ and $\eta$ are trivial, with
trivializations
\[
g\colon U\times\R^n\to\pi_\xi^{-1}(U),
\qquad
h\colon U\times\R^n\to\pi_\eta^{-1}(U).
\]
It suffices to prove that $f^{-1}$ is continuous on $\pi_\eta^{-1}(U)$.

\step{3} \emph{$f$ in coordinates.} The composite
$h^{-1}\circ f\circ g\colon U\times\R^n\to U\times\R^n$ preserves the first
coordinate and is linear in the second, so it has the form
\[
(b',x)\longmapsto\big(b',\,M(b')\,x\big),
\qquad M(b')\in\GL_n(\R).
\]
Each column $M(b')e_i$ is the second component of
$h^{-1}fg(b',e_i)$, hence depends continuously on $b'$; so all entries of
$M(b')$ are continuous functions of $b'$.

\step{4} \emph{Inverting continuously.} By Cramer's rule,
\[
M(b')^{-1}=\frac{1}{\det M(b')}\,\mathrm{adj}\,M(b'),
\]
whose entries are rational functions of the entries of $M(b')$ with non-vanishing
denominator. Hence $b'\mapsto M(b')^{-1}$ is continuous, and
$(b',y)\mapsto(b',M(b')^{-1}y)$ is continuous.

\step{5} \emph{Conclusion.} On $\pi_\eta^{-1}(U)$ we have
$f^{-1}=g\circ\big(h^{-1}fg\big)^{-1}\circ h^{-1}$, a composite of continuous
maps. Therefore $f^{-1}$ is continuous, $f$ is a homeomorphism, and it is a
fibrewise linear isomorphism by hypothesis.
\end{proof}

\begin{definition}
Sections $s_1,\dots,s_n$ of $\xi$ are \emph{nowhere dependent} if the vectors
$s_1(b),\dots,s_n(b)$ are linearly independent for every $b\in B$.
\end{definition}

\begin{theorem}\label{thm:trivial-iff-sections}
An $\R^n$-bundle $\xi$ is trivial if and only if it admits $n$ nowhere dependent
sections.
\end{theorem}

\begin{proof}
\begin{strategy}
In both directions the map relating $B\times\R^n$ and $E$ is the obvious one;
Lemma~\ref{lem:fibrewise-iso} does the work of checking it is an isomorphism.
\end{strategy}

\step{1} \emph{From sections to a trivialization.} Let $s_1,\dots,s_n$ be
nowhere dependent and define
\[
f\colon B\times\R^n\longrightarrow E(\xi),
\qquad
f(b,x)=x_1s_1(b)+\dots+x_ns_n(b).
\]
This is continuous, since addition and scalar multiplication are continuous on
$E$ and each $s_i$ is continuous; it commutes with the projections; and on the
fibre over $b$ it sends the standard basis to the basis $s_1(b),\dots,s_n(b)$ of
$F_b(\xi)$, hence is a linear isomorphism there. By
Lemma~\ref{lem:fibrewise-iso} it is an isomorphism of bundles, so $\xi$ is
trivial.

\step{2} \emph{From a trivialization to sections.} If
$h\colon B\times\R^n\to E(\xi)$ is a trivialization, set $s_i(b)=h(b,e_i)$. Each
$s_i$ is continuous, and for every $b$ the vectors $s_1(b),\dots,s_n(b)$ are the
image of a basis under a linear isomorphism, hence linearly independent.
\end{proof}

\begin{corollary}
$TS^2$ is non-trivial.
\end{corollary}

\begin{proof}
By Theorem~\ref{thm:trivial-iff-sections} a trivialization would provide two
nowhere dependent vector fields on $S^2$, in particular one nowhere zero vector
field, contradicting the hairy ball theorem.
\end{proof}

\section{Constructions}

\begin{guiding}
Which operations of linear algebra --- restriction, pull-back, products, direct
sums, subspaces, orthogonal complements --- make sense for whole families of
vector spaces? Essentially all of them, and this flexibility is what makes the
theory usable.
\end{guiding}

\begin{construction}[Restriction]
If $\xi$ is a bundle over $B$ and $\bar B\subset B$, then $\xi|_{\bar B}$, with
total space $\pi^{-1}(\bar B)$, is a vector bundle over $\bar B$: local
trivializations restrict. For smooth $\bar B\subset B$ this produces
$TB|_{\bar B}$.
\end{construction}

\begin{construction}[Pull-back]\label{con:pullback}
Given $\xi$ over $B$ and a continuous map $f\colon B_1\to B$, the
\emph{induced}, or \emph{pull-back}, bundle $f^*\xi$ over $B_1$ has total space
\[
E_1(f^*\xi)=\big\{(b,e)\in B_1\times E(\xi)\ \big|\ f(b)=\pi(e)\big\},
\]
topologized as a subspace, with projection $(b,e)\mapsto b$ and the vector space
structure of $F_{f(b)}(\xi)$ on the fibre over $b$. It fits into the square
\[
\begin{tikzcd}[row sep=large, column sep=large]
E_1(f^*\xi) \arrow[r,"\hat f"] \arrow[d] & E(\xi) \arrow[d,"\pi"] \\
B_1 \arrow[r,"f"] & B
\end{tikzcd}
\qquad
\hat f(b,e)=e .
\]
\end{construction}

\begin{proposition}
The pull-back $f^*\xi$ is a vector bundle.
\end{proposition}

\begin{proof}
\begin{strategy}
Only local triviality needs an argument, and a trivialization of $\xi$ near
$f(b)$ can simply be pulled back over the preimage.
\end{strategy}

\step{1} Fix $b\in B_1$ and choose a local trivialization
$h\colon U\times\R^n\to\pi^{-1}(U)$ of $\xi$ with $f(b)\in U$. Put
$U_1=f^{-1}(U)$, an open neighbourhood of $b$.

\step{2} Define
\[
h_1\colon U_1\times\R^n\longrightarrow \pi_1^{-1}(U_1),
\qquad
h_1(b_1,x)=\big(b_1,\ h(f(b_1),x)\big).
\]
This lands in $E_1$ because $\pi(h(f(b_1),x))=f(b_1)$, and it is continuous
because $f$ and $h$ are.

\step{3} It is bijective with continuous inverse
$(b_1,e)\mapsto(b_1,\ \mathrm{pr}_2\,h^{-1}(e))$, and for fixed $b_1$ it is the
linear isomorphism $x\mapsto h(f(b_1),x)$ onto $F_{f(b_1)}(\xi)$. Hence $h_1$ is
a local trivialization.
\end{proof}

\begin{definition}
A \emph{bundle map} from $\eta$ to $\xi$ is a continuous map
$g\colon E(\eta)\to E(\xi)$ carrying each fibre $F_b(\eta)$ isomorphically onto
a fibre $F_{b'}(\xi)$. It induces a map $\bar g\colon B(\eta)\to B(\xi)$,
$b\mapsto b'$, on base spaces, which is continuous.
\end{definition}

\begin{lemma}\label{lem:bundle-map-pullback}
If $g\colon E(\eta)\to E(\xi)$ is a bundle map with base map $\bar g$, then
$\eta\iso\bar g^{\,*}\xi$.
\end{lemma}

\begin{proof}
\begin{strategy}
There is only one reasonable map to write down; check its fibrewise behaviour
and appeal to Lemma~\ref{lem:fibrewise-iso}.
\end{strategy}

\step{1} Define
\[
\varphi\colon E(\eta)\longrightarrow E\big(\bar g^{\,*}\xi\big),
\qquad
\varphi(e)=\big(\pi_\eta(e),\,g(e)\big).
\]
This is well defined: by definition of $\bar g$ we have
$\pi_\xi(g(e))=\bar g(\pi_\eta(e))$, which is exactly the condition defining the
pull-back. It is continuous because $\pi_\eta$ and $g$ are, and it commutes with
the projections to $B(\eta)$.

\step{2} On the fibre over $b$, the map $\varphi$ is
$e\mapsto(b,g(e))$, and $g$ restricts to an isomorphism
$F_b(\eta)\to F_{\bar g(b)}(\xi)$, which is by definition the fibre of
$\bar g^{\,*}\xi$ over $b$.

\step{3} By Lemma~\ref{lem:fibrewise-iso}, $\varphi$ is an isomorphism of
bundles.
\end{proof}

\begin{construction}[Products and Whitney sums]
For $\xi_1,\xi_2$ over $B_1,B_2$ the \emph{Cartesian product}
$\xi_1\times\xi_2$ over $B_1\times B_2$ has total space
$E(\xi_1)\times E(\xi_2)$, projection $\pi_1\times\pi_2$ and fibres
$F_{b_1}(\xi_1)\times F_{b_2}(\xi_2)$; products of local trivializations
trivialize it. When $B_1=B_2=B$, pulling back along the diagonal
$d\colon B\to B\times B$, $b\mapsto(b,b)$, gives the \emph{Whitney sum}
\[
\xi_1\oplus\xi_2:=d^*(\xi_1\times\xi_2),
\qquad
F_b(\xi_1\oplus\xi_2)\iso F_b(\xi_1)\oplus F_b(\xi_2).
\]
\end{construction}

\begin{definition}
Let $\xi,\eta$ be bundles over $B$ with $E(\xi)\subset E(\eta)$. We call $\xi$ a
\emph{sub-bundle} of $\eta$ if every $F_b(\xi)$ is a vector subspace of
$F_b(\eta)$.
\end{definition}

\begin{lemma}\label{lem:internal-sum}
If $\xi_1,\xi_2$ are sub-bundles of $\eta$ such that
$F_b(\eta)=F_b(\xi_1)\oplus F_b(\xi_2)$ for every $b$, then
$\eta\iso\xi_1\oplus\xi_2$.
\end{lemma}

\begin{proof}
The map $E(\xi_1\oplus\xi_2)\to E(\eta)$ sending $(b,e_1,e_2)\mapsto e_1+e_2$ is
continuous, commutes with the projections, and on the fibre over $b$ is the
canonical isomorphism $F_b(\xi_1)\oplus F_b(\xi_2)\to F_b(\eta)$ provided by the
hypothesis. Apply Lemma~\ref{lem:fibrewise-iso}.
\end{proof}

\section{Euclidean vector bundles}

\begin{definition}
A \emph{Euclidean vector bundle} is a real vector bundle $\xi$ together with a
continuous function $\mu\colon E(\xi)\to\R$ which is quadratic and positive
definite on every fibre. Polarization turns $\mu$ into an inner product on each
fibre, varying continuously with the base point; we call it a \emph{Euclidean
metric} on $\xi$, or a \emph{Riemannian metric} when $\xi=TM$.
\end{definition}

The space $\R^n$ carries its standard metric, which restricts to any
$M\subset\R^n$. On a trivial bundle with a Euclidean metric, Gram--Schmidt --- a
continuous process --- turns the standard frame into $n$ orthonormal nowhere
vanishing sections.

\begin{definition}
Let $\eta$ be a Euclidean bundle of rank $n$ and $\xi\subset\eta$ a sub-bundle
of rank $k$. The \emph{orthogonal complement} $\xi^\perp\subset\eta$ is the
union over $b$ of the subspaces $F_b(\xi)^\perp\subset F_b(\eta)$.
\end{definition}

\begin{theorem}\label{thm:orthogonal-complement}
$\xi^\perp$ is a vector bundle and $\xi\oplus\xi^\perp\iso\eta$.
\end{theorem}

\begin{proof}
\begin{strategy}
Only local triviality of $\xi^\perp$ is at issue, since the second assertion
then follows from Lemma~\ref{lem:internal-sum}. Near a point, build an
orthonormal frame of $\eta$ whose first $k$ vectors span $\xi$; the remaining
$n-k$ vectors then trivialize $\xi^\perp$. Producing such a frame uses two
standard facts: linear independence is an open condition, and Gram--Schmidt is
continuous.
\end{strategy}

\step{1} \emph{Setting up.} Fix $b\in B$ and choose a neighbourhood $U$ over
which both $\xi$ and $\eta$ are trivial. Using a trivialization of $\xi|_U$ and
Gram--Schmidt, choose sections $s_1,\dots,s_k$ of $\xi|_U$ which are orthonormal
at every point of $U$.

\step{2} \emph{Extending the frame at one point.} The vectors
$s_1(b),\dots,s_k(b)$ are linearly independent in $F_b(\eta)$, so we may complete
them to a basis
\[
s_1(b),\dots,s_k(b),\,v_{k+1},\dots,v_n
\]
of $F_b(\eta)$. Using a trivialization of $\eta|_U$, extend the vectors $v_j$ to
sections $s'_{k+1},\dots,s'_n$ of $\eta|_U$ that are constant in that
trivialization.

\step{3} \emph{Independence persists.} In the trivialization of $\eta|_U$ the
determinant of the $n\times n$ matrix with columns
$s_1,\dots,s_k,s'_{k+1},\dots,s'_n$ is a continuous function of the base point
which is non-zero at $b$. Hence there is a neighbourhood $V\ni b$, $V\subset U$,
on which these $n$ sections are nowhere dependent.

\step{4} \emph{Orthonormalizing.} Apply Gram--Schmidt pointwise on $V$ to the
frame of Step 3. Since the first $k$ vectors are already orthonormal, the
process leaves them unchanged and produces sections
$s_1,\dots,s_k,s_{k+1},\dots,s_n$ which form an orthonormal frame of $\eta|_V$
at every point, depend continuously on the base point, and satisfy
$\spann(s_1,\dots,s_k)=F_{b'}(\xi)$ for $b'\in V$. Consequently
$s_{k+1}(b'),\dots,s_n(b')$ is an orthonormal basis of $F_{b'}(\xi)^\perp$.

\step{5} \emph{Local triviality of $\xi^\perp$.} Define
\[
V\times\R^{n-k}\longrightarrow \xi^\perp\big|_V,
\qquad
(b',x)\longmapsto x_1\,s_{k+1}(b')+\dots+x_{n-k}\,s_n(b').
\]
It is continuous, fibrewise a linear isomorphism onto $F_{b'}(\xi)^\perp$ by
Step 4, and its inverse is $e\mapsto(b',(\langle e,s_{k+j}(b')\rangle)_j)$,
continuous because the metric is. Hence $\xi^\perp$ is a vector bundle of rank
$n-k$.

\step{6} \emph{The sum.} By construction
$F_b(\eta)=F_b(\xi)\oplus F_b(\xi^\perp)$ for every $b$, so
Lemma~\ref{lem:internal-sum} gives $\xi\oplus\xi^\perp\iso\eta$.
\end{proof}

\begin{example}\label{ex:normal-bundle}
If $M\subset N$ are smooth manifolds and $N$ is Riemannian, then $TM$ is a
sub-bundle of $TN|_M$ and the \emph{normal bundle of $M$ in $N$} is
$\nu=(TM)^\perp$, so that
\[
TN|_M\ \iso\ TM\oplus\nu .
\]

\begin{center}
\begin{tikzpicture}[scale=0.95]
\draw[thick] (0,0) .. controls (1.2,0.9) and (2.8,-0.6) .. (4,0.3);
\node at (4.4,0.3) {$M$};
\fill (1.9,0.12) circle (1.5pt);
\draw[-{Stealth}] (1.9,0.12) -- (2.85,-0.14);
\node[below] at (2.9,-0.08) {\small $T_xM$};
\draw[-{Stealth}] (1.9,0.12) -- (2.14,1.0);
\node[right] at (2.12,1.0) {\small $\nu_x$};
\end{tikzpicture}
\end{center}
\end{example}

\begin{definition}
A smooth map $f\colon M\to N$ is an \emph{immersion} if
$Df_x\colon T_xM\to T_{f(x)}N$ is injective for every $x$. The figure eight is
the image of an immersion of $S^1$ in $\R^2$ which is not an embedding.
\end{definition}

\begin{lemma}\label{lem:normal-immersion}
For an immersion $f\colon M\to N$ with $N$ Riemannian there is a Whitney sum
decomposition
\[
f^*TN\ \iso\ TM\oplus\nu_f,
\]
where $\nu_f$ is called the \emph{normal bundle of the immersion} $f$.
\end{lemma}

\begin{proof}
The differential $Df$ realizes $TM$ as a sub-bundle of $f^*TN$: it is fibrewise
injective by the definition of an immersion, and it is continuous, so its image
is a sub-bundle. Pulling back the Riemannian metric of $N$ makes $f^*TN$
Euclidean, and Theorem~\ref{thm:orthogonal-complement} applied to
$TM\subset f^*TN$ gives $f^*TN\iso TM\oplus(TM)^\perp$; set $\nu_f=(TM)^\perp$.
\end{proof}

This decomposition is the engine behind many computations to come: it converts
geometric information --- how $M$ sits inside $N$ --- into bundle-theoretic
information. In the next lecture we manufacture new bundles from old ones
functorially, and meet the first characteristic classes.

\lecture

\begin{guiding}
Linear algebra has functors: duals, homomorphisms, tensor products. Do they act
on vector bundles? And can we attach to every bundle algebraic invariants ---
cohomology classes --- that detect twisting?
\end{guiding}

\section{Continuous functors of vector bundles}

\begin{definition}
Let $\Theta$ be the category whose objects are finite dimensional real vector
spaces and whose morphisms are isomorphisms. A functor
\[
T\colon \Theta\times\dots\times\Theta\longrightarrow\Theta
\]
of $k$ variables is \emph{continuous} if $T(f_1,\dots,f_k)$ depends continuously
on $f_1,\dots,f_k$, the latter being regarded as matrices.
\end{definition}

\begin{theorem}\label{thm:functor-bundle}
Let $T$ be a continuous functor of $k$ variables and let $\xi_1,\dots,\xi_k$ be
vector bundles over $B$. Then there is a vector bundle
$T(\xi_1,\dots,\xi_k)$ over $B$ whose fibre over $b$ is
$T\big(F_b(\xi_1),\dots,F_b(\xi_k)\big)$.
\end{theorem}

\begin{proof}
\begin{strategy}
The set and the fibres are prescribed; what has to be produced is a topology.
Local trivializations of the $\xi_i$ give candidate charts, and the only thing
to check is that the change of chart is continuous --- which is exactly the
continuity hypothesis on $T$.
\end{strategy}

\step{1} \emph{The set.} Put
\[
E=\bigsqcup_{b\in B}T\big(F_b(\xi_1),\dots,F_b(\xi_k)\big),
\]
with the obvious projection $\pi$ to $B$ and the vector space structure of
$T(\dots)$ on each fibre.

\step{2} \emph{Candidate charts.} Suppose $U\subset B$ trivializes all the
$\xi_i$ simultaneously, with trivializations giving, for $b\in U$, isomorphisms
$h_{ib}\colon\R^{n_i}\to F_b(\xi_i)$ depending continuously on $b$. Since $T$
is a functor, applying it to these isomorphisms gives a linear isomorphism
\[
T(h_{1b},\dots,h_{kb})\colon
T(\R^{n_1},\dots,\R^{n_k})\longrightarrow
T\big(F_b(\xi_1),\dots,F_b(\xi_k)\big),
\]
and hence a bijection
\[
H_U\colon U\times T(\R^{n_1},\dots,\R^{n_k})\longrightarrow\pi^{-1}(U),
\qquad
(b,x)\longmapsto T(h_{1b},\dots,h_{kb})(x).
\]

\step{3} \emph{Compatibility of charts.} Let $U_1,U_2$ be two such open sets.
On $U_1\cap U_2$ the composite $H_{U_1}^{-1}\circ H_{U_2}$ is
\[
(b,x)\longmapsto\Big(b,\ T\big(h^{(1)}_{1b}{}^{-1}h^{(2)}_{1b},\ \dots,\
h^{(1)}_{kb}{}^{-1}h^{(2)}_{kb}\big)(x)\Big),
\]
using functoriality. Each $h^{(1)}_{ib}{}^{-1}h^{(2)}_{ib}$ is the transition
isomorphism of $\xi_i$, which depends continuously on $b$; by continuity of $T$
the whole expression depends continuously on $(b,x)$. The same applies to the
inverse.

\step{4} \emph{The topology.} Declare a subset of $E$ open when its preimage
under every $H_U$ is open. By Step 3 the charts $H_U$ are then homeomorphisms
onto their images, and they are fibrewise linear isomorphisms by construction.
Hence $E$ is a vector bundle with the prescribed fibres.
\end{proof}

\begin{example}
The functor $V\mapsto\Hom(V,\R)$ yields the \emph{dual bundle} $\xi^*$; two
variables give $\xi\otimes\eta$, $\Hom(\xi,\eta)$, $\xi\oplus\eta$, and so on.
\end{example}

\begin{remark}
If $\xi$ is Euclidean, the metric gives an isomorphism $\xi\iso\xi^*$, namely
$(b,v)\mapsto(b,\langle v,\cdot\rangle_b)$; it is continuous and fibrewise an
isomorphism, so Lemma~\ref{lem:fibrewise-iso} applies. Over a paracompact base a
Euclidean metric always exists, by patching local metrics with a partition of
unity.
\end{remark}

\section{The Stiefel--Whitney axioms}

We now \emph{postulate} a theory of invariants and derive consequences; the
construction of the classes will occupy Lectures 4--6. This shortcut lets us
prove interesting theorems immediately, which is also how the subject developed
historically.

\begin{definition}[Axioms]\label{def:sw-axioms}
To each $\R^n$-bundle $\xi$ over a base $B$ there are associated cohomology
classes
\[
w_i(\xi)\in H^i(B;\Z_2),\qquad i\ge0,
\]
the \emph{Stiefel--Whitney classes}, subject to:

\begin{enumerate}
\item[(A1)] \textbf{Rank.} $w_0(\xi)=1$ and $w_i(\xi)=0$ for $i>n=\rk\xi$.

\item[(A2)] \textbf{Naturality.} If $\hat f\colon\xi\to\eta$ is a bundle map
with base map $f$, then $w_i(\xi)=f^*w_i(\eta)$; equivalently
$w_i(f^*\eta)=f^*w_i(\eta)$.

\item[(A3)] \textbf{Whitney product formula.} If $\xi,\eta$ have the same base,
\[
w_k(\xi\oplus\eta)=\sum_{i=0}^{k}w_i(\xi)\smile w_{k-i}(\eta).
\]

\item[(A4)] \textbf{Normalization.} $w_1(\gamma^1_1)\neq0$ for the canonical
line bundle $\gamma^1_1\to\RP^1$.
\end{enumerate}
\end{definition}

\begin{proposition}
If $\xi\iso\eta$ then $w_i(\xi)=w_i(\eta)$ for all $i$.
\end{proposition}

\begin{proof}
An isomorphism is a bundle map whose base map is the identity, so (A2) gives
$w_i(\xi)=\mathrm{id}^*w_i(\eta)=w_i(\eta)$.
\end{proof}

\begin{proposition}\label{prop:trivial-vanishes}
If $\xi$ is trivial then $w_i(\xi)=0$ for every $i>0$.
\end{proposition}

\begin{proof}
\begin{strategy}
A trivial bundle is pulled back from a bundle over a point, and a point has no
cohomology in positive degrees.
\end{strategy}

\step{1} Write $\xi\iso B\times\R^n$ and let $\eps'$ denote the bundle
$\mathrm{pt}\times\R^n$ over a point. The map
$B\times\R^n\to\mathrm{pt}\times\R^n$, $(b,x)\mapsto(\mathrm{pt},x)$, is a
bundle map whose base map is the constant map $c\colon B\to\mathrm{pt}$.

\step{2} By (A2), $w_i(\xi)=c^*w_i(\eps')$. Since
$H^i(\mathrm{pt};\Z_2)=0$ for $i>0$, we have $w_i(\eps')=0$ for $i>0$, hence
$w_i(\xi)=0$.
\end{proof}

\begin{proposition}[Stability]\label{prop:stability}
If $\eps$ is a trivial bundle then $w_i(\eps\oplus\eta)=w_i(\eta)$.
\end{proposition}

\begin{proof}
By Proposition~\ref{prop:trivial-vanishes}, $w_0(\eps)=1$ and $w_j(\eps)=0$ for
$j>0$; the Whitney formula (A3) then reduces to the single term
$w_0(\eps)\smile w_i(\eta)=w_i(\eta)$.
\end{proof}

\begin{proposition}\label{prop:sections-kill-top}
If a Euclidean $\R^n$-bundle $\xi$ admits $k$ nowhere dependent sections, then
\[
w_{n}(\xi)=w_{n-1}(\xi)=\dots=w_{n-k+1}(\xi)=0 .
\]
\end{proposition}

\begin{proof}
\begin{strategy}
The sections span a trivial sub-bundle; splitting it off drops the rank by $k$,
and the rank axiom does the rest.
\end{strategy}

\step{1} Let $s_1,\dots,s_k$ be the sections and let $\eps^k\subset\xi$ be the
sub-bundle whose fibre over $b$ is $\spann(s_1(b),\dots,s_k(b))$. It is a
sub-bundle of rank $k$, trivialized by
$(b,x)\mapsto\sum x_i s_i(b)$ as in Theorem~\ref{thm:trivial-iff-sections}.

\step{2} Let $\eta=(\eps^k)^\perp$, a bundle of rank $n-k$ with
$\eps^k\oplus\eta\iso\xi$ by Theorem~\ref{thm:orthogonal-complement}.

\step{3} By Proposition~\ref{prop:stability},
$w_i(\xi)=w_i(\eps^k\oplus\eta)=w_i(\eta)$, and by (A1) applied to $\eta$ this
vanishes for $i>n-k$.
\end{proof}

\begin{remark}
Read contrapositively: if $w_{n-k}(TM)\neq0$ then $M$ admits at most $k$
everywhere independent vector fields. Stiefel--Whitney classes are
\emph{obstructions to finding independent sections}, and this is the historical
origin of the theory.
\end{remark}

\section{The total class and Whitney duality}

Let $H^{\Pi}(B;\Z_2)$ denote the ring of formal series
$a=a_0+a_1+a_2+\cdots$ with $a_i\in H^i(B;\Z_2)$, with the evident sum and
product; the product is commutative and associative. The \emph{total
Stiefel--Whitney class} of an $\R^n$-bundle is
\[
w(\xi)=1+w_1(\xi)+\dots+w_n(\xi)\ \in\ H^{\Pi}(B;\Z_2),
\]
and the Whitney formula (A3) becomes the single identity
\[
w(\xi\oplus\eta)=w(\xi)\,w(\eta).
\]

\begin{proposition}\label{prop:invertible}
The series in $H^\Pi(B;\Z_2)$ with leading term $1$ form an abelian group under
multiplication.
\end{proposition}

\begin{proof}
\begin{strategy}
Solve $w\bar w=1$ degree by degree; each equation determines the next unknown
uniquely.
\end{strategy}

\step{1} Given $w=1+w_1+w_2+\cdots$, we seek $\bar w=1+\bar w_1+\bar w_2+\cdots$
with $w\bar w=1$. Comparing terms of degree $k\ge1$ gives
\[
\bar w_k+\bar w_{k-1}w_1+\dots+\bar w_1 w_{k-1}+w_k=0 .
\]

\step{2} This determines $\bar w_k$ uniquely in terms of $\bar w_1,\dots,
\bar w_{k-1}$ and $w_1,\dots,w_k$, so by induction $\bar w$ exists and is
unique. The first values are
\[
\bar w_1=w_1,
\qquad
\bar w_2=w_1^2+w_2,
\qquad
\bar w_3=w_1^3+w_3 .
\]

\step{3} Associativity and commutativity are inherited from
$H^*(B;\Z_2)$, and the identity element is $1$.
\end{proof}

\begin{lemma}[Whitney duality]\label{lem:whitney-duality}
If $M\subset\R^N$ is a smooth submanifold with normal bundle $\nu$, then
\[
w(\nu)=\bar w(TM),
\qquad\text{that is,}\qquad
w(TM)\,w(\nu)=1 .
\]
\end{lemma}

\begin{proof}
By Example~\ref{ex:normal-bundle}, $TM\oplus\nu\iso T\R^N|_M$, which is trivial.
Proposition~\ref{prop:trivial-vanishes} gives $w(TM\oplus\nu)=1$, and the
Whitney formula turns this into $w(TM)w(\nu)=1$; now apply
Proposition~\ref{prop:invertible}.
\end{proof}

\begin{example}
The normal bundle of $S^n\subset\R^{n+1}$ is trivial, spanned by the outward
unit normal, so $w(\nu)=1$ and therefore $w(TS^n)=1$: all Stiefel--Whitney
classes of spheres vanish. Yet $TS^2$ is non-trivial. The classes are powerful,
but they are not complete invariants.
\end{example}

\section{Computations on projective spaces}

Recall $H^*(\RP^n;\Z_2)=\Z_2[a]/(a^{n+1})$ with $a$ the generator of
$H^1(\RP^n;\Z_2)$.

\begin{example}\label{ex:w-gamma}
$w(\gaman)=1+a$.

Indeed, the inclusion $j\colon\RP^1\hookrightarrow\RP^n$ induced by
$\R^2\subset\R^{n+1}$ is covered by a bundle map
$\gamma^1_1\to\gaman$, so (A2) gives $w_1(\gamma^1_1)=j^*w_1(\gaman)$. By (A4)
the left side is non-zero, so $w_1(\gaman)\neq0$; since
$H^1(\RP^n;\Z_2)=\Z_2\{a\}$, this forces $w_1(\gaman)=a$. The bundle has rank
$1$, so $w_i=0$ for $i\ge2$ by (A1).
\end{example}

\begin{example}\label{ex:w-gamma-perp}
Let $\gamma^\perp$ be the orthogonal complement of $\gaman$ inside the trivial
bundle $\RP^n\times\R^{n+1}$. Then
\[
w(\gamma^\perp)=1+a+a^2+\dots+a^n .
\]
Indeed $\gaman\oplus\gamma^\perp$ is trivial, so
$w(\gamma^\perp)=\overline{(1+a)}$, and one checks directly in
$\Z_2[a]/(a^{n+1})$ that
\[
(1+a)(1+a+\dots+a^n)=1+a^{n+1}=1 .
\]
\end{example}

\begin{lemma}\label{lem:TRPn}
There are isomorphisms
\[
T\RP^n\ \iso\ \Hom(\gaman,\gamma^\perp)
\qquad\text{and}\qquad
T\RP^n\oplus\eps^1\ \iso\
\underbrace{\gaman\oplus\dots\oplus\gaman}_{n+1},
\]
where $\eps^1$ is the trivial line bundle.
\end{lemma}

\begin{proof}
\begin{strategy}
Describe tangent vectors of $\RP^n$ as antipodal pairs of tangent vectors of
$S^n$, and observe that such a pair is exactly a linear map from the line
through $x$ to its orthogonal complement. The second isomorphism then comes from
adding the summand $\Hom(\gaman,\gaman)$, which is trivial, and using
$\gamma^\perp\oplus\gaman=\eps^{n+1}$.
\end{strategy}

\step{1} \emph{Tangent vectors upstairs and downstairs.} Let $q\colon S^n\to
\RP^n$ be the quotient map. Since $q$ is a local diffeomorphism, $Dq$ identifies
\[
T\RP^n=\big\{\,\{(x,v),(-x,-v)\}\ \big|\ x\cdot x=1,\ x\cdot v=0\,\big\},
\]
because $Dq(x,v)=Dq(-x,-v)$ and these are the only identifications.

\step{2} \emph{From pairs to homomorphisms.} Given such a pair, let
$L=\R x\subset\R^{n+1}$ be the corresponding line, that is, the fibre of $\gaman$
over $\{\pm x\}$, and define a linear map
\[
\ell\colon L\longrightarrow L^\perp,
\qquad
\ell(x)=v .
\]
Replacing $(x,v)$ by $(-x,-v)$ gives the same $\ell$, since
$\ell(-x)=-v$. Conversely $\ell$ determines the pair. This assignment is a
bijection
\[
T\RP^n\longrightarrow \Hom(\gaman,\gamma^\perp),
\]
continuous with continuous inverse in local trivializations, and linear on each
fibre; Lemma~\ref{lem:fibrewise-iso} makes it an isomorphism.

\step{3} \emph{The trivial summand.} The line bundle
$\Hom(\gaman,\gaman)$ has the nowhere zero section $b\mapsto\mathrm{id}$, hence
is trivial by Theorem~\ref{thm:trivial-iff-sections}:
$\Hom(\gaman,\gaman)\iso\eps^1$.

\step{4} \emph{Adding up.} Using Step 2, Step 3 and additivity of $\Hom$ in the
second variable,
\[
T\RP^n\oplus\eps^1
\iso\Hom(\gaman,\gamma^\perp)\oplus\Hom(\gaman,\gaman)
\iso\Hom\big(\gaman,\ \gamma^\perp\oplus\gaman\big),
\]
and the right-hand side is $\Hom\big(\gaman,\eps^{n+1}\big)$,
since $\gamma^\perp\oplus\gaman$ is the trivial bundle
$\RP^n\times\R^{n+1}$.

\step{5} \emph{Splitting the trivial bundle.} Finally
\[
\Hom\big(\gaman,\eps^{n+1}\big)
\iso\bigoplus_{i=1}^{n+1}\Hom\big(\gaman,\eps^1\big)
=\bigoplus_{i=1}^{n+1}(\gaman)^*
\iso\bigoplus_{i=1}^{n+1}\gaman,
\]
the last step because a line bundle with a Euclidean metric is isomorphic to its
own dual.
\end{proof}

\begin{theorem}\label{thm:w-RPn}
$w(T\RP^n)=(1+a)^{n+1}$, that is,
$w_i(T\RP^n)=\binom{n+1}{i}a^i\bmod 2$.
\end{theorem}

\begin{proof}
By Proposition~\ref{prop:stability} and Lemma~\ref{lem:TRPn},
\[
w(T\RP^n)=w\big(T\RP^n\oplus\eps^1\big)
=w\big((\gaman)^{\oplus(n+1)}\big)
=w(\gaman)^{n+1}=(1+a)^{n+1},
\]
using the Whitney formula and Example~\ref{ex:w-gamma}.
\end{proof}

The mod-2 Pascal triangle
\[
\begin{array}{ccccccccc}
 & & & &1& & & &\\
 & & &1& &1& & &\\
 & &1& &0& &1& &\\
 &1& &1& &1& &1&\\
1& &0& &0& &0& &1
\end{array}
\]
already gives $w(T\RP^3)=1$ and $w(T\RP^4)=1+a+a^4$.

\begin{corollary}[Stiefel]\label{cor:parallelizable}
If $\RP^n$ is parallelizable then $n+1$ is a power of $2$.
\end{corollary}

\begin{proof}
\begin{strategy}
Parallelizable means $T\RP^n$ trivial, hence $w=1$. Over $\Z_2$ the Frobenius
identity $(1+a)^{2^r}=1+a^{2^r}$ lets us factor $(1+a)^{n+1}$ and exhibit a
surviving term unless $n+1$ is a power of $2$.
\end{strategy}

\step{1} If $T\RP^n$ is trivial then $w(T\RP^n)=1$ by
Proposition~\ref{prop:trivial-vanishes}, so $(1+a)^{n+1}=1$ in
$\Z_2[a]/(a^{n+1})$.

\step{2} Write $n+1=2^r m$ with $m$ odd. Since squaring is additive mod $2$,
\[
(1+a)^{2^r}=1+a^{2^r},
\]
and therefore
\[
(1+a)^{n+1}=\big(1+a^{2^r}\big)^{m}
=1+m\,a^{2^r}+\binom{m}{2}a^{2\cdot 2^r}+\cdots
\]

\step{3} Suppose $m>1$. Then $2^r<n+1$, so $2^r\le n$ and the class $a^{2^r}$ is
non-zero in $\Z_2[a]/(a^{n+1})$. Its coefficient is $m\equiv1\bmod2$, so
$(1+a)^{n+1}\neq1$, contradicting Step 1. Hence $m=1$ and $n+1=2^r$.
\end{proof}

\begin{remark}
The converse is false: the only parallelizable real projective spaces are
$\RP^1,\RP^3,\RP^7$, corresponding to $\C$, the quaternions and the octonions.
This requires deeper tools (Bott--Milnor, Kervaire, Adams) and is not proved
here.
\end{remark}

\begin{example}[Immersions of $\RP^9$]\label{ex:immersion-RP9}
Suppose $M^n$ immerses in $\R^{n+k}$. By Lemma~\ref{lem:normal-immersion} the
normal bundle $\nu_f$ has rank $k$, and $TM\oplus\nu_f$ is trivial, so that
$w_i(\nu_f)=\bar w_i(TM)$. Therefore
\[
\bar w_i(TM)=0\qquad\text{for } i>k .
\]
For $M=\RP^9$: since $10=8+2$,
\[
w(T\RP^9)=(1+a)^{10}=(1+a^{8})(1+a^{2})=1+a^{2}+a^{8}
\]
in $\Z_2[a]/(a^{10})$. Inverting, with $b=a^2$ and $b^5=a^{10}=0$, the equation
$(1+b+b^4)\,\bar w=1$ gives successively the coefficients $1,1,1,1,0$, so
\[
\bar w(T\RP^9)=1+a^{2}+a^{4}+a^{6}.
\]
The top non-zero term sits in degree $6$, so $k\ge6$: the space $\RP^9$ cannot be
immersed in $\R^{9+k}$ for $k<6$.
\end{example}

\begin{example}[Immersions of $\RP^{2^r}$]
Let $n=2^r$. Then
\[
w(T\RP^n)=(1+a)^{2^r+1}=\big(1+a^{2^r}\big)(1+a)=1+a+a^{n}
\]
in $\Z_2[a]/(a^{n+1})$, and one checks
\[
\bar w(T\RP^{n})=1+a+a^2+\dots+a^{n-1},
\]
since $(1+a)(1+a+\dots+a^{n-1})=1+a^n$ telescopes and
$a^{n}(1+a+\dots+a^{n-1})=a^n$ modulo $a^{n+1}$, so the two extra terms cancel.
The top non-zero term is $a^{n-1}$, so any immersion of $\RP^{2^r}$ needs
$k\ge n-1$, that is, an ambient $\R^{2n-1}$. This is exactly the bound of
Whitney's immersion theorem, so projective spaces of $2$-power dimension are the
worst case.
\end{example}

\section{Stiefel--Whitney numbers}

\begin{guiding}
Classes live in the cohomology of $M$, which makes them awkward to compare
between different manifolds. Can we extract from them plain \emph{numbers}, and
what do those numbers see?
\end{guiding}

\begin{definition}
Let $M$ be a closed smooth $n$-manifold with fundamental class
$\mu_M\in H_n(M;\Z_2)$. For every partition
$r_1+2r_2+\dots+nr_n=n$ the associated \emph{Stiefel--Whitney number} is
\[
w_1^{r_1}\cdots w_n^{r_n}[M]
:=\big\langle\, w_1(TM)^{r_1}\smile\cdots\smile w_n(TM)^{r_n},\ \mu_M\,\big\rangle
\ \in\ \Z_2 .
\]
Two $n$-manifolds \emph{have the same Stiefel--Whitney numbers} if these agree
for every such $(r_1,\dots,r_n)$.
\end{definition}

\begin{example}
For $n$ even, $w_n(T\RP^n)=\binom{n+1}{n}a^n=(n+1)a^n\neq0$ and likewise
$w_1(T\RP^n)^n=(n+1)^na^n\neq0$, so $\RP^n$ has non-zero Stiefel--Whitney
numbers. For $n$ odd they all vanish, as we verify at the start of the next
lecture --- and this is no accident: odd projective spaces are boundaries.
\end{example}

In Lecture 3 we prove that Stiefel--Whitney numbers vanish on boundaries, and
state Thom's theorem that they classify closed manifolds up to unoriented
cobordism.

\lecture

\begin{guiding}
Stiefel--Whitney numbers are computable invariants of closed manifolds. Which
equivalence relation do they detect? And is there a universal space from which
\emph{all} characteristic classes originate?
\end{guiding}

\section{Stiefel--Whitney numbers and cobordism}

\begin{example}\label{ex:odd-RP-numbers}
All Stiefel--Whitney numbers of $\RP^{2n-1}$ vanish.

Indeed $w(T\RP^{2n-1})=(1+a)^{2n}=(1+a^2)^n$ contains only terms of even degree,
so $w_i(T\RP^{2n-1})=0$ whenever $i$ is odd. A monomial
$w_1^{r_1}\cdots w_{2n-1}^{r_{2n-1}}$ of total degree $\sum i\,r_i=2n-1$ must
have $i\,r_i$ odd for at least one $i$, hence both $i$ and $r_i$ odd; that
factor is $w_i^{r_i}=0$.
\end{example}

\begin{theorem}\label{thm:boundary-numbers-vanish}
If $M^n$ is the boundary of a compact smooth $(n+1)$-manifold $B$, then all
Stiefel--Whitney numbers of $M$ vanish.
\end{theorem}

\begin{proof}
\begin{strategy}
The tangent bundle of $M$ is the restriction of the tangent bundle of $B$, up to
a trivial summand, so every characteristic class of $M$ in top degree is
restricted from $B$. On the other hand the fundamental class of $M$ is the
boundary of the fundamental class of $B$, and the evaluation pairing turns
$\partial$ into $\delta$. Exactness of the sequence of the pair kills the
result.
\end{strategy}

\step{1} \emph{Classes come from $B$.} Along $M=\partial B$ there is a
well-defined outward pointing direction, giving a nowhere zero section of
$TB|_M$ transverse to $TM$, hence a trivial line sub-bundle $\eps^1$ with
\[
TB|_M\ \iso\ TM\oplus\eps^1 .
\]
By Proposition~\ref{prop:stability} and naturality, for the inclusion
$i\colon M\hookrightarrow B$,
\[
i^*w_j(TB)=w_j\big(TB|_M\big)=w_j(TM).
\]

\step{2} \emph{Fundamental classes.} With $\Z_2$ coefficients,
$H_{n+1}(B,M;\Z_2)$ contains the fundamental class $\mu_B$ of the compact
manifold with boundary, and the connecting homomorphism of the pair sends
\[
\partial\colon H_{n+1}(B,M;\Z_2)\longrightarrow H_n(M;\Z_2),
\qquad
\mu_B\longmapsto\mu_M .
\]

\step{3} \emph{Transferring the evaluation.} Let
$v=w_1(TB)^{r_1}\smile\cdots\smile w_n(TB)^{r_n}\in H^n(B;\Z_2)$ with
$\sum i\,r_i=n$. By Step 1, $i^*v$ is the corresponding monomial in the classes
of $TM$. Using Step 2 and the definition of the connecting maps,
\[
\big\langle i^*v,\ \mu_M\big\rangle
=\big\langle i^*v,\ \partial\mu_B\big\rangle
=\big\langle \delta\,i^*v,\ \mu_B\big\rangle .
\]

\step{4} \emph{Exactness.} The cohomology sequence of the pair $(B,M)$ contains
\[
H^n(B;\Z_2)\ \xrightarrow{\ i^*\ }\ H^n(M;\Z_2)
\ \xrightarrow{\ \delta\ }\ H^{n+1}(B,M;\Z_2),
\]
so $\delta\circ i^*=0$. Hence the expression in Step 3 vanishes, and every
Stiefel--Whitney number of $M$ is zero.
\end{proof}

\begin{center}
\begin{tikzpicture}[scale=0.9]
\draw[thick] (0,0) ellipse (0.45 and 0.16);
\draw[thick] (2,0) ellipse (0.45 and 0.16);
\draw[thick] (1,-2.2) ellipse (0.55 and 0.18);
\draw[thick] (-0.45,0) .. controls (-0.4,-1.2) .. (0.45,-2.2);
\draw[thick] (2.45,0) .. controls (2.4,-1.2) .. (1.55,-2.2);
\draw[thick] (0.45,0) .. controls (0.7,-0.7) and (1.3,-0.7) .. (1.55,0);
\node at (4.1,-1.1) {\small a cobordism $\partial B=M_1\sqcup M_2$};
\end{tikzpicture}
\end{center}

\begin{theorem}[Thom; quoted without proof]
Conversely, if all Stiefel--Whitney numbers of a closed smooth manifold $M$
vanish, then $M$ is a boundary.
\end{theorem}

This converse is genuinely hard; it is proved by the cobordism machinery
sketched in Lecture 12. See Milnor--Stasheff, \S 4 and \S 18, or Thom's original
paper.

\begin{definition}
Two closed smooth $n$-manifolds $M_1,M_2$ are \emph{(unoriented) cobordant} if
there is a compact smooth manifold $B$ with $\partial B=M_1\sqcup M_2$.
\end{definition}

\begin{corollary}
Two closed smooth manifolds are cobordant if and only if their
Stiefel--Whitney numbers are equal. In particular the numbers determine the
unoriented cobordism class.
\end{corollary}

\begin{proof}
If they are cobordant, Theorem~\ref{thm:boundary-numbers-vanish} applied to
$M_1\sqcup M_2$ gives $w_I[M_1]+w_I[M_2]=0$ for every partition $I$, hence
equality since we work mod $2$. Conversely, if the numbers agree then all
numbers of $M_1\sqcup M_2$ vanish and Thom's theorem makes $M_1\sqcup M_2$ a
boundary, that is, a cobordism between $M_1$ and $M_2$.
\end{proof}

\section{Classifying spaces}

\begin{guiding}
Where do characteristic classes come from? The slogan to aim at: bundles over
$B$ are the same thing as homotopy classes of maps from $B$ to one fixed
universal space, so that every cohomology class of that space produces an
invariant of bundles.
\end{guiding}

Recall that a \emph{fibre bundle} $(\pi,E,B,F)$ is a continuous surjection
$\pi\colon E\to B$ with local trivializations $U\times F\iso\pi^{-1}(U)$, and
that a \emph{principal $G$-bundle} is a fibre bundle carrying a $G$-action which
preserves the fibres and is free and transitive on each of them, so that each
fibre is homeomorphic to $G$.

Given a topological group $G$ there is a contractible space $EG$ on which $G$
acts freely; the quotient map
\[
\pi\colon EG\longrightarrow BG=EG/G
\]
is the \emph{universal principal $G$-bundle}, and $BG$ is the \emph{classifying
space}. We quote the following.

\begin{theorem}[quoted]
If $M$ is paracompact and $P\to M$ is a principal $G$-bundle, then there is a
map $f\colon M\to BG$ with $f^*EG\iso P$.
\end{theorem}

For vector bundles the relevant group is $\GL_n(\R)$, or $\Or(n)$ in the
presence of a Euclidean metric: to an $\R^n$-bundle one associates its
\emph{frame bundle}, the principal $\GL_n(\R)$-bundle whose fibre over $b$ is
the set of frames of $F_b(\xi)$; orthonormal frames give a principal
$\Or(n)$-bundle. Conversely a principal bundle $P$ recovers the vector bundle
$P\times_{\GL_n}\R^n$. So the classifying space we want is $B\Or(n)$, and it has
a concrete model: the Grassmannian, with $E\Or(n)$ modelled on a Stiefel
manifold.

\section{Grassmannians and universal bundles}

The motivation is geometric. A curve $C\subset\R^{k+1}$ has a tangent line at
each point, giving a map $C\to\RP^k$; an $n$-manifold $M\subset\R^{n+k}$ has a
tangent $n$-plane at each point, giving a map to the space of all $n$-planes.

\begin{definition}
The \emph{Grassmann manifold} $G_n(\R^{n+k})$ is the set of $n$-dimensional
linear subspaces of $\R^{n+k}$. In particular $G_1(\R^{k+1})=\RP^k$.
\end{definition}

\begin{definition}
An \emph{$n$-frame} in $\R^{n+k}$ is an $n$-tuple of linearly independent
vectors. The \emph{Stiefel manifold} $V_n(\R^{n+k})\subset(\R^{n+k})^n$ is the
set of $n$-frames, an open subset since it is
$\{A\mid\det AA^{\mathsf T}\neq0\}$, the preimage of $\R\smallsetminus\{0\}$
under a continuous map. Let $V_n^0(\R^{n+k})\subset V_n(\R^{n+k})$ be the
subspace of \emph{orthonormal} frames.
\end{definition}

We topologize $G_n(\R^{n+k})$ as the quotient of $V_n(\R^{n+k})$ under
$q=\spann$.

\begin{proposition}
$G_n(\R^{n+k})$ is compact.
\end{proposition}

\begin{proof}
The set $V_n^0(\R^{n+k})$ is closed and bounded in $(\R^{n+k})^n$ --- it is cut
out by the equations $x_i\cdot x_j=\delta_{ij}$ --- hence compact. Gram--Schmidt
is a continuous retraction $V_n\to V_n^0$ commuting with $q$, so
$q|_{V_n^0}\colon V_n^0\to G_n(\R^{n+k})$ is onto. A continuous image of a
compact space is compact.
\end{proof}

\begin{theorem}\label{thm:grassmannian-manifold}
$G_n(\R^{n+k})$ is a topological manifold of dimension $nk$, and
$X\mapsto X^\perp$ is a homeomorphism
$G_n(\R^{n+k})\iso G_k(\R^{n+k})$.
\end{theorem}

\begin{proof}
\begin{strategy}
Hausdorffness is proved by separating two planes with an explicit continuous
function, namely the distance to a vector lying in one and not the other.
Charts come from the observation that the planes close to a fixed $X_0$ are
exactly the graphs of linear maps $X_0\to X_0^\perp$; the identification with
$\Hom(X_0,X_0^\perp)$ is the chart.
\end{strategy}

\step{1} \emph{A supply of continuous functions.} For $w\in\R^{n+k}$ let
$\rho_w(X)$ be the squared distance from $w$ to the plane $X$. If
$(x_1,\dots,x_n)$ is an orthonormal frame of $X$ then
\[
\rho_w(X)=w\cdot w-(w\cdot x_1)^2-\dots-(w\cdot x_n)^2,
\]
which is continuous in the frame. Hence $\rho_w\circ q$ is continuous on
$V^0_n$, and since $G_n$ carries the quotient topology, $\rho_w$ is continuous
on $G_n(\R^{n+k})$.

\step{2} \emph{Hausdorff.} Let $X_1\neq X_2$ and choose
$w\in X_1\smallsetminus X_2$. Then $\rho_w(X_1)=0$ while $\rho_w(X_2)>0$, so the
preimages under $\rho_w$ of two disjoint intervals separate $X_1$ from $X_2$.

\step{3} \emph{The chart domain.} Fix $X_0$ and split
$\R^{n+k}=X_0\oplus X_0^{\perp}$. Let
\[
U=\big\{Y\in G_n(\R^{n+k})\ \big|\
\proj_{X_0}\!\big|_Y\colon Y\to X_0 \text{ is an isomorphism}\big\}
\]
that is, $U=\{Y\mid Y\cap X_0^{\perp}=0\}$. This is
an open neighbourhood of $X_0$: writing $Y=q(y_1,\dots,y_n)$, the condition is
that a certain determinant in the frame be non-zero.

\step{4} \emph{The chart.} For $Y\in U$ every element of $Y$ is determined by
its $X_0$-component, so $Y$ is the graph
\[
Y=\big\{\,x+T(Y)x\ \big|\ x\in X_0\,\big\}
\]
of a unique linear map $T(Y)\colon X_0\to X_0^\perp$. This gives a bijection
\[
T\colon U\longrightarrow\Hom(X_0,X_0^\perp)\iso\R^{nk}.
\]

\step{5} \emph{Continuity both ways.} Choose an orthonormal basis
$x_1,\dots,x_n$ of $X_0$. For $Y\in U$ the vectors $y_i=x_i+T(Y)x_i$ form a
frame of $Y$, and conversely any frame of $Y$ close to that of $X_0$ can be
normalized to this form by an invertible linear change, depending continuously
on the frame. Hence $y_i$ depends continuously on $Y$ and vice versa, and
$T(Y)$ is determined by the $y_i$ through $T(Y)x_i=y_i-x_i$. So both $T$ and
$T^{-1}$ are continuous, and $U\iso\R^{nk}$.

\step{6} \emph{The duality homeomorphism.} The assignment $Y\mapsto Y^{\perp}$
is a bijection $G_n(\R^{n+k})\to G_k(\R^{n+k})$, and it is continuous: an
orthonormal frame of $Y$ can be completed continuously, locally in $Y$, to an
orthonormal frame of $\R^{n+k}$ by the Gram--Schmidt argument of
Theorem~\ref{thm:orthogonal-complement}, and the last $k$ vectors give a frame
of $Y^\perp$. Applying the same to $Y^\perp$ gives continuity of the inverse.
\end{proof}

\begin{definition}
The \emph{tautological bundle} $\gamma^n(\R^{n+k})$ over $G_n(\R^{n+k})$ has
total space
\[
E=\big\{(X,\vec v)\ \big|\ X\in G_n(\R^{n+k}),\ \vec v\in X\big\}
\]
and projection $(X,\vec v)\mapsto X$. For $n=1$ this is
$\gamma^1(\R^{k+1})=\gamma^1_k\to\RP^k$.
\end{definition}

\begin{lemma}\label{lem:tautological-locally-trivial}
$\gamma^n(\R^{n+k})$ is locally trivial, hence a vector bundle of rank $n$.
\end{lemma}

\begin{proof}
Work over the chart domain $U$ of $X_0$ from Step 3 above. For $Y\in U$ the
orthogonal projection $P_Y\colon Y\to X_0$ is an isomorphism, so we may define
\[
h\colon U\times X_0\longrightarrow\pi^{-1}(U),
\qquad
h(Y,x)=\big(Y,\ P_Y^{-1}(x)\big)=\big(Y,\ x+T(Y)x\big).
\]
By Step 5 of Theorem~\ref{thm:grassmannian-manifold} the map $T$ is continuous,
so $h$ is continuous; its inverse is $(Y,v)\mapsto(Y,\proj_{X_0}v)$, also
continuous; and for fixed $Y$ it is a linear isomorphism $X_0\to Y$. Identifying
$X_0\iso\R^n$ gives a local trivialization.
\end{proof}

\begin{definition}
For a smooth $M^n\subset\R^{n+k}$, the \emph{generalized Gauss map} is
\[
\bar g\colon M\longrightarrow G_n(\R^{n+k}),
\qquad
x\longmapsto T_xM,
\]
the base map of the bundle map
$g\colon TM\to\gamma^n(\R^{n+k})$, $(x,v)\mapsto(T_xM,v)$.
\end{definition}

So tangent bundles of submanifolds of Euclidean space are pulled back from
tautological bundles. The following says this is neither an accident nor a
restriction.

\begin{lemma}\label{lem:compact-classify}
Let $\xi$ be an $\R^n$-bundle over a \emph{compact} base $B$. Then for $k$ large
enough there is a bundle map $\xi\to\gamma^n(\R^{n+k})$.
\end{lemma}

\begin{proof}
\begin{strategy}
A bundle map into a tautological bundle is the same thing as a map from the
total space to a Euclidean space, linear and injective on each fibre; the plane
component of the bundle map is then forced. Such a map is assembled from
finitely many local trivializations, damped by bump functions so that they
extend by zero, and stacked into a single large Euclidean space.
\end{strategy}

\step{1} \emph{Reduction.} Suppose we have a continuous
$\hat f\colon E(\xi)\to\R^m$ which is linear and injective on every fibre. Then
\[
f(e)=\Big(\ \hat f\big(F_{\pi(e)}(\xi)\big),\ \hat f(e)\ \Big)
\]
defines a map $E(\xi)\to E\big(\gamma^n(\R^m)\big)$: the first entry is an
$n$-plane in $\R^m$ because $\hat f$ is injective and linear on the fibre, and
the second lies in it. It is fibrewise a linear isomorphism onto that plane. Its
continuity follows from the continuity of $\hat f$ together with the local
description of the chart in Lemma~\ref{lem:tautological-locally-trivial}.

\step{2} \emph{Bump functions.} By compactness choose a finite trivializing
cover $U_1,\dots,U_r$ of $B$, sets $W_i\subset\overline{W_i}\subset V_i\subset
\overline{V_i}\subset U_i$ with $\{W_i\}$ still covering $B$, and continuous
$\lambda_i\colon B\to[0,1]$ with $\lambda_i=1$ on $W_i$ and $\lambda_i=0$
outside $V_i$.

\step{3} \emph{Damped local coordinates.} Each trivialization
$h_i\colon U_i\times\R^n\to\pi^{-1}(U_i)$ gives
$\tilde h_i=\proj_{\R^n}\circ\,h_i^{-1}\colon\pi^{-1}(U_i)\to\R^n$, linear and
injective on each fibre over $U_i$. Define
\[
h'_i\colon E(\xi)\longrightarrow\R^n,
\qquad
h'_i(e)=\lambda_i\big(\pi(e)\big)\cdot\tilde h_i(e)
\]
for $\pi(e)\in U_i$, and $h'_i(e)=0$ otherwise. This is continuous: the two
definitions agree on the overlap $U_i\smallsetminus \overline{V_i}$, where
$\lambda_i=0$.

\step{4} \emph{Stacking.} Set $m=rn$ and
\[
\hat f=(h'_1,\dots,h'_r)\colon E(\xi)\longrightarrow
\underbrace{\R^n\oplus\dots\oplus\R^n}_{r}=\R^{m}.
\]
It is continuous and linear on fibres. It is injective on the fibre over $b$:
some $W_i$ contains $b$, so $\lambda_i(b)=1$ and the $i$-th component is
$\tilde h_i$, already injective there.

\step{5} By Step 1 this $\hat f$ produces a bundle map
$\xi\to\gamma^n(\R^{m})$, and $m=rn$ is of the form $n+k$ with $k=(r-1)n$.
\end{proof}

To remove the dependence on $k$, let $\R^\infty=\bigcup_m\R^m$ carry the direct
limit topology and define the \emph{infinite Grassmannian}
\[
G_n=G_n(\R^\infty)=\varinjlim_k G_n(\R^{n+k}),
\]
with tautological bundle $\gamma^n$ defined by the same formula. This is the
\emph{universal $n$-plane bundle}. In the next lecture we prove that every
$\R^n$-bundle over a paracompact base admits a bundle map to $\gamma^n$, unique
up to homotopy, so that
\[
\big\{\R^n\text{-bundles over }B\big\}\big/\!\iso
\ \ \longleftrightarrow\ \
\big[\,B,\ G_n\,\big],
\]
and characteristic classes are then simply elements of $H^*(G_n)$ pulled back
along classifying maps.

\lecture

\begin{guiding}
For which base spaces does the slogan ``bundles are the same as maps to the
Grassmannian'' hold, and what does the Grassmannian look like as a space? The
technical hypothesis is paracompactness; the geometric answer is that the
Grassmannian is a CW complex with exactly one cell for each Schubert symbol.
\end{guiding}

\section{Paracompactness}

\begin{definition}
A space is \emph{paracompact} if every open cover admits a locally finite
refinement, that is, a refinement such that every point has a neighbourhood
meeting only finitely many of its members.
\end{definition}

\begin{theorem}[quoted]
Every metric space is paracompact and Hausdorff. A regular space which is a
countable union of compact subsets is paracompact and Hausdorff. Consequently
the direct limit of a sequence of compact Hausdorff spaces is paracompact.
\end{theorem}

\begin{remark}
The infinite Grassmannian $G_n=\varinjlim_k G_n(\R^{n+k})$ is therefore
paracompact, being a direct limit of compact spaces. All bases we shall meet ---
manifolds, CW complexes, Grassmannians --- are paracompact.
\end{remark}

The following technical lemma replaces an arbitrary trivializing cover by a
countable one. It is the device that makes the classification theorem work for
non-compact bases.

\begin{lemma}\label{lem:countable-cover}
Let $\xi$ be a vector bundle over a paracompact Hausdorff base $B$. Then there
is a countable open cover $U_1,U_2,\dots$ of $B$ such that $\xi|_{U_k}$ is
trivial for every $k$.
\end{lemma}

\begin{proof}
\begin{strategy}
Start from a locally finite trivializing cover $\{V_\alpha\}$ with subordinate
bump functions $\lambda_\alpha$. For each \emph{finite} index set $S$ collect the
points where the functions indexed by $S$ dominate all the others; this produces
an open set contained in $\bigcap_{\alpha\in S}V_\alpha$, hence trivializing.
Two different sets of the same cardinality give disjoint pieces, so we may take
their union and index the resulting cover by the cardinality $k=|S|$ alone.
\end{strategy}

\step{1} \emph{Set-up.} Choose a locally finite open cover $\{V_\alpha\}$ of $B$
such that $\xi|_{V_\alpha}$ is trivial, and continuous functions
$\lambda_\alpha\colon B\to[0,1]$ with $\lambda_\alpha=0$ outside $V_\alpha$ and
with the open sets $W_\alpha=\lambda_\alpha^{-1}\big((0,1]\big)$ still covering
$B$. Such data exists by paracompactness.

\step{2} \emph{The sets $U(S)$.} For a finite non-empty set $S$ of indices put
\[
U(S)=\Big\{b\in B\ \Big|\ \min_{\alpha\in S}\lambda_\alpha(b)
\ >\ \max_{\beta\notin S}\lambda_\beta(b)\Big\}.
\]
Near any point only finitely many $\lambda$'s are non-zero, so both the minimum
and the maximum are continuous there, and $U(S)$ is open.

\step{3} \emph{$U(S)$ is trivializing.} If $b\in U(S)$ and $\alpha\in S$ then
$\lambda_\alpha(b)>0$, so $b\in V_\alpha$. Hence
\[
U(S)\ \subset\ \bigcap_{\alpha\in S}V_\alpha\ \subset\ V_\alpha ,
\]
and $\xi$ restricted to $U(S)$ is trivial.

\step{4} \emph{Sets of equal size are disjoint.} Let $S\neq S'$ with
$|S|=|S'|=k$. Choose $\alpha\in S\smallsetminus S'$ and
$\alpha'\in S'\smallsetminus S$. If $b$ were in $U(S)\cap U(S')$ we would get
$\lambda_\alpha(b)>\lambda_{\alpha'}(b)$ from $b\in U(S)$, since
$\alpha\in S$ and $\alpha'\notin S$, and symmetrically
$\lambda_{\alpha'}(b)>\lambda_\alpha(b)$ from $b\in U(S')$: a contradiction.

\step{5} \emph{The cover.} Put
\[
U_k=\bigsqcup_{|S|=k}U(S),
\]
a disjoint union of open trivializing sets by Steps 3 and 4, hence itself open
and trivializing, since a bundle trivial on each member of a disjoint family is
trivial on their union. Finally the $U_k$ cover $B$: given $b$, let
$S=\{\alpha\mid\lambda_\alpha(b)\ \text{is maximal}\}$, which is finite and
non-empty by local finiteness and Step 1; then $b\in U(S)\subset U_{|S|}$.
\end{proof}

\section{The universal bundle theorem}

\begin{theorem}\label{thm:universal}
Let $\xi$ be an $\R^n$-bundle over a paracompact base $B$. Then

\begin{enumerate}
\item there exists a bundle map $f\colon\xi\to\gamma^n$;
\item any two bundle maps $\xi\to\gamma^n$ are homotopic through bundle maps.
\end{enumerate}

Consequently the homotopy class of the base map
$\bar f\in\Maps(B,G_n)$ depends only on $\xi$, and $\xi\iso\bar f^{\,*}\gamma^n$.
\end{theorem}

\begin{proof}
\begin{strategy}
Existence repeats the construction of Lemma~\ref{lem:compact-classify}, with the
countable cover of Lemma~\ref{lem:countable-cover} replacing the finite one; the
target is $\R^\infty$ instead of a fixed $\R^m$. For uniqueness, the naive
straight-line homotopy between two fibrewise injective maps
$\hat f,\hat g\colon E(\xi)\to\R^\infty$ works whenever their values are never
negative multiples of each other. In general one first pushes $\hat f$ into the
odd coordinates and $\hat g$ into the even ones, where the obstruction cannot
occur, and chains three such homotopies.
\end{strategy}

\step{1} \emph{Existence.} Let $U_1,U_2,\dots$ be as in
Lemma~\ref{lem:countable-cover}, with trivializations
$h_i\colon U_i\times\R^n\to\pi^{-1}(U_i)$ and bump functions $\lambda_i$ whose
positivity sets cover $B$. Setting
$\tilde h_i=\proj_{\R^n}\circ h_i^{-1}$ and
$h'_i(e)=\lambda_i(\pi(e))\,\tilde h_i(e)$, extended by zero, define
\[
\hat f\colon E(\xi)\longrightarrow \R^n\oplus\R^n\oplus\cdots=\R^\infty,
\qquad
\hat f(e)=\big(h'_1(e),\,h'_2(e),\,\dots\big).
\]
Locally only finitely many terms are non-zero, so $\hat f$ is continuous; it is
linear on fibres, and injective on the fibre over $b$ because
$\lambda_i(b)>0$ for some $i$. As in Step 1 of
Lemma~\ref{lem:compact-classify}, $\hat f$ determines a bundle map
$f\colon\xi\to\gamma^n$ by
$f(e)=\big(\hat f(F_{\pi(e)}(\xi)),\ \hat f(e)\big)$.

\step{2} \emph{The easy homotopy.} Suppose $\hat f,\hat g\colon E(\xi)\to
\R^\infty$ are continuous, linear and injective on fibres, and satisfy
\[
\hat f(e)\ \neq\ -\lambda^2\,\hat g(e)
\qquad\text{for all }e\neq0\text{ and all }\lambda\in\R .
\tag{$\ast$}
\]
Then $\hat h_t=(1-t)\hat f+t\hat g$ is fibrewise linear, and injective on
fibres: if $\hat h_t(v)=0$ with $v\neq0$ and $0<t<1$ then
$\hat f(v)=-\tfrac{t}{1-t}\hat g(v)$, contradicting $(\ast)$; for $t=0,1$ it is
$\hat f$ or $\hat g$. So $\hat h_t$ gives a homotopy of bundle maps from $f$ to
$g$.

\step{3} \emph{Two shift operators.} Let $d_1,d_2\colon\R^\infty\to\R^\infty$ be
the linear injections
\[
d_1(e_i)=e_{2i-1},
\qquad
d_2(e_i)=e_{2i}.
\]
Being injective and linear, each induces a bundle map $\gamma^n\to\gamma^n$
sending $(X,v)\mapsto(d_jX,d_jv)$; and $d_j\circ\hat f$ is again continuous,
linear and fibrewise injective.

\step{4} \emph{$\hat f$ and $d_1\hat f$ satisfy $(\ast)$.} Let $v=\hat f(e)\neq0$
and let $k$ be the index of its first non-zero coordinate. Suppose
$d_1v=-\lambda^2 v$ with $\lambda\neq0$. If $k=1$ then $(d_1v)_1=v_1$, so
$v_1=-\lambda^2v_1$ forces $v_1=0$, a contradiction. If $k>1$ then the first
non-zero coordinate of $d_1v$ sits in position $2k-1>k$, so $(d_1v)_k=0$, giving
$-\lambda^2v_k=0$ and again $v_k=0$. The same argument applies to $\hat g$ and
$d_2\hat g$, using $2k>k$.

\step{5} \emph{$d_1\hat f$ and $d_2\hat g$ satisfy $(\ast)$.} The vector
$d_1\hat f(e)$ is supported in odd coordinates and $d_2\hat g(e)$ in even ones.
If $d_1\hat f(e)=-\lambda^2d_2\hat g(e)$ then both sides vanish, so
$\hat f(e)=0$ and hence $e=0$.

\step{6} \emph{Chaining.} By Steps 2, 4 and 5,
\[
f\ \simeq\ d_1\circ f\ \simeq\ d_2\circ g\ \simeq\ g
\]
through bundle maps, which proves (2).

\step{7} \emph{Conclusion.} A bundle map $f\colon\xi\to\gamma^n$ gives
$\xi\iso\bar f^{\,*}\gamma^n$ by Lemma~\ref{lem:bundle-map-pullback}, and by (2)
the class $[\bar f]$ is well defined.
\end{proof}

\begin{remark}
The converse also holds: homotopic maps $B\to G_n$ pull $\gamma^n$ back to
isomorphic bundles, so
\[
\big\{\R^n\text{-bundles over }B\big\}\big/\!\iso
\ \longleftrightarrow\
[\,B,\,G_n\,].
\]
This gives a clean definition of characteristic classes. For
$c\in H^i(G_n)$ and a bundle $\xi$ with classifying map $f_\xi$, put
\[
c(\xi):=f_\xi^*\,c\ \in\ H^i(B).
\]
Naturality is automatic: if $\eta=g^*\xi$ then $f_\eta\simeq f_\xi\circ g$, so
$c(\eta)=g^*c(\xi)$.
\[
\begin{tikzcd}[row sep=large, column sep=large]
\eta \arrow[d] \arrow[r] & \xi \arrow[d] \arrow[r] & \gamma^n\arrow[d]\\
B_1 \arrow[r,"g"] \arrow[rr, bend right=20, "f_\eta"'] & B_2 \arrow[r,"f_\xi"] & G_n
\end{tikzcd}
\]
The whole theory therefore reduces to one computation: what is $H^*(G_n)$?
\end{remark}

\section{A cell structure on the Grassmannian}

\begin{definition}
A \emph{CW complex} is a Hausdorff space $K$ partitioned into disjoint open
cells $\{e_\alpha\}$ such that

\begin{enumerate}
\item each $e_\alpha$ carries a characteristic map
$f\colon D^{n(\alpha)}\to K$ taking the open disc homeomorphically onto
$e_\alpha$;
\item $\bar e_\alpha\smallsetminus e_\alpha$ is contained in a finite union of
cells of strictly lower dimension;
\item every point lies in a finite subcomplex;
\item $K$ carries the direct limit topology of its finite subcomplexes.
\end{enumerate}
\end{definition}

Fix the standard flag $\R^0\subset\R^1\subset\dots\subset\R^m$, where $\R^k$ is
spanned by the first $k$ coordinates, and let $X\subset\R^m$ be an $n$-plane.
Consider the sequence of integers
\[
0=\dim(X\cap\R^0)\ \le\ \dim(X\cap\R^1)\ \le\ \dots\ \le\ \dim(X\cap\R^m)=n .
\]

\begin{claim}\label{claim:jumps}
Consecutive terms differ by at most $1$.
\end{claim}

\begin{proof}
The sequence
\[
0\longrightarrow X\cap\R^{k-1}\longrightarrow X\cap\R^{k}
\ \xrightarrow{\ x\,\mapsto\, x_k\ }\ \R
\]
is exact: the kernel of the $k$-th coordinate function on $X\cap\R^k$ is exactly
$X\cap\R^{k-1}$. Hence
\[
\dim(X\cap\R^k)=\dim(X\cap\R^{k-1})+\rk(x\mapsto x_k)
\le\dim(X\cap\R^{k-1})+1 .
\]
\end{proof}

So the sequence increases by exactly $1$ at $n$ places, the \emph{jumps}
\[
1\le\sigma_1<\sigma_2<\dots<\sigma_n\le m,
\]
called the \emph{Schubert symbol} of $X$. Define
\[
e(\sigma)=\big\{X\in G_n(\R^m)\ \big|\
\dim(X\cap\R^{\sigma_i})=i \ \text{ and }\
\dim(X\cap\R^{\sigma_i-1})=i-1,
\]
\[
\text{for } i=1,\dots,n\big\}.
\]
Every $X$ lies in exactly one $e(\sigma)$, and there are $\binom{m}{n}$ symbols.

\begin{claim}\label{claim:half-space}
Let $H^{k}\subset\R^k\subset\R^m$ be the half-space
$\{x\in\R^k\mid x_k>0\}$. Then $X\in e(\sigma)$ if and only if $X$ has a basis
$x_1,\dots,x_n$ with $x_i\in H^{\sigma_i}$; moreover such a basis can be chosen
\emph{orthonormal}, and then it is unique.
\end{claim}

\begin{proof}
\begin{strategy}
One direction is a rank computation: a vector whose last non-zero coordinate is
in position $\sigma_i$ forces the dimension to jump there. The other direction
constructs the basis one vector at a time, taking at each jump the unique unit
vector orthogonal to the previous ones and normalized to have positive
$\sigma_i$-th coordinate.
\end{strategy}

\step{1} \emph{Sufficiency.} Suppose $X$ has a basis with $x_i\in H^{\sigma_i}$.
Then $x_1,\dots,x_i\in\R^{\sigma_i}$, so $\dim(X\cap\R^{\sigma_i})\ge i$; and
the $\sigma_i$-th coordinate of $x_i$ is non-zero while
$x_1,\dots,x_{i-1}\in\R^{\sigma_{i-1}}\subset\R^{\sigma_i-1}$, so the map
$X\cap\R^{\sigma_i}\to\R$, $x\mapsto x_{\sigma_i}$, has rank $1$ and
$\dim(X\cap\R^{\sigma_i})>\dim(X\cap\R^{\sigma_i-1})$. Counting the $n$ jumps
this forces $\dim(X\cap\R^{\sigma_i})=i$ exactly, so $X\in e(\sigma)$.

\step{2} \emph{Necessity, first vector.} Let $X\in e(\sigma)$. Then
$X\cap\R^{\sigma_1}$ is a line and $X\cap\R^{\sigma_1-1}=0$, so every non-zero
vector of that line has non-zero $\sigma_1$-th coordinate. There is exactly one
unit vector $x_1$ on it with positive $\sigma_1$-th coordinate, and
$x_1\in H^{\sigma_1}$.

\step{3} \emph{Necessity, induction.} Suppose orthonormal
$x_1,\dots,x_{i-1}$ have been found with $x_j\in H^{\sigma_j}$, spanning
$X\cap\R^{\sigma_{i-1}}$. The space $X\cap\R^{\sigma_i}$ has dimension $i$ and
contains $X\cap\R^{\sigma_{i-1}}$ as a hyperplane, so its orthogonal complement
inside $X\cap\R^{\sigma_i}$ is a line. Since
$\dim(X\cap\R^{\sigma_i-1})=i-1=\dim(X\cap\R^{\sigma_{i-1}})$, no non-zero
vector of that line lies in $\R^{\sigma_i-1}$; hence its two unit vectors have
non-zero $\sigma_i$-th coordinate, and exactly one of them, call it $x_i$, has
it positive. Then $x_i\in H^{\sigma_i}$, and $x_1,\dots,x_i$ is orthonormal.

\step{4} \emph{Uniqueness.} At each stage the choice was forced: the line was
determined by $X$ and the previous vectors, and the sign by the positivity
condition.
\end{proof}

In matrix terms: writing a basis of $X$ as the rows of an $n\times m$ matrix and
row reducing, $X\in e(\sigma)$ if and only if $X$ is the row space of a unique
matrix in reduced echelon form whose pivots occupy the columns
$\sigma_1,\dots,\sigma_n$,
\[
\begin{pmatrix}
*&\cdots&*&1&0&\cdots& & &0\\
*&\cdots&*&*&*&1&0&\cdots&0\\
\vdots& & & & & & &\ddots&\\
*&\cdots&*&*&*&*&\cdots&1&0\cdots 0
\end{pmatrix}
\]
The number of free entries is
\[
d(\sigma)=(\sigma_1-1)+(\sigma_2-2)+\dots+(\sigma_n-n),
\]
which will be the dimension of the cell. Following
Claim~\ref{claim:half-space}, set
\[
e'(\sigma)=V^0_n(\R^m)\cap\big(H^{\sigma_1}\times\dots\times H^{\sigma_n}\big),
\qquad
\bar e'(\sigma)=V^0_n(\R^m)\cap\big(\bar H^{\sigma_1}\times\dots\times
\bar H^{\sigma_n}\big),
\]
so that $q=\spann$ maps $e'(\sigma)$ bijectively onto $e(\sigma)$. In the next
lecture we prove that $\bar e'(\sigma)$ is a closed cell of dimension
$d(\sigma)$ and that these cells assemble into a CW structure.

\lecture

\begin{guiding}
We want $H^*(G_n;\Z_2)$, together with generators we can compute with. The plan
has four movements: finish the cell structure and count cells; prove the
Leray--Hirsch theorem; use projective bundles to obtain a splitting principle;
and conclude that $H^*(G_n;\Z_2)$ is a polynomial ring, which finally
\emph{constructs} the Stiefel--Whitney classes and proves they are unique.
\end{guiding}

\section{Schubert cells, completed}

Recall the notation of Lecture 4: for a symbol
$\sigma=(\sigma_1<\dots<\sigma_n)$ we have $e(\sigma)\subset G_n(\R^m)$, the set
$e'(\sigma)$ of orthonormal frames with $x_i\in H^{\sigma_i}$, its closure
$\bar e'(\sigma)$ where $x_i\in\bar H^{\sigma_i}$, and
$d(\sigma)=\sum_i(\sigma_i-i)$.

We first record the tool that drives the induction.

\begin{lemma}\label{lem:rotation}
For unit vectors $u,v\in\R^m$ let $T(u,v)\in\Or(m)$ be the rotation carrying $u$
to $v$ and fixing $(\R u+\R v)^\perp$ pointwise. Then $T(u,v)(x)$ depends
continuously on $(u,v,x)$, and if $u,v\in\R^k$ then
$T(u,v)x\equiv x \bmod \R^k$ for every $x$.
\end{lemma}

\begin{proof}
On the plane $\R u+\R v$ the rotation is given by an explicit formula in $u,v$
(Gram--Schmidt in that plane followed by the planar rotation through the angle
between $u$ and $v$), and it is the identity on the orthogonal complement; both
pieces depend continuously on $(u,v)$. If $u,v\in\R^k$ then $T(u,v)$ is the
identity on $(\R^k)^\perp$ and preserves $\R^k$, so $T(u,v)x-x\in\R^k$.
\end{proof}

\begin{lemma}\label{lem:closed-cell}
$\bar e'(\sigma)$ is a closed cell of dimension $d(\sigma)$ with interior
$e'(\sigma)$.
\end{lemma}

\begin{proof}
\begin{strategy}
Induct on $n$. For $n=1$ the set is a closed hemisphere. For the inductive step,
observe that a frame in $\bar e'(\sigma_1,\dots,\sigma_{n+1})$ is a frame in
$\bar e'(\sigma_1,\dots,\sigma_n)$ together with one extra unit vector,
orthogonal to the previous ones and lying in $\bar H^{\sigma_{n+1}}$. Rotations
identify the space of possible extra vectors, which varies with the frame, with
one fixed closed cell; this exhibits the whole space as a product.
\end{strategy}

\step{1} \emph{Base case $n=1$.} Here
$\bar e'(\sigma_1)=\{x\in\R^{\sigma_1}\mid |x|=1,\ x_{\sigma_1}\ge0\}$, a closed
hemisphere of $S^{\sigma_1-1}$, hence a closed cell of dimension
$\sigma_1-1=d(\sigma_1)$, with interior the open hemisphere $e'(\sigma_1)$.

\step{2} \emph{The model for the new vector.} Fix
$\sigma_1<\dots<\sigma_{n+1}$ and let $b_i\in\R^m$ be the standard basis vector
with a $1$ in position $\sigma_i$, so that
$(b_1,\dots,b_n)\in e'(\sigma_1,\dots,\sigma_n)$. Put
\[
D=\big\{\,u\in\bar H^{\sigma_{n+1}}\ \big|\ |u|=1,\
b_1\cdot u=\dots=b_n\cdot u=0\,\big\}.
\]
The conditions say that $u$ lies in $\R^{\sigma_{n+1}}$, has last coordinate
$\ge0$, and has vanishing coordinates in the $n$ positions
$\sigma_1,\dots,\sigma_n$. So $D$ is a closed hemisphere of a sphere of
dimension $\sigma_{n+1}-n-1$, hence a closed cell of dimension
$\sigma_{n+1}-n-1$.

\step{3} \emph{Moving a frame to the standard one.} Given
$(x_1,\dots,x_n)\in\bar e'(\sigma_1,\dots,\sigma_n)$, set
\[
T_{(x)}=T(b_n,x_n)\circ\dots\circ T(b_1,x_1)\in\Or(m).
\]
Then $T_{(x)}b_i=x_i$ for every $i$. Indeed $T(b_j,x_j)$ with $j<i$ fixes
$b_i$, because $b_i$ is orthogonal to both $b_j$ and $x_j$: orthogonality to
$x_j$ holds since $x_j\in\R^{\sigma_j}$ has vanishing coordinate in position
$\sigma_i>\sigma_j$. Then $T(b_i,x_i)$ carries $b_i$ to $x_i$, and the later
rotations $T(b_l,x_l)$ with $l>i$ fix $x_i$ for the same reason. The map
$(x)\mapsto T_{(x)}$ is continuous by Lemma~\ref{lem:rotation}.

\step{4} \emph{The product decomposition.} Define
\[
\Psi\colon\bar e'(\sigma_1,\dots,\sigma_n)\times D
\longrightarrow \bar e'(\sigma_1,\dots,\sigma_{n+1})
\]
by
\[
\Psi\big((x_1,\dots,x_n),\,u\big)=\big(x_1,\dots,x_n,\ T_{(x)}u\big).
\]
This lands where claimed. The vector $T_{(x)}u$ is a unit vector; it is
orthogonal to each $x_i$, since
$x_i\cdot T_{(x)}u=T_{(x)}b_i\cdot T_{(x)}u=b_i\cdot u=0$ as $T_{(x)}$ is
orthogonal; and it lies in $\bar H^{\sigma_{n+1}}$, because by
Lemma~\ref{lem:rotation} each rotation used moves a vector only within
$\R^{\sigma_i}\subset\R^{\sigma_{n+1}}$, so
$T_{(x)}u\equiv u\bmod\R^{\sigma_{n+1}-1}$ and the $\sigma_{n+1}$-th coordinate
of $T_{(x)}u$ equals that of $u$, which is $\ge0$.

\step{5} \emph{$\Psi$ is a homeomorphism.} It is continuous by Step 3. It is
injective, since the first $n$ entries determine $T_{(x)}$ and hence
$u=T_{(x)}^{-1}(T_{(x)}u)$. It is surjective: given
$(x_1,\dots,x_{n+1})$ in the target, put $u=T_{(x)}^{-1}x_{n+1}$ and reverse the
computations of Step 4. Source and target are compact Hausdorff, so $\Psi$ is a
homeomorphism.

\step{6} \emph{Conclusion.} By induction $\bar e'(\sigma_1,\dots,\sigma_n)$ is a
closed cell of dimension $\sum_{i\le n}(\sigma_i-i)$; by Step 2, $D$ is a closed
cell of dimension $\sigma_{n+1}-(n+1)$. A product of closed cells is a closed
cell, and dimensions add, giving $d(\sigma_1,\dots,\sigma_{n+1})$. Under $\Psi$
the interiors correspond, since $T_{(x)}u$ has positive $\sigma_{n+1}$-th
coordinate exactly when $u$ does; hence the interior is
$e'(\sigma_1,\dots,\sigma_{n+1})$.
\end{proof}

\begin{lemma}\label{lem:q-characteristic}
The map $q|_{\bar e'(\sigma)}\colon\bar e'(\sigma)\to G_n(\R^m)$ has image
$\overline{e(\sigma)}$ and restricts to a homeomorphism
$e'(\sigma)\to e(\sigma)$. In particular it is a characteristic map for the cell
$e(\sigma)$.
\end{lemma}

\begin{proof}
By Claim~\ref{claim:half-space} the restriction $q\colon e'(\sigma)\to e(\sigma)$
is a bijection; it is continuous, and its inverse assigns to $X$ its unique
positive orthonormal basis, which is continuous by the explicit inductive
construction in the proof of that claim. Since $\bar e'(\sigma)$ is compact and
$G_n(\R^m)$ Hausdorff, $q(\bar e'(\sigma))$ is compact, hence closed, and it
contains $e(\sigma)$, so it contains $\overline{e(\sigma)}$; conversely
$\bar e'(\sigma)$ is the closure of $e'(\sigma)$, so its image lies in the
closure of $e(\sigma)$.
\end{proof}

\begin{theorem}\label{thm:cw-grassmannian}
The $\binom{m}{n}$ sets $e(\sigma)$ are the cells of a CW structure on
$G_n(\R^m)$, with $\dim e(\sigma)=d(\sigma)$. Passing to the limit $m\to\infty$
gives a CW structure on $G_n$ with one cell for each partition into at most $n$
parts.
\end{theorem}

\begin{proof}
\begin{strategy}
Everything except the boundary condition has been proved. For that condition one
shows that a plane in the closure of $e(\sigma)$ has a symbol $\tau$ which is
termwise smaller, hence a strictly smaller cell dimension unless $\tau=\sigma$.
\end{strategy}

\step{1} By Lemma~\ref{lem:closed-cell} and Lemma~\ref{lem:q-characteristic},
each $e(\sigma)$ is an open cell of dimension $d(\sigma)$ with characteristic
map $q|_{\bar e'(\sigma)}$, and the $e(\sigma)$ partition $G_n(\R^m)$.

\step{2} \emph{Boundaries drop dimension.} Let
$X\in\overline{e(\sigma)}\smallsetminus e(\sigma)$, say $X=q(x_1,\dots,x_n)$
with $(x_i)\in\bar e'(\sigma)\smallsetminus e'(\sigma)$. Since
$x_i\in\bar H^{\sigma_i}\subset\R^{\sigma_i}$, we get
$x_1,\dots,x_i\in X\cap\R^{\sigma_i}$ and therefore
\[
\dim\big(X\cap\R^{\sigma_i}\big)\ \ge\ i\qquad\text{for every } i .
\]
If $\tau$ is the symbol of $X$, this says exactly $\tau_i\le\sigma_i$ for all
$i$, whence
\[
d(\tau)=\sum_i(\tau_i-i)\ \le\ \sum_i(\sigma_i-i)=d(\sigma).
\]

\step{3} \emph{Strictness.} Equality would force $\tau_i=\sigma_i$ for all $i$,
i.e.\ $X\in e(\sigma)$, which is excluded. Hence
$\overline{e(\sigma)}\smallsetminus e(\sigma)$ is contained in the union of
cells of strictly smaller dimension, of which there are finitely many.

\step{4} \emph{Limit.} The inclusions $G_n(\R^m)\subset G_n(\R^{m+1})$ are
inclusions of subcomplexes, since a symbol for $\R^m$ is one for $\R^{m+1}$ and
the cells match. The limit therefore inherits a CW structure, with one cell for
each sequence $1\le\sigma_1<\dots<\sigma_n$, that is, for each partition
$0\le\sigma_1-1\le\dots\le\sigma_n-n$ into at most $n$ parts.
\end{proof}

\begin{corollary}\label{cor:HGn-rank}
For every $r$,
\[
\dim_{\Z_2}H^r\big(G_n;\Z_2\big)\ \le\
p_n(r):=\#\{\text{partitions of }r\text{ into at most }n\text{ parts}\}.
\]
\end{corollary}

\begin{proof}
The cellular cochain complex computing $H^*(G_n;\Z_2)$ has exactly $p_n(r)$
generators in degree $r$, by Theorem~\ref{thm:cw-grassmannian}; cohomology of a
complex is a subquotient of the complex.
\end{proof}

\begin{remark}[quoted]
In fact the mod $2$ cellular boundary operator vanishes identically, so the
inequality above is an equality. We shall not need this: the reverse inequality
will come for free from Theorem~\ref{thm:HGn} below. See Milnor--Stasheff,
\S 6, for the direct verification that all incidence numbers of Schubert cells
are even.
\end{remark}

\section{The Leray--Hirsch theorem}

\begin{theorem}[Leray--Hirsch]\label{thm:leray-hirsch}
Let $\pi\colon E\to B$ be a fibre bundle with fibre $F$ and
$i\colon F\hookrightarrow E$ the inclusion of a fibre. Assume that $B$ admits a
finite trivializing cover, that $H^*(F;\Lambda)$ is a free $\Lambda$-module of
finite rank, and that there is a $\Lambda$-linear map
\[
s\colon H^*(F;\Lambda)\longrightarrow H^*(E;\Lambda)
\qquad\text{with}\qquad
i^*\circ s=\mathrm{id}_{H^*(F)} .
\]
Then
\[
\Phi\colon H^*(F)\otimes_\Lambda H^*(B)\longrightarrow H^*(E),
\qquad
\alpha\otimes\beta\longmapsto s(\alpha)\smile\pi^*\beta,
\]
is an isomorphism. In particular $\pi^*\colon H^*(B)\to H^*(E)$ is injective.
\end{theorem}

\begin{proof}
\begin{strategy}
The map $\Phi$ is a homomorphism of $H^*(B)$-modules and is defined globally, so
it can be compared over the members of a cover. Over a trivializing open set
K\"unneth applies, and the matrix of $\Phi$ in the K\"unneth basis is triangular
with respect to the degree of the fibre classes, with invertible diagonal
blocks; hence $\Phi$ is an isomorphism there. Mayer--Vietoris and the five lemma
then propagate the statement from a cover to a union, and induction on the size
of the cover finishes.
\end{strategy}

\step{1} \emph{Restriction.} For open $U\subset B$ let $E_U=\pi^{-1}(U)$ and let
$\Phi_U$ be defined by the same formula, using the restrictions of $s(\alpha)$ to
$E_U$. Restriction of cohomology classes commutes with $\Phi$, so we obtain maps
of long exact sequences whenever we use Mayer--Vietoris.

\step{2} \emph{The trivial case.} Assume $E_U\iso U\times F$. Choose a
homogeneous basis $\alpha_1,\dots,\alpha_N$ of $H^*(F)$ ordered so that
$\deg\alpha_1\ge\dots\ge\deg\alpha_N$. By K\"unneth,
\[
H^*(E_U)\ \iso\ H^*(U)\otimes H^*(F),
\]
a free $H^*(U)$-module with basis $1\otimes\alpha_1,\dots,1\otimes\alpha_N$.
Write
\[
s(\alpha_j)\big|_{E_U}=\sum_{k}c_{jk}\otimes\alpha_k,
\qquad c_{jk}\in H^{\deg\alpha_j-\deg\alpha_k}(U).
\]
Since degrees are non-negative, $c_{jk}=0$ unless
$\deg\alpha_j\ge\deg\alpha_k$, so the matrix $(c_{jk})$ is block triangular for
our ordering. Restricting to a fibre kills all classes of positive degree
pulled back from $U$, so $i^*s=\mathrm{id}$ gives, for
$\deg\alpha_j=\deg\alpha_k$, that $c_{jk}\in H^0(U)$ has the value $\delta_{jk}$
at every point; hence the diagonal blocks are invertible over $H^0(U)$.

\step{3} \emph{Invertibility.} A block triangular matrix over a commutative
ring with invertible diagonal blocks is invertible, because its strictly
triangular part is nilpotent (it strictly decreases the finite ordering). Hence
$(c_{jk})$ is invertible over $H^*(U)$, and therefore
\[
\Phi_U(\alpha_j\otimes\beta)=s(\alpha_j)\smile\pi^*\beta
=\sum_k\big(c_{jk}\,\beta\big)\otimes\alpha_k
\]
is an isomorphism of $H^*(U)$-modules.

\step{4} \emph{Mayer--Vietoris.} Suppose $U=U'\cup U''$ with the theorem known
for $U',U''$ and $U'\cap U''$. Tensoring the Mayer--Vietoris sequence of
$(U',U'')$ with the free module $H^*(F)$ keeps it exact, and $\Phi$ maps it to
the Mayer--Vietoris sequence of $(E_{U'},E_{U''})$, commutatively by Step 1.
Three out of every four consecutive vertical maps are isomorphisms, so the five
lemma gives that $\Phi_U$ is an isomorphism.

\step{5} \emph{Induction.} Let $U_1,\dots,U_r$ be a finite trivializing cover of
$B$. The theorem holds over each $U_i$ by Step 2, hence over
$U_1\cup U_2$ by Step 4, since $U_1\cap U_2$ is trivializing as well. Assuming
it over $U_1\cup\dots\cup U_j$, apply Step 4 to
$U'=U_1\cup\dots\cup U_j$ and $U''=U_{j+1}$, whose intersection is
$\bigcup_{i\le j}(U_i\cap U_{j+1})$, a union of $j$ trivializing sets to which
the inductive hypothesis applies. After $r$ steps we reach $B$.

\step{6} \emph{Injectivity of $\pi^*$.} Taking $\alpha=1\in H^0(F)$ and noting
$s(1)=1$, since $i^*$ is injective in degree $0$ on the relevant summand, the
composite $\beta\mapsto\Phi(1\otimes\beta)=\pi^*\beta$ is the inclusion of a
direct summand, hence injective.
\end{proof}

\begin{remark}
For the infinite Grassmannians we apply the theorem to the compact
approximations $G_n(\R^{n+k})$, which do admit finite trivializing covers, and
pass to the limit; the cohomology of $G_n$ in each fixed degree agrees with that
of $G_n(\R^{n+k})$ for $k$ large, since the CW structures are compatible and no
new cells appear in low dimensions.
\end{remark}

\section{Projective bundles and the splitting principle}

\begin{definition}
For a real bundle $\pi\colon E\to B$ of rank $n$, the \emph{projective bundle}
$p\colon\mathbb P(E)\to B$ is the fibre bundle whose fibre over $b$ is the space
$\mathbb P(\pi^{-1}(b))\iso\RP^{n-1}$ of lines in the fibre.
\end{definition}

\begin{lemma}\label{lem:PE-injective}
For every rank $n$ bundle $E\to B$ over a base with a finite trivializing cover,
\[
H^*\big(\mathbb P(E);\Z_2\big)\ \iso\ H^*(B;\Z_2)\otimes H^*(\RP^{n-1};\Z_2)
\]
as $H^*(B)$-modules; in particular
$p^*\colon H^*(B;\Z_2)\to H^*(\mathbb P(E);\Z_2)$ is injective.
\end{lemma}

\begin{proof}
\begin{strategy}
Produce the splitting $s$ required by Leray--Hirsch from a canonical line bundle
living on $\mathbb P(E)$: its classifying map restricts on each fibre to the
inclusion $\RP^{n-1}\hookrightarrow\RP^\infty$, which is injective on
cohomology in the relevant range.
\end{strategy}

\step{1} \emph{The tautological line bundle.} Over the point
$(b,\ell)\in\mathbb P(E)$, with $\ell\subset\pi^{-1}(b)$ a line, take the line
$\ell$ itself. This defines a line bundle $\gamma\to\mathbb P(E)$; local
triviality is inherited from that of $E$ together with the local triviality of
$\gamma^1$ over $\RP^{n-1}$.

\step{2} \emph{Its classifying map.} Let
$f\colon\mathbb P(E)\to G_1=\RP^\infty$ classify $\gamma$. On a fibre
$F=\mathbb P(\pi^{-1}(b))\iso\RP^{n-1}$ the bundle $\gamma$ restricts to the
canonical line bundle of $\RP^{n-1}$, so $f|_F$ classifies it and is therefore
homotopic to the standard inclusion $j\colon\RP^{n-1}\hookrightarrow\RP^\infty$.

\step{3} \emph{The splitting.} With
$H^*(\RP^\infty;\Z_2)=\Z_2[a]$ and
$H^*(\RP^{n-1};\Z_2)=\Z_2[a]/(a^{n})$, the map $j^*$ sends $a^k$ to $a^k$ and is
an isomorphism in degrees $\le n-1$. Define
\[
s\colon H^*(\RP^{n-1};\Z_2)\to H^*(\mathbb P(E);\Z_2),
\qquad s(a^k)=f^*(a^k)\quad(0\le k\le n-1),
\]
extended linearly. By Step 2, $i_F^*s(a^k)=j^*(a^k)=a^k$, so
$i_F^*\circ s=\mathrm{id}$.

\step{4} Apply Theorem~\ref{thm:leray-hirsch}.
\end{proof}

\begin{definition}
A \emph{splitting map} for a rank $n$ bundle $E\to B$ is a map $f\colon B'\to B$
such that $f^*E\iso L_1\oplus\dots\oplus L_n$ with all $L_i$ line bundles, and
$f^*\colon H^*(B)\to H^*(B')$ injective.
\end{definition}

\begin{corollary}[Splitting principle]\label{cor:splitting}
Every real vector bundle admits a splitting map.
\end{corollary}

\begin{proof}
\begin{strategy}
Split off one line at a time on the projectivization, where the tautological
line bundle is by construction a sub-bundle of the pull-back; the rank drops by
one at each stage and injectivity is preserved because a composite of injections
is injective.
\end{strategy}

\step{1} On $\mathbb P(E)$ the tautological line bundle $\gamma$ of Step 1 above
is a sub-bundle of $p^*E$: over $(b,\ell)$ the fibre of $p^*E$ is
$\pi^{-1}(b)\supset\ell$. Choosing a Euclidean metric,
Theorem~\ref{thm:orthogonal-complement} gives
\[
p^*E\ \iso\ \gamma\oplus\gamma^{\perp},
\qquad \rk\gamma^\perp=n-1 .
\]

\step{2} By Lemma~\ref{lem:PE-injective}, $p^*$ is injective on cohomology.

\step{3} Repeat the construction with $\gamma^{\perp}$ in place of $E$: form
$\mathbb P(\gamma^\perp)$, split off a further line, and so on. After $n-1$
steps the pull-back of $E$ to
\[
B'=\mathbb P\big(\cdots\mathbb P(\gamma^\perp)\cdots\big)
\]
is a sum of $n$ line bundles, and the composite of the projections is injective
on cohomology, being a composite of injective maps.
\end{proof}

Any identity between characteristic classes which holds for sums of line bundles
therefore holds in general. This is the workhorse of the subject.

\section{Cohomology of the Grassmannian}

\begin{definition}
A polynomial $p\in\Z_2[a_1,\dots,a_n]$ is \emph{symmetric} if it is invariant
under all permutations of the variables. The \emph{elementary symmetric
functions} are
\[
\sigma_1=\sum_i a_i,
\qquad
\sigma_2=\sum_{i<j}a_ia_j,
\qquad\dots,\qquad
\sigma_n=a_1a_2\cdots a_n .
\]
\end{definition}

\begin{theorem}[fundamental theorem of symmetric polynomials, quoted]
The subring of symmetric polynomials in $\Z_2[a_1,\dots,a_n]$ is freely
generated by $\sigma_1,\dots,\sigma_n$:
\[
\Z_2[a_1,\dots,a_n]^{S_n}=\Z_2[\sigma_1,\dots,\sigma_n].
\]
\end{theorem}

Recall that $H^*(\RP^\infty;\Z_2)=\Z_2[a]$ and that, by K\"unneth,
\[
H^*\big((\RP^\infty)^n;\Z_2\big)=\Z_2[a_1,\dots,a_n]
\]
with each $a_i$ of degree $1$.

\begin{theorem}\label{thm:HGn}
Let $h_n\colon(\RP^\infty)^n\to G_n$ classify
$\gamma^1\oplus\dots\oplus\gamma^1$, one summand pulled back from each factor.
Then
\[
h_n^*\colon H^*(G_n;\Z_2)\longrightarrow\Z_2[a_1,\dots,a_n]
\]
is injective with image the ring of symmetric polynomials. Consequently
\[
H^*(G_n;\Z_2)=\Z_2\big[\tilde\sigma_1,\dots,\tilde\sigma_n\big],
\qquad
\tilde\sigma_i:=(h_n^*)^{-1}(\sigma_i)\in H^i(G_n;\Z_2),
\]
and $\dim_{\Z_2}H^r(G_n;\Z_2)=p_n(r)$ for every $r$.
\end{theorem}

\begin{proof}
\begin{strategy}
Injectivity comes from the splitting principle: a splitting map for the
universal bundle factors through $h_n$, and an injective map cannot factor
through a non-injective one. Symmetry of the image comes from permuting the
factors, which does not change the bundle being classified. Finally the image is
all of the symmetric polynomials by comparing dimensions with the cell count of
Corollary~\ref{cor:HGn-rank}.
\end{strategy}

\step{1} \emph{Injectivity.} Let $f\colon B'\to G_n$ be a splitting map for
$\gamma^n$ (Corollary~\ref{cor:splitting}), so
$f^*\gamma^n\iso L_1\oplus\dots\oplus L_n$ and $f^*$ is injective. Let
$g\colon B'\to(\RP^\infty)^n$ have $i$-th component a classifying map for
$L_i$; then
\[
g^*\big(\gamma^1\oplus\dots\oplus\gamma^1\big)\iso L_1\oplus\dots\oplus L_n
\iso f^*\gamma^n .
\]
So $f$ and $h_n\circ g$ both classify the same bundle over $B'$, hence are
homotopic by Theorem~\ref{thm:universal}, giving $f^*=g^*\circ h_n^*$. Since
$f^*$ is injective, so is $h_n^*$.

\step{2} \emph{The image is symmetric.} A permutation $\rho$ of the factors of
$(\RP^\infty)^n$ satisfies
$\rho^*(\gamma^1\oplus\dots\oplus\gamma^1)\iso\gamma^1\oplus\dots\oplus\gamma^1$,
since it merely reorders the summands. Hence $h_n\circ\rho\simeq h_n$ and
$\rho^*h_n^*=h_n^*$. As $\rho^*$ permutes $a_1,\dots,a_n$, every element of the
image is invariant under permutations, so
\[
\im h_n^*\ \subset\ \Z_2[\sigma_1,\dots,\sigma_n].
\]

\step{3} \emph{Dimension count.} In degree $r$, the symmetric polynomials have
$\Z_2$-dimension equal to the number of monomials
$\sigma_1^{d_1}\cdots\sigma_n^{d_n}$ with $\sum i\,d_i=r$, that is, to the
number $p_n(r)$ of partitions of $r$ into at most $n$ parts (take $d_i$ parts
equal to $i$). By Steps 1 and 2 together with
Corollary~\ref{cor:HGn-rank},
\[
\dim H^r(G_n;\Z_2)\ \le\ p_n(r)
\qquad\text{and}\qquad
\dim \im h_n^*=\dim H^r(G_n;\Z_2).
\]

\step{4} \emph{Surjectivity onto the symmetric part.} It remains to produce
classes mapping onto $\sigma_1,\dots,\sigma_n$. The mod-$2$ cellular boundary of
the Schubert complex vanishes (quoted above, Milnor--Stasheff \S 6), so equality
holds in Corollary~\ref{cor:HGn-rank}:
$\dim H^r(G_n;\Z_2)=p_n(r)$. Combined with Step 3 this forces
$\im h_n^*$ to be the whole degree-$r$ symmetric part, for every $r$. Hence
$h_n^*$ is an isomorphism onto $\Z_2[\sigma_1,\dots,\sigma_n]$ and we may define
$\tilde\sigma_i$ as its preimage of $\sigma_i$.
\end{proof}

\begin{remark}
Step 4 is the one place where we import a computation. It can be avoided: the
independent construction of Stiefel--Whitney classes by Steenrod squares in
Lecture 7 produces classes $w_i(\gamma^n)\in H^i(G_n;\Z_2)$ with
$h_n^*w_i(\gamma^n)=\sigma_i$, which gives the surjectivity of Step 4 directly,
and then the equality $\dim H^r=p_n(r)$ becomes a corollary rather than an
input. Either route closes the argument; only the order of exposition changes.
\end{remark}

\section{Existence and uniqueness of the SW classes}

\begin{theorem}
There is exactly one family of classes satisfying the axioms of
Definition~\ref{def:sw-axioms}, namely
\[
w_i(\xi)=f_\xi^*\big(\tilde\sigma_i\big),
\]
where $f_\xi\colon B\to G_n$ is a classifying map for $\xi$.
\end{theorem}

\begin{proof}
\begin{strategy}
Existence: three of the four axioms are immediate from the definition, and the
fourth --- the product formula --- is proved in Lecture 6 by computing the
effect on cohomology of the map $G_n\times G_{n'}\to G_{n+n'}$ classifying the
direct sum. Uniqueness: naturality reduces everything to the universal bundle,
the splitting principle reduces the universal bundle to sums of line bundles,
and the product formula reduces sums of line bundles to a single line bundle,
where normalization pins the answer down.
\end{strategy}

\step{1} \emph{Rank.} $\tilde\sigma_0=1$ and $H^*(G_n)$ has no generators beyond
$\tilde\sigma_n$, so $w_0=1$ and $w_i=0$ for $i>n$ by definition.

\step{2} \emph{Naturality.} If $\eta=g^*\xi$ then $f_\eta\simeq f_\xi\circ g$ by
Theorem~\ref{thm:universal}, so
$w_i(\eta)=f_\eta^*\tilde\sigma_i=g^*f_\xi^*\tilde\sigma_i=g^*w_i(\xi)$.

\step{3} \emph{Normalization.} The inclusion $\RP^1\hookrightarrow\RP^\infty=G_1$
classifies $\gamma^1_1$, and $\tilde\sigma_1=a$ generates
$H^1(\RP^\infty;\Z_2)$, whose restriction to $H^1(\RP^1;\Z_2)$ is the generator;
hence $w_1(\gamma^1_1)\neq0$.

\step{4} \emph{Product formula.} Proved in Theorem~\ref{thm:whitney} of the next
lecture.

\step{5} \emph{Uniqueness.} Let $w_i'$ be any family satisfying the axioms. By
naturality both families are determined by their values on the universal bundle
$\gamma^n$, so it suffices to prove $w_i(\gamma^n)=w_i'(\gamma^n)$. Since
$h_n^*$ is injective (Theorem~\ref{thm:HGn}), it suffices to compare after
applying $h_n^*$, that is, to compare the classes of
$\gamma^1\oplus\dots\oplus\gamma^1$ over $(\RP^\infty)^n$. By the product
formula, both are determined by the classes of a single $\gamma^1\to\RP^\infty$.
There only $w_1$ is in question, by the rank axiom, and naturality applied to
$\gamma^1_1\to\gamma^1$ together with the normalization axiom forces
$w_1'(\gamma^1)$ to be non-zero, hence equal to the generator $a$, which is
$w_1(\gamma^1)$.
\end{proof}

\lecture

\begin{guiding}
One axiom is still owed: the Whitney product formula. The universal way to prove
an identity between characteristic classes is to verify it on universal
examples, where cohomology is a polynomial ring and everything reduces to
symmetric functions. The same circle of ideas then produces the \emph{Chern
classes} of complex bundles, with integer rather than mod-$2$ coefficients.
\end{guiding}

\section{Proof of the Whitney product formula}

Recall from Lecture 5 that $H^*(G_n;\Z_2)=\Z_2[\tilde\sigma_1,\dots,
\tilde\sigma_n]$, embedded by $h_n^*$ as the symmetric polynomials in
$\Z_2[a_1,\dots,a_n]$, and that $w_i(E)=f_E^*(\tilde\sigma_i)$.

Write $G=G_n$, $G'=G_{n'}$, $G''=G_{n+n'}$, and let
\[
F\colon G\times G'\longrightarrow G''
\]
classify $p_1^*\gamma^n\oplus p_2^*\gamma^{n'}$, where $p_1,p_2$ are the two
projections.

\begin{lemma}\label{lem:F-star}
$\displaystyle F^*(\tilde\sigma_i)
=\sum_{s=0}^{i}\tilde\sigma'_s\times\tilde\sigma''_{i-s}$,
where $\tilde\sigma'$ and $\tilde\sigma''$ denote the generators of $H^*(G)$ and
$H^*(G')$.
\end{lemma}

\begin{proof}
\begin{strategy}
Both sides live in $H^*(G\times G')$, where $(h_n\times h_{n'})^*$ is injective;
so it is enough to compare their images in a polynomial ring. There the identity
becomes the classical decomposition of an elementary symmetric function in
$n+n'$ variables into products of elementary symmetric functions in the two
groups of variables.
\end{strategy}

\step{1} \emph{A homotopy commutative square.} Let
\[
\tilde F\colon(\RP^\infty)^n\times(\RP^\infty)^{n'}
\longrightarrow(\RP^\infty)^{n+n'}
\]
be the identity under the evident identification. Both composites in
\[
\begin{tikzcd}[row sep=large, column sep=large]
(\RP^\infty)^n\times(\RP^\infty)^{n'}
\arrow[r,"\tilde F"] \arrow[d,"h_n\times h_{n'}"']
& (\RP^\infty)^{n+n'} \arrow[d,"h_{n+n'}"]\\
G\times G' \arrow[r,"F"] & G''
\end{tikzcd}
\]
classify $\gamma^1\oplus\dots\oplus\gamma^1$ with $n+n'$ summands, hence are
homotopic by Theorem~\ref{thm:universal}. Therefore
\[
(h_n\times h_{n'})^*\circ F^* \;=\; \tilde F^*\circ h_{n+n'}^* .
\]

\step{2} \emph{Computing the right-hand side.} By definition we have
$h_{n+n'}^*(\tilde\sigma_i)=\sigma_i(a_1,\dots,a_{n+n'})$. Under $\tilde F^*$
the first $n$ variables become the variables $a'_1,\dots,a'_n$ of the first
factor, and the last $n'$ become $a''_1,\dots,a''_{n'}$. Sorting a monomial
$a_{j_1}\cdots a_{j_i}$ by how many of its indices are at most $n$, we obtain
\[
\sigma_i(a',a'')=\sum_{s=0}^{i}\sigma_s(a')\,\sigma_{i-s}(a'').
\]

\step{3} \emph{Computing the left-hand side.} By K\"unneth,
$(h_n\times h_{n'})^*(x\times y)=h_n^*(x)\times h_{n'}^*(y)$, so
\[
(h_n\times h_{n'})^*\Big(\sum_s\tilde\sigma'_s\times\tilde\sigma''_{i-s}\Big)
=\sum_s\sigma_s(a')\,\sigma_{i-s}(a'').
\]

\step{4} \emph{Conclusion.} The two expressions agree, and
$(h_n\times h_{n'})^*$ is injective, being a tensor product of injective maps
between free modules (Theorem~\ref{thm:HGn} and K\"unneth). Hence
$F^*(\tilde\sigma_i)=\sum_s\tilde\sigma'_s\times\tilde\sigma''_{i-s}$.
\end{proof}

\begin{theorem}[Whitney product formula]\label{thm:whitney}
For vector bundles $E,E'$ over the same base $B$,
\[
w_i(E\oplus E')=\sum_{s=0}^{i}w_s(E)\smile w_{i-s}(E').
\]
\end{theorem}

\begin{proof}
\begin{strategy}
Write the classifying map of $E\oplus E'$ as the composite of $F$ with
$f_E\times f_{E'}$ and the diagonal, then feed in Lemma~\ref{lem:F-star} and
convert cross products into cup products.
\end{strategy}

\step{1} \emph{The classifying map of the sum.} Let $d\colon B\to B\times B$ be
the diagonal. The bundle $(f_E\times f_{E'})^*\big(p_1^*\gamma^n\oplus
p_2^*\gamma^{n'}\big)$ over $B\times B$ restricts along $d$ to $E\oplus E'$, so
$F\circ(f_E\times f_{E'})\circ d$ is a classifying map for $E\oplus E'$.

\step{2} \emph{Applying the lemma.}
\[
w_i(E\oplus E')
=d^*(f_E\times f_{E'})^*F^*(\tilde\sigma_i)
=d^*(f_E\times f_{E'})^*\Big(\sum_s\tilde\sigma'_s\times\tilde\sigma''_{i-s}\Big).
\]

\step{3} \emph{Naturality of the cross product.}
$(f_E\times f_{E'})^*(x\times y)=f_E^*x\times f_{E'}^*y$, so the expression
becomes $d^*\big(\sum_s w_s(E)\times w_{i-s}(E')\big)$.

\step{4} \emph{Diagonal.} Finally $d^*(x\times y)=x\smile y$, which gives the
formula.
\end{proof}

This completes the axiomatic circle begun in Lecture 2: the Stiefel--Whitney
classes exist, satisfy the four axioms, and are the only classes that do.

\section{Complex vector bundles}

\begin{guiding}
What changes when the fibres are complex vector spaces? Two things, both to our
advantage: orientations come for free, and the relevant Grassmannians have cells
only in even dimensions, so the theory can be run over $\Z$ instead of $\Z_2$.
\end{guiding}

\begin{definition}
A \emph{complex vector bundle} of rank $n$ over $B$ is a map $\pi\colon E\to B$
in which every fibre carries the structure of a complex vector space and local
trivializations $\pi^{-1}(U)\iso U\times\C^n$ can be chosen complex linear on
each fibre. All the constructions of Lecture 1 --- pull-back, Whitney sum,
sub-bundles --- carry over verbatim.
\end{definition}

\begin{remark}
A complex bundle of rank $n$ is in particular a real bundle of rank $2n$.
Conversely a real $\R^{2n}$-bundle $E$ can be promoted to a complex bundle
exactly when it admits a \emph{complex structure}: a fibrewise linear
$J\colon E\to E$ with $J(J(v))=-v$, scalar multiplication being defined by
$(a+ib)v=av+bJv$. For instance $TS^2$ admits one, rotation by a quarter turn,
whereas the tangent bundles of most even-dimensional spheres do not.
\end{remark}

\begin{definition}
A \emph{Hermitian metric} on a complex bundle is a Euclidean metric on the
underlying real bundle satisfying $|iv|=|v|$. It determines a Hermitian inner
product $\langle\cdot,\cdot\rangle$, complex linear in the first variable and
conjugate linear in the second, and hence a notion of orthogonality. If
$L\subset E$ is a complex sub-bundle then $L^\perp$ is again complex and
$L\oplus L^\perp\iso E$ as complex bundles, by the argument of
Theorem~\ref{thm:orthogonal-complement} applied with complex scalars.
\end{definition}

\begin{definition}
The \emph{complex Grassmannian} $G_n(\C^m)$ is the space of complex $n$-planes
in $\C^m$, topologized as the quotient of the complex Stiefel manifold
$V_n(\C^m)$ of complex $n$-frames. It is a compact manifold of real dimension
$2n(m-n)$, and $G_1(\C^m)=\CP^{m-1}$. The tautological bundle
$\gamma_n\to G_n(\C^m)$, the limit $G_n(\C^\infty)$ and its tautological bundle
are defined exactly as in the real case, and the classification theorem holds
with the same proof: complex $n$-bundles over a paracompact base $B$ correspond
to homotopy classes $[B,G_n(\C^\infty)]$.
\end{definition}

\begin{theorem}\label{thm:complex-cells}
$G_n(\C^m)$ is a CW complex with one cell of real dimension
$2\sum_i(\sigma_i-i)$ for each Schubert symbol $\sigma$. All cells are
even-dimensional, so the cellular boundary operator vanishes and
\[
H^{2r}\big(G_n(\C^\infty);\Z\big)\iso\Z^{\,p_n(r)},
\qquad
H^{\mathrm{odd}}\big(G_n(\C^\infty);\Z\big)=0,
\]
where $p_n(r)$ is the number of partitions of $r$ into at most $n$ parts.
\end{theorem}

\begin{proof}
\begin{strategy}
Repeat the Schubert analysis over $\C$. The only change is in the normalization
of a basis vector: a complex line contains a whole circle of unit vectors rather
than two, so requiring the last non-zero coordinate to be real and positive cuts
the relevant sphere down by one further dimension. The rotation argument of
Lecture 5 goes through unchanged, and the parity of the resulting dimensions
makes the chain complex collapse.
\end{strategy}

\step{1} \emph{Symbols.} For $X\in G_n(\C^m)$ the sequence
$\dim_\C(X\cap\C^k)$ increases from $0$ to $n$ in steps of at most $1$, by the
same exactness argument as in Claim~\ref{claim:jumps}, so it determines a symbol
$1\le\sigma_1<\dots<\sigma_n\le m$ and a partition of $G_n(\C^m)$ into sets
$e(\sigma)$.

\step{2} \emph{Normalized bases.} Put
\[
H^{\sigma_i}_\C=\big\{(x_1,\dots,x_{\sigma_i},0,\dots,0)\in\C^m
\ \big|\ x_{\sigma_i}\in\R_{>0}\big\}.
\]
Given $x$ in a complex line, the vectors $e^{i\theta}x$ all have the same norm,
so there is a unique unit vector on the line whose $\sigma_i$-th coordinate is
real and positive. Repeating the induction of Claim~\ref{claim:half-space} with
complex orthogonality gives: $X\in e(\sigma)$ if and only if $X$ has a unique
unitary basis $(x_1,\dots,x_n)$ with $x_i\in H^{\sigma_i}_\C$.

\step{3} \emph{Cell dimensions.} For $n=1$ the closure of
$H^{\sigma_1}_\C\cap S^{2\sigma_1-1}$ consists of the unit vectors of
$\C^{\sigma_1}$ with last coordinate in $\R_{\ge0}$; this is a closed cell of
real dimension $2\sigma_1-2$, since the last coordinate contributes one real
parameter instead of two. For the inductive step, the space of admissible new
vectors is the set of unit vectors in $\bar H^{\sigma_{n+1}}_\C$ orthogonal to
$n$ fixed ones, a closed cell of real dimension
$2(\sigma_{n+1}-n)-2$; the map $\Psi$ built from unitary rotations
$T(u,v)$ as in Lemma~\ref{lem:closed-cell} identifies the product with
$\bar e'(\sigma_1,\dots,\sigma_{n+1})$. Adding up gives real dimension
$2\sum_i(\sigma_i-i)$.

\step{4} \emph{Boundaries.} The argument of
Theorem~\ref{thm:cw-grassmannian} shows that the boundary of a cell lies in
cells of strictly smaller dimension, so we do obtain a CW structure.

\step{5} \emph{Cohomology.} Since all cells have even dimension, the cellular
chain complex has zero groups in odd degrees, hence all differentials vanish and
$H_{2r}$ is free on the $2r$-cells. Counting cells as in Lecture 5 gives
$p_n(r)$ of them, and universal coefficients gives the statement for
cohomology.
\end{proof}

\section{Chern classes}

By K\"unneth, $H^*\big((\CP^\infty)^n;\Z\big)=\Z[a_1,\dots,a_n]$ with $a_i$ of
degree $2$.

\begin{theorem}\label{thm:HGnC}
Let $h_n\colon(\CP^\infty)^n\to G_n(\C^\infty)$ classify
$\gamma_1\oplus\dots\oplus\gamma_1$. Then $h_n^*$ is injective with image the
symmetric polynomials, so
\[
H^*\big(G_n(\C^\infty);\Z\big)=\Z\big[\tilde\sigma_1,\dots,\tilde\sigma_n\big],
\qquad \tilde\sigma_i\in H^{2i}.
\]
\end{theorem}

\begin{proof}
The proof of Theorem~\ref{thm:HGn} applies word for word, with $\RP^\infty$
replaced by $\CP^\infty$, degrees doubled, and $\Z_2$ replaced by $\Z$. The
three ingredients are available: the splitting principle over $\C$
(Remark~\ref{rem:complex-splitting} below), the invariance of the image under
permutation of the factors, and the cell count of
Theorem~\ref{thm:complex-cells}, which here is an equality without any imported
input, since the chain complex has no odd cells and hence zero differentials.
\end{proof}

\begin{definition}[Chern classes]
For a complex bundle $E\to B$ of rank $n$ with classifying map
$f_E\colon B\to G_n(\C^\infty)$, set
\[
c_i(E):=f_E^*\big(\tilde\sigma_i\big)\ \in\ H^{2i}(B;\Z),
\qquad c_0=1,\qquad c_i=0\ \text{ for } i>n .
\]
The \emph{total Chern class} is $c(E)=1+c_1(E)+\dots+c_n(E)$.
\end{definition}

\begin{theorem}
The Chern classes are the unique family of classes satisfying:

\begin{enumerate}
\item[(C1)] $c_0(E)=1$ and $c_i(E)=0$ for $i>\rk_\C E$;
\item[(C2)] $c_i(f^*E)=f^*c_i(E)$;
\item[(C3)] $c_k(E_1\oplus E_2)=\sum_{j}c_j(E_1)\smile c_{k-j}(E_2)$;
\item[(C4)] $c_1(\gamma_1\to\CP^1)$ is a generator of $H^2(\CP^1;\Z)$.
\end{enumerate}
\end{theorem}

\begin{proof}
Every step of Lectures 5 and 6 transposes: (C1) and (C2) hold by definition and
by homotopy invariance of classifying maps; (C3) is proved exactly as
Theorem~\ref{thm:whitney}, using the complex analogue of
Lemma~\ref{lem:F-star}, whose proof is the same decomposition of elementary
symmetric functions; and (C4) holds because $\tilde\sigma_1$ restricts to the
generator on $G_1(\C^\infty)=\CP^\infty$ and then to $\CP^1$. Uniqueness follows
as before: naturality reduces to the universal bundle, the complex splitting
principle to sums of line bundles, (C3) to a single line bundle, and (C4) pins
that case down.
\end{proof}

\begin{lemma}\label{lem:pi-star-injective}
Let $E\to B$ be a complex bundle of rank $n$ over a base with a finite
trivializing cover, and let $\mathbb P(E)\to B$ be its complex projectivization,
with fibres $\CP^{n-1}$. Then
\[
H^*\big(\mathbb P(E);\Z\big)\iso H^*(B;\Z)\otimes H^*(\CP^{n-1};\Z),
\]
and in particular $H^*(B;\Z)\to H^*(\mathbb P(E);\Z)$ is injective.
\end{lemma}

\begin{proof}
Identical to Lemma~\ref{lem:PE-injective}. The tautological complex line bundle
$\gamma$ over $\mathbb P(E)$ restricts on each fibre to the canonical bundle of
$\CP^{n-1}$, so its classifying map
$f\colon\mathbb P(E)\to\CP^\infty$ restricts on the fibre to a map homotopic to
the standard inclusion, which is an isomorphism on $H^{2k}$ for $k\le n-1$.
Setting $s(a^k)=f^*(a^k)$ gives $i_F^*\circ s=\mathrm{id}$, and
Theorem~\ref{thm:leray-hirsch} applies over $\Lambda=\Z$, since
$H^*(\CP^{n-1};\Z)$ is free of finite rank.
\end{proof}

\begin{remark}\label{rem:complex-splitting}
Consequently every complex bundle admits a splitting map: on $\mathbb P(E)$ the
tautological line bundle is a complex sub-bundle of $p^*E$, and a Hermitian
metric splits $p^*E\iso\gamma\oplus\gamma^\perp$ with $\gamma^\perp$ of rank
$n-1$; iterate as in Corollary~\ref{cor:splitting}. Formally one writes
\[
c(E)=\prod_{i=1}^{n}(1+x_i),
\qquad x_i=c_1(L_i),
\]
and calls the $x_i$ the \emph{Chern roots} of $E$. Every identity between Chern
classes may be verified on these roots.
\end{remark}

In the next lecture we change viewpoint entirely and construct the
Stiefel--Whitney classes a second time, intrinsically, out of Steenrod squares
and the Thom isomorphism. That route also hands us the Euler class, the integral
refinement of the top class.

\lecture

\begin{guiding}
Milnor and Stasheff \emph{define} the Stiefel--Whitney classes by the formula
$w_i(\xi)=\phi^{-1}\Sq^i\phi(1)$. What are the two ingredients $\Sq^i$ and
$\phi$, why does the formula satisfy our axioms, and what integral information
appears along the way?
\end{guiding}

\section{Steenrod squares}

\begin{definition}
The $n$-th homotopy group of a space $X$ is $\pi_n(X)=[S^n,X]$. An
\emph{Eilenberg--MacLane space} $K(G,n)$ is a space with $\pi_n=G$ and all other
homotopy groups trivial.
\end{definition}

\begin{theorem}[representability, quoted]
There is a natural bijection
\[
B\colon\langle X,K(G,n)\rangle\longrightarrow H^n(X;G),
\qquad
[f]\longmapsto f^*\alpha,
\]
where $\langle\ ,\ \rangle$ denotes homotopy classes of maps and
$\alpha\in H^n(K(G,n);G)$ is a fixed fundamental class.
\end{theorem}

We sketch the construction of the squares, following the course. Fix
$\alpha\in H^n(X;\Z_2)$ and let $f_\alpha\colon X\to K(\Z_2,n)$ represent it.
The cross square $\alpha\times\alpha\in H^{2n}(X\times X;\Z_2)$ is represented by
a map $f_{\alpha\times\alpha}$, and the twist $T(x,y)=(y,x)$ satisfies
$T^*(\alpha\times\alpha)=\alpha\times\alpha$ over $\Z_2$, so
\[
f_{\alpha\times\alpha}\circ T\ \simeq\ f_{\alpha\times\alpha}.
\]
Concatenating such a homotopy with its reverse gives a map
$\tilde F\colon X\times X\times S^1\to K(\Z_2,2n)$ which is null-homotopic on
$X\times X\times\{\mathrm{pt}\}$; extending over $D^2$, composing with $T$ and
gluing produces maps on $X\times X\times S^2$, and iterating yields
\[
H\colon X\times X\times S^\infty\longrightarrow K(\Z_2,2n),
\qquad
H(x,y,s)=H(y,x,-s).
\]
Restricting to the diagonal $x=y$, the symmetry makes $H$ factor through
$X\times\RP^\infty$, so we obtain a class in $H^{2n}(X\times\RP^\infty;\Z_2)$
which we expand using $H^*(\RP^\infty;\Z_2)=\Z_2[\omega]$:
\[
H_\Delta^*(\alpha)=\sum_{i=0}^{2n}\alpha_i\times\omega^{\,2n-i},
\qquad \alpha_i\in H^i(X;\Z_2).
\]

\begin{definition}
$\Sq^i(\alpha):=\alpha_{n+i}\in H^{n+i}(X;\Z_2)$.
\end{definition}

\begin{theorem}[properties, quoted]\label{thm:sq-properties}
The operations
\[
\Sq^i\colon H^n(X;\Z_2)\longrightarrow H^{n+i}(X;\Z_2)
\]
are natural homomorphisms satisfying

\begin{enumerate}
\item $\Sq^i(f^*\alpha)=f^*(\Sq^i\alpha)$ and
$\Sq^i(\alpha+\beta)=\Sq^i\alpha+\Sq^i\beta$;
\item $\Sq^0=\mathrm{id}$; for $\alpha\in H^n$, $\Sq^n\alpha=\alpha\smile\alpha$
and $\Sq^i\alpha=0$ for $i>n$;
\item the Cartan formula
$\Sq^i(\alpha\smile\beta)=\sum_j\Sq^j\alpha\smile\Sq^{i-j}\beta$.
\end{enumerate}
\end{theorem}

The verification of these properties from the construction above is a long
homotopy-theoretic argument which we do not reproduce; see Milnor--Stasheff
\S 8 or Steenrod--Epstein.

\section{Orientations}

\begin{definition}
An \emph{orientation} of a real $n$-dimensional vector space $V$ is an
equivalence class of bases, two bases being equivalent when the transition
matrix has positive determinant. There are exactly two such classes.
\end{definition}

A linearly embedded simplex $\Delta^n\subset V$ with ordered vertices
$v_0,\dots,v_n$ determines the orientation given by
$\overrightarrow{v_0v_1},\dots,\overrightarrow{v_{n-1}v_n}$. Writing
$V_0=V\smallsetminus\{0\}$ and assuming $0$ lies in the interior of the simplex,
the map $\sigma\colon\Delta^n\to V$ represents a generator of
$H_n(V,V_0)\iso\Z$: this is the \emph{preferred generator} $\mu_V$ attached to
the orientation. Dually there is a preferred $u_V\in H^n(V,V_0)$ with
$\langle u_V,\mu_V\rangle=1$.

\begin{definition}
An \emph{orientation of a vector bundle} $\xi$ is a compatible choice of
orientation on every fibre: local trivializations can be chosen orientation
preserving on each fibre, equivalently all transition matrices have positive
determinant. It determines a preferred generator $u_F\in H^n(F,F_0)$ on every
fibre pair.
\end{definition}

\section{The Thom isomorphism}

\begin{theorem}[Thom isomorphism theorem]\label{thm:thom}
Let $\pi\colon E\to B$ be a vector bundle of rank $n$, oriented if we use $\Z$
coefficients; over $\Z_2$ no hypothesis is needed. Write
$E_0=E\smallsetminus(\text{zero section})$. Then

\begin{enumerate}
\item there is a unique class $u\in H^n(E,E_0)$, the \emph{Thom class}, whose
restriction to every fibre pair $(F,F_0)$ is the preferred generator;
\item $H^i(E,E_0)=0$ for $i<n$, and
\[
\phi\colon H^i(B)\ \xrightarrow{\ \pi^*\ }\ H^i(E)
\ \xrightarrow{\ \smile\, u\ }\ H^{i+n}(E,E_0)
\]
is an isomorphism for every $i$.
\end{enumerate}
\end{theorem}

The proof occupies the second half of Lecture 8. Here we record the formal
consequences, which are what the applications actually use.

\begin{theorem}\label{thm:thom-properties}
Let $u,\phi$ be as above.

\begin{enumerate}
\item \emph{Pull-back.} If $f\colon\xi\to\xi'$ is an orientation preserving
bundle map then $f^*u'=u$.
\item \emph{Naturality.} $\phi_\xi\circ\bar f^{\,*}=f^*\circ\phi_{\xi'}$.
\item \emph{Cross product.} $u_{\xi\times\xi'}=u_\xi\times u_{\xi'}$ and
$\phi_{\xi\times\xi'}(a\times b)=\phi_\xi(a)\times\phi_{\xi'}(b)$.
\end{enumerate}
\end{theorem}

\begin{proof}
\begin{strategy}
All three are consequences of the uniqueness clause in
Theorem~\ref{thm:thom}: to identify a class in $H^n(E,E_0)$ it is enough to
check what it restricts to on one fibre pair.
\end{strategy}

\step{1} \emph{Pull-back.} The map $f$ induces
$f^*\colon H^n(E',E'_0)\to H^n(E,E_0)$. On a fibre pair, $f$ restricts to an
orientation preserving isomorphism $(F,F_0)\to(F',F'_0)$, which carries the
preferred generator to the preferred generator. Hence $f^*u'$ restricts to the
preferred generator on every fibre of $\xi$, and uniqueness gives $f^*u'=u$.

\step{2} \emph{Naturality.} For $y\in H^*(B')$,
\[
f^*\phi_{\xi'}(y)=f^*\big(\pi'^*y\smile u'\big)
=\pi^*\bar f^{\,*}y\smile f^*u'
=\pi^*\bar f^{\,*}y\smile u
=\phi_\xi\big(\bar f^{\,*}y\big),
\]
using that $f^*$ is a ring map compatible with the relative cup product,
$\pi'\circ f=\bar f\circ\pi$, and Step 1.

\step{3} \emph{Cross product.} The class $u_\xi\times u_{\xi'}$ lies in
$H^{n+m}\big(E\times E',(E\times E')_0\big)$, where
$(E\times E')_0=E_0\times E'\cup E\times E'_0$. Restrict it to a fibre pair
$(F^{n+m},F^{n+m}_0)\iso(F^n\times F'^m,\dots)$: by the definition of the cross
product and the product orientation,
\[
\big(u_\xi\times u_{\xi'}\big)\big(\Delta^{n}\times\Delta^{m}\big)
=u_\xi(\Delta^n)\cdot u_{\xi'}(\Delta^m)=1 ,
\]
so it is the preferred generator there. Uniqueness gives
$u_{\xi\times\xi'}=u_\xi\times u_{\xi'}$, and the formula for $\phi$ follows by
expanding both sides.
\end{proof}

\begin{definition}[Stiefel--Whitney classes, intrinsic definition]
For an $\R^n$-bundle $\xi$ with $\Z_2$ Thom class $u=\phi(1)$, put
\[
w_i(\xi):=\phi^{-1}\big(\Sq^i u\big).
\]
\end{definition}

\begin{proposition}
These classes satisfy the axioms of Definition~\ref{def:sw-axioms}, hence agree
with the classes constructed in Lecture 5.
\end{proposition}

\begin{proof}
\begin{strategy}
Match each axiom with the corresponding property of $\Sq$ from
Theorem~\ref{thm:sq-properties} and of $\phi$ from
Theorem~\ref{thm:thom-properties}; then invoke the uniqueness theorem of
Lecture 5.
\end{strategy}

\step{1} \emph{Rank.} $\Sq^0u=u$ gives $w_0=1$. Since $u$ has degree $n$,
property (2) gives $\Sq^iu=0$ for $i>n$, so $w_i=0$ for $i>n$.

\step{2} \emph{Naturality.} For a bundle map $f$ with base map $\bar f$, using
naturality of $\Sq$ and Theorem~\ref{thm:thom-properties},
\[
w_i(\xi)=\phi_\xi^{-1}\Sq^i u=\phi_\xi^{-1}\Sq^i f^*u'
=\phi_\xi^{-1}f^*\Sq^iu'=\bar f^{\,*}\phi_{\xi'}^{-1}\Sq^iu'
=\bar f^{\,*}w_i(\xi').
\]

\step{3} \emph{Whitney formula.} By Theorem~\ref{thm:thom-properties}(3),
$u_{\xi\times\xi'}=u_\xi\times u_{\xi'}$, and the Cartan formula applied to a
cross product gives
$\Sq^i(u_\xi\times u_{\xi'})=\sum_j\Sq^ju_\xi\times\Sq^{i-j}u_{\xi'}$. Applying
$\phi^{-1}_{\xi\times\xi'}$ and using the product formula for $\phi$ yields
$w_i(\xi\times\xi')=\sum_j w_j(\xi)\times w_{i-j}(\xi')$; pulling back along the
diagonal converts this into the Whitney formula for $\xi\oplus\xi'$.

\step{4} \emph{Normalization.} For $\gamma^1_1\to\RP^1$ the pair
$(E,E_0)$ is the M\"obius band relative to its boundary; there
$\Sq^1u=u\smile u\neq0$, as computed in Lecture 8, so $w_1(\gamma^1_1)\neq0$.

\step{5} By the uniqueness theorem of Lecture 5, these classes coincide with
$f_\xi^*\tilde\sigma_i$.
\end{proof}

\section{The Euler class}

\begin{guiding}
The Thom class lives in the relative group $H^n(E,E_0)$. What invariant of the
base do we obtain by simply forgetting that it is relative?
\end{guiding}

\begin{definition}
Let $\xi$ be an oriented $\R^n$-bundle over $B$ with Thom class $u$. The
\emph{Euler class} is
\[
e(\xi):=(\pi^*)^{-1}\,i^*u\ \in\ H^n(B;\Z),
\]
where $i^*\colon H^n(E,E_0)\to H^n(E)$ forgets the pair and
$\pi^*\colon H^*(B)\to H^*(E)$ is an isomorphism because $\pi$ is a homotopy
equivalence.
\end{definition}

\begin{theorem}\label{thm:euler-properties}
The Euler class has the following properties.

\begin{enumerate}
\item \emph{Naturality:} $\bar f^{\,*}e(\xi')=e(\xi)$ for an orientation
preserving bundle map $f\colon\xi\to\xi'$.
\item If $\xi$ is trivial then $e(\xi)=0$.
\item Reversing the orientation of $\xi$ changes the sign of $e(\xi)$;
consequently $2e(\xi)=0$ when $n$ is odd.
\item $e(\xi)\bmod2=w_n(\xi)$.
\item $e(\xi\times\xi')=e(\xi)\times e(\xi')$ and
$e(\xi\oplus\xi')=e(\xi)\smile e(\xi')$.
\item If $\xi$ has a nowhere zero section then $e(\xi)=0$.
\end{enumerate}
\end{theorem}

\begin{proof}
\step{1} Apply $i^*$ to $f^*u'=u$ and use
$\pi^*\circ\bar f^{\,*}=f^*\circ\pi'^*$.

\step{2} A trivial bundle admits a bundle map to the bundle
$\mathrm{pt}\times\R^n$ over a point, whose Euler class lies in
$H^{n}(\mathrm{pt})=0$ for $n>0$; now use (1).

\step{3} The map $(x,v)\mapsto(x,-v)$ is a bundle map over the identity which
reverses the orientation of every fibre when $n$ is odd, hence carries $u$ to
$-u$ and $e$ to $-e$. Naturality then gives $e(\xi)=-e(\xi)$.

\step{4} Over $\Z_2$ we have $i^*u=\phi^{-1}$ applied to $u\smile u=\Sq^nu$ up
to the identification $\phi$, so
\[
e(\xi)\bmod2=\phi^{-1}\big(u\smile u\big)=\phi^{-1}\Sq^n\phi(1)=w_n(\xi).
\]

\step{5} By Theorem~\ref{thm:thom-properties}(3),
$u_{\xi\times\xi'}=u_\xi\times u_{\xi'}$; applying $i^*$, which is compatible
with cross products, and then $(\pi^*)^{-1}$, gives
$e(\xi\times\xi')=e(\xi)\times e(\xi')$. Since
$\xi\oplus\xi'=\Delta^*(\xi\times\xi')$ and
$\Delta^*(x\times y)=x\smile y$, the second formula follows from (1).

\step{6} A nowhere zero section spans a trivial line sub-bundle $\ell\subset\xi$;
with a metric, $\xi\iso\ell\oplus\ell^{\perp}$, so by (5) and (2),
$e(\xi)=e(\ell)\smile e(\ell^\perp)=0$.
\end{proof}

\begin{corollary}
If $2e(\xi)\neq0$ then $\xi$ does not split off an oriented summand of odd rank;
in particular it has no nowhere zero section.
\end{corollary}

In Lecture 8 we compute Euler classes of the canonical bundles, obtain the
universal identity $e(\gamma^n)=\tilde\sigma_n$, and then prove
Theorem~\ref{thm:thom}.

\lecture

\begin{guiding}
Two tasks. First, ground the Euler class in examples: canonical bundles over
projective lines, and the universal identity $e(\gamma^n)=\tilde\sigma_n$.
Second, pay the outstanding debt and prove the Thom isomorphism theorem.
\end{guiding}

\section{Euler classes of canonical bundles}

\begin{example}[The M\"obius band]\label{ex:mobius-euler}
Consider $\RP^1\hookrightarrow\RP^2$, $[x_0,x_1]\mapsto[x_0,x_1,0]$. The
complement of $\RP^1$ is the affine chart $\{x_2\neq0\}\iso\R^2$, an open disc,
and a neighbourhood of $\RP^1$ in $\RP^2$ is the total space of its normal
bundle, which is $\gamma^1_1$. Hence
\[
\RP^2=\big(\text{M\"obius band}\big)\cup_{\partial}D^2 .
\]

\begin{center}
\begin{tikzpicture}[scale=0.9]
\draw[thick] (0,0) circle (1.35);
\draw[thick,dashed] (0,0) circle (0.6);
\node at (0,0) {\small $D^2$};
\node at (0,0.98) {\small M\"ob};
\node at (3.2,0) {\small $\RP^2=\mathrm{M\ddot ob}\cup_{S^1}D^2$};
\end{tikzpicture}
\end{center}

The pair $(E,E_0)$ for $\gamma^1_1$ deformation retracts onto (M\"obius band,
boundary), so with $\Z_2$ coefficients
\[
H^{*>0}\big(\gamma^1_1,(\gamma^1_1)_0\big)
\iso H^{*>0}\big(\RP^2\big)=\Z_2[a]/(a^3),
\]
by excising the complementary disc. Under this identification the Thom class $u$
corresponds to $a$, so $u\smile u$ corresponds to $a^2\neq0$. This is the
computation quoted in Step 4 of the previous lecture, and it gives
$w_1(\gamma^1_1)\neq0$ once more, now by Thom-theoretic means.
\end{example}

\begin{example}[The complex canonical bundle]\label{ex:e-gamma-CP1}
Let $\gamma_1\to\CP^1$ be the tautological complex line bundle. Exactly as above,
$\CP^1\subset\CP^2$ has normal bundle $\gamma_1$ and complementary affine chart
a $4$-disc, so
\[
\CP^2=E\big(\gamma_1\big)\cup_{S^3}D^4 ,
\]
and $H^*(E,E_0;\Z)\iso H^{*>0}(\CP^2;\Z)$, generated by $u$ in degree $2$ with
$u^2$ generating degree $4$. Chasing $u$ through
$H^2(E,E_0)\to H^2(E)\iso H^2(\CP^1)$ identifies
\[
e\big(\gamma_1\to\CP^1\big)=a,
\]
a generator of $H^2(\CP^1;\Z)$. Note that every complex bundle is canonically
oriented, by declaring $v_1,Jv_1,\dots,v_n,Jv_n$ positive for any complex basis
$v_1,\dots,v_n$; so the Euler class of a complex bundle is defined without extra
choices.
\end{example}

\begin{theorem}\label{thm:e-gamma-n}
$e(\gamma_n)=\tilde\sigma_n=c_n(\gamma_n)$ for the universal complex bundle,
and over $\Z_2$, $e(\gamma^n)=\tilde\sigma_n=w_n(\gamma^n)$ for the universal
real bundle. Consequently, for every bundle $\xi$,
\[
e(\xi)=c_n(\xi)\ \ (\xi \text{ complex}),
\qquad
e(\xi)\bmod2=w_n(\xi)\ \ (\xi \text{ real oriented}).
\]
\end{theorem}

\begin{proof}
\begin{strategy}
Compute the Euler class of a sum of $n$ line bundles over $(\CP^\infty)^n$,
where multiplicativity turns it into $a_1\cdots a_n=\sigma_n$; then use that
$h_n^*$ is injective to transport the identity to the universal bundle.
\end{strategy}

\step{1} \emph{Stabilizing.} The inclusion $\CP^1\hookrightarrow\CP^\infty$ is
covered by a bundle map $\gamma_1\to\gamma_1$ and induces an isomorphism on
$H^2$. By naturality and Example~\ref{ex:e-gamma-CP1},
$e(\gamma_1\to\CP^\infty)=a$, the generator of $H^2(\CP^\infty;\Z)=\Z[a]$.

\step{2} \emph{Sums of lines.} Over $(\CP^\infty)^n$, by
Theorem~\ref{thm:euler-properties}(5),
\[
e\big(\gamma_1\oplus\dots\oplus\gamma_1\big)
=e(\gamma_1)\smile\dots\smile e(\gamma_1)
=a_1a_2\cdots a_n=\sigma_n(a_1,\dots,a_n).
\]

\step{3} \emph{Transport.} Let $h_n\colon(\CP^\infty)^n\to G_n(\C^\infty)$
classify that sum, so $h_n^*\gamma_n\iso\gamma_1\oplus\dots\oplus\gamma_1$. By
naturality of the Euler class,
$h_n^*e(\gamma_n)=\sigma_n=h_n^*(\tilde\sigma_n)$, and $h_n^*$ is injective by
Theorem~\ref{thm:HGnC}; hence $e(\gamma_n)=\tilde\sigma_n=c_n(\gamma_n)$.

\step{4} \emph{Real case and general bundles.} The same three steps over
$\RP^\infty$ with $\Z_2$ coefficients give $e(\gamma^n)=\tilde\sigma_n$. For an
arbitrary bundle, pull back along a classifying map and use naturality of both
$e$ and the characteristic classes.
\end{proof}

\section{Proof of the Thom isomorphism theorem}

We now prove Theorem~\ref{thm:thom}. Throughout, $\pi\colon E\to B$ has rank $n$
and $E_0$ is the complement of the zero section. The model computation is the
following consequence of K\"unneth.

\begin{lemma}\label{lem:suspension}
Let $e^1\in H^1(\R,\R_0;\Lambda)$ be a generator, where $\R_0=\R\smallsetminus0$
and $\Lambda$ is any coefficient ring. Then for every space $B$ and every open
$B'\subset B$ the map
\[
H^{j}(B,B')\longrightarrow
H^{j+1}\big(B\times\R,\ B'\times\R\cup B\times\R_0\big),
\qquad
y\longmapsto y\times e^1,
\]
is an isomorphism; iterating $n$ times, with $e^n=e^1\times\dots\times e^1$,
\[
H^{j}(B,B')\ \xrightarrow{\ \cong\ }\
H^{j+n}\big(B\times\R^n,\ B'\times\R^n\cup B\times\R^n_0\big),
\qquad y\longmapsto y\times e^n .
\]
\end{lemma}

\begin{proof}
The pair $(\R,\R_0)$ has cohomology $\Lambda$ concentrated in degree $1$,
generated by $e^1$, and it is free; the relative K\"unneth theorem therefore
gives
\[
H^{j+1}\big(B\times\R,\ B'\times\R\cup B\times\R_0\big)
\iso\bigoplus_{p+q=j+1}H^p(B,B')\otimes H^q(\R,\R_0)
=H^{j}(B,B')\otimes\Lambda ,
\]
and the isomorphism is exactly $y\mapsto y\times e^1$. Iterate.
\end{proof}

\begin{proof}[Proof of Theorem~\ref{thm:thom}]
\begin{strategy}
Four stages of increasing generality. For a trivial bundle the statement is
Lemma~\ref{lem:suspension}. If the base splits into two open pieces over which
the theorem is known, Mayer--Vietoris and the five lemma give it over the union;
a finite trivializing cover then handles compact bases. Arbitrary bases follow
by writing cohomology as a limit over compact subsets, which works with field
coefficients. Finally an algebraic lemma about mapping cones upgrades field
coefficients to $\Z$.
\end{strategy}

\step{1} \emph{Case 1: $E$ trivial.} Here
$(E,E_0)\iso(B\times\R^n,B\times\R^n_0)$. Taking $B'=\emptyset$ in
Lemma~\ref{lem:suspension} gives
\[
H^j(B)\ \iso\ H^{j+n}(E,E_0),\qquad y\mapsto y\times e^n .
\]
In particular $H^{j+n}(E,E_0)=0$ for $j<0$, and $H^n(E,E_0)\iso H^0(B)$. Put
$u=1\times e^n$. Its restriction to a fibre pair is $e^n$, the preferred
generator, so $u$ is a Thom class; and it is unique because a class restricting
to $0$ on every fibre corresponds to $y\times e^n$ with $y$ vanishing at every
point of $B$, hence $y=0$. Finally, elements of $H^j(E)=H^j(B\times\R^n)$ are of
the form $y\times1$, and
\[
(y\times1)\smile u=(y\times1)\smile(1\times e^n)=y\times e^n,
\]
so $\smile u$ is the isomorphism just described.

\step{2} \emph{Case 2: Mayer--Vietoris.} Suppose $B=B'\cup B''$ with $B',B''$
open and the theorem known over $B'$, $B''$ and $B^\cap=B'\cap B''$. Write
$E',E'',E^\cap$ for the restrictions of $E$. Consider the Mayer--Vietoris
sequence of the pairs,
\[
H^{n-1}(E^\cap,E_0^\cap)\longrightarrow H^n(E,E_0)
\xrightarrow{\ (i'^*,\,i''^*)\ }H^n(E',E'_0)\oplus H^n(E'',E''_0)
\]
\[
\xrightarrow{\ i'^*-i''^*\ }H^n(E^\cap,E^\cap_0).
\]
By hypothesis $H^{n-1}(E^\cap,E^\cap_0)\iso H^{-1}(B^\cap)=0$. The classes
$u',u''$ both restrict to the unique Thom class $u^\cap$ of $E^\cap$, so the
pair $(u',u'')$ dies under $i'^*-i''^*$ and lifts to a unique
$u\in H^n(E,E_0)$; it restricts to the preferred generator on every fibre, since
every fibre lies over $B'$ or $B''$. Now $\smile u$ maps the Mayer--Vietoris
sequence of $(B',B'')$ to that of $(E',E'')$, commutatively; the maps over
$B',B'',B^\cap$ are isomorphisms by hypothesis, so the five lemma gives the
result over $B$.

\step{3} \emph{Case 3: $B$ compact.} Choose a finite trivializing cover
$U_1,\dots,U_r$. By Case 1 the theorem holds over each $U_i$ and over each
intersection, which is again trivializing. By Case 2 it holds over
$U_1\cup U_2$; inductively, if it holds over $U_1\cup\dots\cup U_j$ then apply
Case 2 to that set together with $U_{j+1}$, whose intersection with it is a
union of $j$ trivializing sets covered by the inductive hypothesis. After $r$
steps we reach $B$.

\step{4} \emph{Case 4: general base, field coefficients.} Every singular chain
has compact image, so
$H_j(B)=\varinjlim_C H_j(C)$ over the compact subsets $C\subset B$. With
\emph{field} coefficients $\Lambda$, cohomology is the dual of homology, and
duals turn direct limits into inverse limits:
\[
H^j(B;\Lambda)=\Hom\big(\varinjlim_C H_j(C),\Lambda\big)
=\varprojlim_C\Hom\big(H_j(C),\Lambda\big)=\varprojlim_C H^j(C;\Lambda),
\]
and similarly for the pairs $(E,E_0)$ restricted over $C$. By Case 3 there is a
compatible Thom class over each $C$, and compatible isomorphisms; passing to the
inverse limit gives a Thom class $u$ over $B$, its uniqueness, and the
isomorphism $\smile u$. This already proves the theorem over $\Z_2$ for
arbitrary bundles, which is all that Stiefel--Whitney theory requires.

\step{5} \emph{Integral coefficients.} Two further ingredients are needed, and
they occupy the rest of the lecture: the algebraic Lemma~\ref{lem:alg} below,
and the vanishing $H_{n-1}(E,E_0;\Z)=0$. Granting the latter, universal
coefficients,
\[
0\to\Ext\big(H_{n-1}(E,E_0),\Z\big)\to H^n(E,E_0;\Z)\to
\Hom\big(H_n(E,E_0),\Z\big)\to0,
\]
gives $H^n(E,E_0;\Z)\iso\Hom(H_n(E,E_0),\Z)$, and the compatible fibrewise
normalization determines a unique integral $u$. Choose a singular cocycle
$z\in Z^n(E,E_0)$ representing $u$. Cap product with $z$ is a chain map of
degree $-n$,
\[
z\,\cap\ \colon\ C_{n+i}(E,E_0)\longrightarrow C_i(E),
\qquad
\partial(z\cap\gamma)=(-1)^n z\cap\partial\gamma,
\]
whose dual on cochains is $c\mapsto c\smile z$. By Step 4 the induced map on
cohomology is an isomorphism over every field. By Lemma~\ref{lem:alg} it is then
an isomorphism over $\Z$, and over any coefficient ring.
\end{proof}

\begin{lemma}\label{lem:alg}
Let $f\colon K\to K'$ be a chain map of free chain complexes over $\Z$ which
induces isomorphisms $H^*(K';\Lambda)\to H^*(K;\Lambda)$ for every field
$\Lambda$. Then $f$ induces isomorphisms on homology and cohomology with any
coefficient ring.
\end{lemma}

\begin{proof}
\begin{strategy}
Encode ``$f$ is a homology isomorphism'' as the vanishing of the homology of the
algebraic mapping cone. That homology is finitely constrained by universal
coefficients: rational coefficients show it is torsion, and $\Z_p$ coefficients
show it has no $p$-torsion for any prime $p$.
\end{strategy}

\step{1} \emph{The mapping cone.} Let $d$ be the degree of $f$ and put
$K^f_i=K_{i-d-1}\oplus K'_i$ with
\[
\partial^f(k,k')=\big((-1)^{d+1}\partial k,\ f(k)+\partial k'\big).
\]
A direct computation gives $\partial^f\circ\partial^f=0$, using
$\partial f=(-1)^d f\partial$. The complex $K^f$ is free.

\step{2} \emph{The long exact sequence.} There is a short exact sequence of
complexes $0\to K'\to K^f\to K[\,d+1\,]\to0$, whose long exact sequence has
connecting homomorphism induced by $f$:
\[
\cdots\to H_i(K')\to H_i(K^f)\to H_{i-d-1}(K)
\xrightarrow{\ f_*\ }H_{i-1}(K')\to\cdots
\]
Hence $H_*(K^f)=0$ if and only if $f_*$ is an isomorphism on homology.

\step{3} \emph{Vanishing over every field.} By hypothesis $f$ induces
isomorphisms on cohomology with field coefficients, so the corresponding
sequence gives $H^*(K^f;\Lambda)=0$ for every field $\Lambda$. Universal
coefficients,
\[
0\to\Ext\big(H_{q-1}(K^f),\Lambda\big)\to H^q(K^f;\Lambda)\to
\Hom\big(H_q(K^f),\Lambda\big)\to0,
\]
then forces $\Hom(H_q(K^f),\Lambda)=0$ for every field.

\step{4} \emph{Conclusion.} Taking $\Lambda=\Q$ shows $H_q(K^f)$ has no free
part, so it is a torsion group. Taking $\Lambda=\Z_p$ and using also the $\Ext$
term shows $H_q(K^f)$ has no $p$-torsion, for every prime $p$. Hence
$H_q(K^f)=0$ for all $q$, and by Step 2 the map $f$ is a homology isomorphism
over $\Z$; the same then holds with any coefficients, again by universal
coefficients applied to the free complexes.
\end{proof}

\begin{corollary}\label{cor:cap-thom}
For an oriented rank $n$ bundle, the cap product with the Thom class,
\[
H_{n+i}(E,E_0;\Lambda)\longrightarrow H_i(E;\Lambda),
\qquad
\eta\longmapsto u\cap\eta,
\]
is an isomorphism for every coefficient ring $\Lambda$, compatibly with the
cohomological Thom isomorphism through
$\langle c,u\cap\gamma\rangle=\langle c\smile u,\gamma\rangle$.
\end{corollary}

\begin{remark}
It remains to justify $H_{n-1}(E,E_0;\Z)=0$, used in Step 5. For compact $B$
this follows from Corollary~\ref{cor:cap-thom} in the already established cases,
which gives $H_{n-1}(E,E_0)\iso H_{-1}(E)=0$; for general $B$ write
$H_{n-1}(E,E_0;\Z)=\varinjlim_C H_{n-1}(E|_C,E_0|_C;\Z)=0$, since homology
commutes with direct limits.
\end{remark}

With the Thom isomorphism established, the next lecture brings it into contact
with geometry: tubular neighbourhoods, dual classes of submanifolds, and the
diagonal.

\lecture

\begin{guiding}
The Thom class lives on the total space of a bundle. Geometry enters through one
observation: a submanifold has a neighbourhood that looks exactly like its
normal bundle. This transports the Thom class into the cohomology of the
ambient manifold, where it becomes the \emph{dual class} of the submanifold.
Applied to the diagonal in $M\times M$ it will compute the Euler class of $TM$.
\end{guiding}

Throughout, $M^n\subset A^{n+k}$ is a closed smooth $n$-manifold embedded in a
smooth Riemannian $(n+k)$-manifold, and
\[
\nu=\big\{v\in T_xA\ \big|\ x\in M,\ v\perp T_xM\big\}
\]
is its normal bundle, so that $TA|_M\iso TM\oplus\nu$.

\section{The tubular neighbourhood theorem}

\begin{theorem}\label{thm:tnt}
There is an open neighbourhood of $M$ in $A$ diffeomorphic to the total space
$E(\nu)$, by a diffeomorphism restricting to the identity on the zero section.
\end{theorem}

\begin{proof}
\begin{strategy}
Write down a map from the normal bundle to $A$ which is the identity on $M$ and
whose differential along $M$ is invertible, then quote the global inverse
function theorem. The map is the exponential map of the Riemannian metric.
\end{strategy}

\step{1} \emph{The inverse function theorem, global form (quoted).} If
$f\colon N\to A$ is smooth, injective on a compact submanifold $M\subset N$, and
$df_x$ is an isomorphism for every $x\in M$, then $f$ maps some neighbourhood of
$M$ in $N$ diffeomorphically onto a neighbourhood of $f(M)$ in $A$.

\step{2} \emph{The exponential map.} For $\eps>0$ let
$E(\eps)=\{v\in E(\nu)\mid |v|<\eps\}$, an open neighbourhood of the zero
section, and define
\[
\mathrm{Exp}\colon E(\eps)\longrightarrow A,
\qquad
(x,v)\longmapsto\gamma_v(1),
\]
where $\gamma_v$ is the unique geodesic with $\gamma_v(0)=x$ and
$\dot\gamma_v(0)=v$. For $\eps$ small this is defined and smooth, by the theory
of ordinary differential equations.

\step{3} \emph{Behaviour along $M$.} On the zero section $\mathrm{Exp}$ is the
identity. Its differential at a point of the zero section, under the canonical
splitting $T\big(E(\nu)\big)\big|_M\iso TM\oplus\nu\iso TA|_M$, is the identity:
the $TM$ directions are unmoved and the $\nu$ directions differentiate the
geodesics, giving back $v$.

\step{4} \emph{Conclusion.} Step 1 applies with $N=E(\eps)$ and $f=\mathrm{Exp}$,
so a neighbourhood of the zero section in $E(\eps)$ maps diffeomorphically onto
a neighbourhood $U$ of $M$ in $A$. Finally $E(\nu)\iso E(\eps)$ by the fibrewise
rescaling $v\mapsto \eps v/\sqrt{1+|v|^2}$.
\end{proof}

\begin{center}
\begin{tikzpicture}[scale=0.95]
\draw[very thick] (0,0) .. controls (1.3,0.8) and (2.7,-0.5) .. (4,0.25);
\node at (4.35,0.25) {$M$};
\draw[dashed] (0,0.42) .. controls (1.3,1.22) and (2.7,-0.08) .. (4,0.67);
\draw[dashed] (0,-0.42) .. controls (1.3,0.38) and (2.7,-0.92) .. (4,-0.17);
\node at (6.1,0.25) {\small tube $\iso E(\nu)$, inside $A$};
\end{tikzpicture}
\end{center}

\begin{corollary}\label{cor:excision-tube}
If $M$ is closed in $A$ then, for any coefficient ring,
\[
H^*\big(E,E_0\big)\ \iso\ H^*\big(A,\,A\smallsetminus M\big),
\qquad E=E(\nu).
\]
\end{corollary}

\begin{proof}
Let $U\subset A$ be a tubular neighbourhood, so $(U,U\smallsetminus M)\iso
(E,E_0)$ by Theorem~\ref{thm:tnt}. Excising the closed set
$A\smallsetminus U$ from the pair $(A,A\smallsetminus M)$ gives
$H^*(A,A\smallsetminus M)\iso H^*(U,U\smallsetminus M)$.
\end{proof}

\section{The dual class of a submanifold}

Assume from now on either $\Z_2$ coefficients, or that $TA$ and $TM$ are
oriented and coefficients are $\Z$; in the latter case $\nu$ inherits an
orientation from $TA|_M\iso TM\oplus\nu$. By Theorem~\ref{thm:thom} there is a
unique Thom class $u\in H^k(E,E_0)$, and Corollary~\ref{cor:excision-tube}
carries it to a class
\[
u'\in H^k\big(A,\,A\smallsetminus M\big),
\]
whose image in $H^k(A)$ we denote $u''$. Both are called the \emph{dual class}
of $M$ in $A$.

\begin{remark}
The class $u'$ does not depend on the Riemannian metric: two metrics $g_1,g_2$
are joined by the path $tg_1+(1-t)g_2$ of metrics, and the corresponding
exponential identifications are homotopic.
\end{remark}

\begin{theorem}\label{thm:dual-restricts-euler}
The inclusions $M\hookrightarrow A\hookrightarrow(A,A\smallsetminus M)$ induce
\[
H^k\big(A,A\smallsetminus M\big)\longrightarrow H^k(A)\longrightarrow H^k(M),
\]
under which $u'\mapsto u''$ and $u''$ restricts to
\[
\begin{cases}
e(\nu)&\text{in the oriented case},\\[2pt]
w_k(\nu)&\text{over }\Z_2 .
\end{cases}
\]
\end{theorem}

\begin{proof}
Restrict the whole picture to the tube. The composite
$H^k(E,E_0)\to H^k(E)\iso H^k(M)$ sends $u$ to $e(\nu)$, by the very definition
of the Euler class; over $\Z_2$ it sends $u$ to $w_k(\nu)$ by
Theorem~\ref{thm:euler-properties}(4). The square
\[
\begin{tikzcd}[row sep=large, column sep=large]
H^k(A,A{\smallsetminus}M)\arrow[r]\arrow[d,"\iso"'] & H^k(A)\arrow[r]\arrow[d]
& H^k(M)\arrow[d,equal]\\
H^k(E,E_0)\arrow[r] & H^k(E)\arrow[r,"\iso"] & H^k(M)
\end{tikzcd}
\]
commutes, since all maps are induced by inclusions and the tubular
identification.
\end{proof}

\begin{corollary}\label{cor:embedding-obstruction}
If $M^n$ is closed and smoothly embedded in $\R^{n+k}$, then $w_k(\nu)=0$, and
in the oriented case $e(\nu)=0$.
\end{corollary}

\begin{proof}
Here $u''$ lies in $H^k(\R^{n+k})=0$ for $k>0$, so its restriction to $M$
vanishes; apply Theorem~\ref{thm:dual-restricts-euler}.
\end{proof}

\begin{example}
Combined with Whitney duality this refines the immersion obstructions of
Lecture 2 into \emph{embedding} obstructions. Since
$\bar w(T\RP^{2^r})=1+a+\dots+a^{2^r-1}$ has non-zero term in degree $2^r-1$,
an embedding $\RP^{2^r}\subset\R^{2^r+k}$ would give
$w_j(\nu)=\bar w_j(T\RP^{2^r})=0$ for $j>k$ and, by
Corollary~\ref{cor:embedding-obstruction}, also $w_k(\nu)=0$; hence
$k\ge 2^r$. So $\RP^{2^r}$ does not embed in $\R^{2\cdot2^r-1}$, which is sharp
against Whitney's embedding theorem $M^n\hookrightarrow\R^{2n}$.
\end{example}

\section{The diagonal}

Apply the machine to the diagonal embedding
$\Delta\colon M\to M\times M$, $x\mapsto(x,x)$, of a closed Riemannian manifold;
the product carries the product metric.

\begin{lemma}\label{lem:normal-diagonal}
The normal bundle $N(\Delta)$ of $\Delta(M)\subset M\times M$ is canonically
isomorphic to $TM$.
\end{lemma}

\begin{proof}
At $(x,x)$ the tangent space of the diagonal is
$T\Delta=\{(v,v)\}\subset T_xM\times T_xM$, whose orthogonal complement is
$N(\Delta)=\{(-v,v)\}$. The map
\[
TM\longrightarrow N(\Delta),
\qquad
(x,v)\longmapsto\big((x,x),\,(-v,v)\big),
\]
is continuous, fibrewise a linear isomorphism, hence an isomorphism of bundles
by Lemma~\ref{lem:fibrewise-iso}.
\end{proof}

\begin{corollary}\label{cor:diagonal-euler}
Let $u''\in H^n(M\times M)$ be the dual class of the diagonal. Then
\[
\Delta^*u''=e(TM)
\qquad\big(\text{resp. } w_n(TM)\ \text{over }\Z_2\big).
\]
\end{corollary}

\begin{proof}
Theorem~\ref{thm:dual-restricts-euler} identifies the restriction of $u''$ to
$\Delta(M)$ with $e(N(\Delta))$, and
Lemma~\ref{lem:normal-diagonal} identifies $N(\Delta)$ with $TM$. Restricting
to $\Delta(M)$ is the same as applying $\Delta^*$.
\end{proof}

The dual class of the diagonal interacts with the product structure very
rigidly.

\begin{lemma}\label{lem:ax1}
For every $a\in H^*(M)$,
\[
(a\times1)\smile u''=(1\times a)\smile u''\ \in H^*(M\times M).
\]
\end{lemma}

\begin{proof}
\begin{strategy}
Both products are computed from classes supported near the diagonal, and near
the diagonal the two projections $M\times M\to M$ agree. Make this precise by
factoring the cup product through the relative group and restricting to a
tubular neighbourhood, which retracts onto the diagonal.
\end{strategy}

\step{1} \emph{Where the product is computed.} Let $U$ be a tubular
neighbourhood of $\Delta(M)$ in $M\times M$. Cup product with the relative class
$u'$ gives a commutative square
\[
\begin{tikzcd}[row sep=large, column sep=large]
H^i(M\times M)\arrow[r]\arrow[d,"\smile u'"'] & H^i(U)\arrow[d,"\smile u'|_U"]\\
H^{i+n}\big(M\times M,\ M\times M\smallsetminus\Delta\big)
\arrow[r,"\iso"',"\text{excision}"] & H^{i+n}(U,\,U\smallsetminus\Delta)
\end{tikzcd}
\]
so $x\smile u'$ depends only on the restriction $x|_U$.

\step{2} \emph{The two restrictions agree.} The diagonal is a deformation
retract of $U$, and $p_1\circ\Delta=p_2\circ\Delta=\mathrm{id}$; hence the two
maps $p_1,p_2\colon U\to M$ are homotopic, and for $a\in H^*(M)$,
\[
(a\times1)\big|_U=p_1^*a\big|_U=p_2^*a\big|_U=(1\times a)\big|_U .
\]

\step{3} \emph{Conclusion.} By Steps 1 and 2,
$(a\times1)\smile u'=(1\times a)\smile u'$ in
$H^*(M\times M,\,M\times M\smallsetminus\Delta)$. Mapping to the absolute group
$H^*(M\times M)$ turns $u'$ into $u''$ and preserves the identity.
\end{proof}

\section{The slant product}

We work with field coefficients, so that $H^*(X\times Y)=H^*(X)\otimes H^*(Y)$.

\begin{definition}
The \emph{slant product}
\[
/\ \colon\ H^{p+q}(X\times Y)\otimes H_q(Y)\longrightarrow H^p(X)
\]
is defined on decomposable classes by
$(\alpha\times\beta)/\mu=\langle\beta,\mu\rangle\,\alpha$ and extended linearly.
\end{definition}

\begin{lemma}\label{lem:slant-identity}
For $a\in H^*(X)$, $P\in H^*(X\times Y)$ and $\mu\in H_*(Y)$,
\[
\big((a\times1)\smile P\big)\big/\mu\ =\ a\smile\big(P/\mu\big).
\]
\end{lemma}

\begin{proof}
Write $P=\sum_{i,j}a_{ij}\,p_i\times q_j$ with $p_i\in H^*(X)$,
$q_j\in H^*(Y)$. Then
\[
(a\times1)\smile P=\sum_{i,j}a_{ij}\,(a\smile p_i)\times q_j ,
\]
so slanting with $\mu$ gives
$\sum_{i,j}a_{ij}\langle q_j,\mu\rangle\,(a\smile p_i)
=a\smile\big(\sum_{i,j}a_{ij}\langle q_j,\mu\rangle\,p_i\big)=a\smile(P/\mu)$.
\end{proof}

\begin{lemma}\label{lem:u-over-mu}
Let $\mu\in H_n(M)$ be the fundamental class. Then
\[
u''/\mu=1\ \in H^0(M).
\]
\end{lemma}

The proof is a local computation which opens the next lecture. Together with
Lemma~\ref{lem:ax1} and Corollary~\ref{cor:diagonal-euler} it will determine
$u''$ completely and produce the identity
$\langle e(TM),\mu\rangle=\chi(M)$.

\lecture

\begin{guiding}
We finish the diagonal computation: the dual class of the diagonal is a sum over
dual bases, and evaluating it turns $e(TM)$ into the Euler characteristic. Then
we introduce the last family of characteristic classes, the Pontryagin classes
of a real bundle.
\end{guiding}

Throughout $M$ is a closed $n$-manifold with fundamental class $\mu$;
coefficients are a field, always admissible over $\Z_2$ and, when $M$ is
oriented, over any field.

\section{The local computation}

\begin{proof}[Proof of Lemma~\ref{lem:u-over-mu}]
\begin{strategy}
The class $u''/\mu$ lies in $H^0(M)$, so it is determined by its value at each
point $x$. Naturality of the slant product turns that value into the evaluation
of the dual class against the local fundamental class at $x$, and the
fibre-restriction property of the Thom class makes it $1$.
\end{strategy}

\step{1} \emph{Reduction to a point.} Fix $x\in M$ and let
$i_x\colon M\to M\times M$, $y\mapsto(x,y)$. Naturality of the slant product in
the first variable gives that the value of $u''/\mu$ at $x$ is
\[
\big(i_x^*u''\big)/\mu=\big\langle i_x^*u'',\ \mu\big\rangle ,
\]
where we use that $i_x^*u''\in H^n(M)$ and that slanting an $n$-class with the
$n$-cycle $\mu$ is the Kronecker pairing.

\step{2} \emph{Passing to the relative class.} Let
$j_x\colon(M,M\smallsetminus x)\to(M\times M,\ M\times M\smallsetminus\Delta)$ be
$y\mapsto(x,y)$; it is well defined, since $(x,y)\in\Delta$ forces $y=x$. Then
$i_x^*u''=i_M^*j_x^*u'$, where $i_M\colon M\to(M,M\smallsetminus x)$, so with
$\mu_x\in H_n(M,M\smallsetminus x)$ the image of $\mu$,
\[
\big\langle i_x^*u'',\mu\big\rangle
=\big\langle j_x^*u',\ \mu_x\big\rangle .
\]

\step{3} \emph{Identifying $j_x^*u'$.} Under the exponential chart of
Theorem~\ref{thm:tnt}, the map $j_x$ corresponds to the inclusion of the fibre
of $N(\Delta)$ over $(x,x)$:
\[
\big(N_x,\ N_x\smallsetminus0\big)
\xrightarrow{\ \mathrm{Exp}\ }
\big(M,\ M\smallsetminus x\big)
\xrightarrow{\ j_x\ }
\big(M\times M,\ M\times M\smallsetminus\Delta\big),
\]
the first map being a homeomorphism onto a neighbourhood, and the composite
being homotopic to $v\mapsto(x,\exp_x v)$. Since $u'$ comes from the Thom class,
its restriction to a fibre pair is the preferred generator; hence $j_x^*u'$ is
the preferred generator of $H^n(T_xM,T_xM\smallsetminus0)$.

\step{4} \emph{Conclusion.} The correspondence between orientations of $T_xM$
and local orientations of $M$ is exactly the one induced by $\exp$, so $\mu_x$
corresponds to the preferred generator in homology and
$\langle j_x^*u',\mu_x\rangle=1$. Therefore $u''/\mu$ takes the value $1$ at
every $x$, that is, $u''/\mu=1$.
\end{proof}

\section{The duality theorem and the Euler characteristic}

\begin{theorem}[Poincar\'e duality, quoted]\label{thm:duality}
For $M$ closed, and with field coefficients (any field if $M$ is oriented,
$\Z_2$ otherwise), the cup product pairing
$H^i(M)\otimes H^{n-i}(M)\to\Lambda$, $(x,y)\mapsto\langle x\smile y,\mu\rangle$,
is non-degenerate. Equivalently, every homogeneous basis $b_1,\dots,b_r$ of
$H^*(M)$ admits a \emph{dual basis} $b_1^\#,\dots,b_r^\#$ with
\[
\big\langle b_i\smile b_j^{\#},\ \mu\big\rangle=\delta_{ij}.
\]
\end{theorem}

\begin{theorem}\label{claim:diag-class-formula}
In terms of such bases the dual class of the diagonal is
\[
u''=\sum_{i=1}^{r}(-1)^{\dim b_i}\ b_i\times b_i^{\#}.
\]
\end{theorem}

\begin{proof}
\begin{strategy}
Expand $u''$ in the K\"unneth basis with unknown coefficients, then exploit the
two facts proved so far: $u''$ slants to $1$ against $\mu$, and multiplying by
$a\times1$ or $1\times a$ gives the same answer. Evaluating on basis elements
identifies the unknowns as the dual basis.
\end{strategy}

\step{1} \emph{Expansion.} By K\"unneth,
$H^n(M\times M)=\bigoplus_i H^{\dim b_i}(M)\otimes H^{n-\dim b_i}(M)$, so we may
write
\[
u''=\sum_{i=1}^{r} b_i\times c_i,
\qquad \dim b_i+\dim c_i=n .
\]

\step{2} \emph{A formula for every class.} Let $a\in H^*(M)$. Using
Lemma~\ref{lem:slant-identity}, then Lemma~\ref{lem:ax1}, then
Lemma~\ref{lem:u-over-mu},
\[
a=a\smile\big(u''/\mu\big)
=\big((a\times1)\smile u''\big)\big/\mu
=\big((1\times a)\smile u''\big)\big/\mu .
\]

\step{3} \emph{Expanding the right-hand side.} Substituting the expansion of
$u''$ and moving $a$ past $b_i$, which costs the Koszul sign
$(-1)^{\dim a\,\dim b_i}$,
\[
\big((1\times a)\smile u''\big)\big/\mu
=\sum_{i}(-1)^{\dim a\,\dim b_i}\,\big(b_i\times(a\smile c_i)\big)\big/\mu
\]
\[
=\sum_{i}(-1)^{\dim a\,\dim b_i}\,
\big\langle a\smile c_i,\ \mu\big\rangle\ b_i .
\]

\step{4} \emph{Identification.} Take $a=b_j$. Comparing with Step 2, which says
the sum equals $b_j$, and using that the $b_i$ are a basis,
\[
(-1)^{\dim b_j\dim b_i}\big\langle b_j\smile c_i,\mu\big\rangle=\delta_{ij}.
\]
For $i=j$ the sign is $(-1)^{(\dim b_i)^2}=(-1)^{\dim b_i}$, so
$\langle b_i\smile c_i,\mu\rangle=(-1)^{\dim b_i}$, while for $i\neq j$ the
pairing vanishes. Comparing with Theorem~\ref{thm:duality} gives
$c_i=(-1)^{\dim b_i}b_i^{\#}$.
\end{proof}

\begin{theorem}\label{thm:euler-char}
For $M$ closed and oriented,
\[
\big\langle e(TM),\ \mu\big\rangle=\chi(M)=\sum_k(-1)^k\dim H^k(M);
\]
without orientation, $\langle w_n(TM),\mu\rangle\equiv\chi(M)\bmod 2$.
\end{theorem}

\begin{proof}
By Corollary~\ref{cor:diagonal-euler} and
Theorem~\ref{claim:diag-class-formula},
\[
e(TM)=\Delta^*u''=\sum_i(-1)^{\dim b_i}\,\Delta^*\big(b_i\times b_i^\#\big)
=\sum_i(-1)^{\dim b_i}\,b_i\smile b_i^{\#} .
\]
Evaluating on $\mu$ and using $\langle b_i\smile b_i^\#,\mu\rangle=1$,
\[
\big\langle e(TM),\mu\big\rangle=\sum_i(-1)^{\dim b_i}
=\sum_k(-1)^k\dim H^k(M)=\chi(M),
\]
since the $b_i$ of a given dimension $k$ form a basis of $H^k(M)$.
\end{proof}

\begin{example}
$\chi(S^2)=2$, so $e(TS^2)\neq0$: a third proof that $TS^2$ is non-trivial, and
the statement underlying the hairy ball theorem. For odd-dimensional $M$ we have
$\chi(M)=0$, consistent with $2e(\xi)=0$ in odd rank.
\end{example}

\section{Chern classes through the Gysin sequence}

We describe the inductive construction of Chern classes used in
Milnor--Stasheff, an alternative to the Grassmannian route of Lecture 6.

\begin{theorem}[Gysin sequence]\label{thm:gysin}
For an oriented $\R^n$-bundle $\xi$ with total space $E\to B$ there is a long
exact sequence with integer coefficients
\[
\cdots\to H^{i-n}(B)\ \xrightarrow{\ \smile e\ }\ H^{i}(B)
\ \xrightarrow{\ \pi_0^*\ }\ H^{i}(E_0)\ \to\ H^{i-n+1}(B)\to\cdots
\]
\end{theorem}

\begin{proof}
Start from the long exact sequence of the pair $(E,E_0)$ and substitute
$H^{j}(E,E_0)\iso H^{j-n}(B)$, by the Thom isomorphism, and
$H^j(E)\iso H^j(B)$, by homotopy invariance. The map
$H^{j-n}(B)\to H^{j}(B)$ so obtained is computed by
\[
i^*\phi(y)=i^*\big(\pi^*y\smile u\big)=\pi^*y\smile i^*u=\pi^*\big(y\smile e\big),
\]
so under the identification it becomes $\smile e$.
\end{proof}

\begin{construction}[Chern classes, inductively]
Let $E\to B$ be a complex bundle of rank $n$, with its canonical orientation.
Set $c_n(E):=e(E)$. Over the sphere bundle $\pi_0\colon S(E)=E_0\to B$ there is
a tautological complex line sub-bundle $L\subset\pi_0^*E$, spanned at a point
$(x,v)$ by $v$; let $W=L^{\perp}$, of rank $n-1$. In degrees $i<2n-1$ the Gysin
sequence shows $\pi_0^*\colon H^i(B)\to H^i(E_0)$ is an isomorphism, so we may
define
\[
c_k(E):=(\pi_0^*)^{-1}c_k(W),\qquad k<n,
\]
inductively on the rank. The two definitions of $c_k$ agree, as one checks on
sums of line bundles using the splitting principle.
\end{construction}

\begin{proposition}
The total Chern class $c(E)=1+c_1+\dots+c_n$ is invertible in $H^\Pi(B;\Z)$.
\end{proposition}

\begin{proof}
The proof of Proposition~\ref{prop:invertible} works verbatim over $\Z$: the
inverse is constructed degree by degree, each coefficient being determined by
the previous ones.
\end{proof}

\section{Conjugate bundles and Pontryagin classes}

\begin{definition}
For a complex bundle $\xi$, the \emph{conjugate bundle} $\bar\xi$ is the same
real bundle equipped with the opposite complex structure $-J$.
\end{definition}

\begin{proposition}\label{prop:conjugate}
$c_k(\bar\xi)=(-1)^k c_k(\xi)$.
\end{proposition}

\begin{proof}
\begin{strategy}
Check it for line bundles, where conjugation reverses the canonical orientation
and hence the sign of the Euler class, then propagate by the splitting principle
and injectivity of $h_n^*$.
\end{strategy}

\step{1} \emph{Line bundles.} For a complex line bundle $L$, a complex basis $v$
of a fibre gives the oriented real basis $(v,Jv)$; passing to $-J$ replaces it
by $(v,-Jv)$, which has the opposite orientation. Hence
$e(\bar L)=-e(L)$ and, since $c_1=e$ for line bundles by
Theorem~\ref{thm:e-gamma-n}, $c_1(\bar L)=-c_1(L)$.

\step{2} \emph{Universal case.} Over $(\CP^\infty)^n$, conjugating
$\gamma_1\oplus\dots\oplus\gamma_1$ conjugates each summand, so by Step 1 and
the Whitney formula
\[
h_n^*c_k(\bar\gamma_n)=\sigma_k(-a_1,\dots,-a_n)=(-1)^k\sigma_k(a_1,\dots,a_n)
=(-1)^k h_n^*c_k(\gamma_n).
\]
Injectivity of $h_n^*$ gives $c_k(\bar\gamma_n)=(-1)^kc_k(\gamma_n)$.

\step{3} \emph{General case.} If $f$ classifies $\xi$ then it also classifies
$\bar\xi$ as a conjugate, in the sense that $f^*\bar\gamma_n\iso\bar\xi$; apply
naturality.
\end{proof}

\begin{definition}
For a real bundle $\xi$, the \emph{complexification} is
$\xi\otimes\C$, whose underlying real bundle is $\xi\oplus\xi$ with complex
structure $J(x,y)=(-y,x)$.
\end{definition}

\begin{lemma}\label{lem:complexification-selfconj}
$\xi\otimes\C\iso\overline{\xi\otimes\C}$. Consequently the odd Chern classes of
a complexification are $2$-torsion:
$2c_{2i+1}(\xi\otimes\C)=0$.
\end{lemma}

\begin{proof}
The map $(x,y)\mapsto(x,-y)$ is a real bundle isomorphism
$\xi\oplus\xi\to\xi\oplus\xi$, and it is conjugate linear for $J$:
\[
J(x,-y)=(y,x),
\qquad\text{while}\qquad
-\overline{J}(x,y)\ \text{corresponds to}\ (y,x),
\]
so it is an isomorphism of complex bundles from $\xi\otimes\C$ onto its
conjugate. By Proposition~\ref{prop:conjugate},
$c_{2i+1}(\xi\otimes\C)=c_{2i+1}(\overline{\xi\otimes\C})
=-c_{2i+1}(\xi\otimes\C)$.
\end{proof}

\begin{definition}[Pontryagin classes]
For a real bundle $E\to B$ set
\[
p_i(E):=(-1)^i\,c_{2i}\big(E\otimes\C\big)\ \in H^{4i}(B;\Z),
\qquad p_0=1,
\]
and $p(E)=1+p_1(E)+p_2(E)+\cdots$. Naturality
$f^*p_i(E)=p_i(f^*E)$ is inherited from the Chern classes.
\end{definition}

\begin{theorem}\label{thm:pontryagin-whitney}
$2\big(p(E_1\oplus E_2)-p(E_1)p(E_2)\big)=0$; that is, the total Pontryagin
class is multiplicative modulo $2$-torsion.
\end{theorem}

\begin{proof}
Complexification is additive:
$(E_1\oplus E_2)\otimes\C\iso(E_1\otimes\C)\oplus(E_2\otimes\C)$. The Whitney
formula for Chern classes gives
$c\big((E_1\oplus E_2)\otimes\C\big)=c(E_1\otimes\C)\,c(E_2\otimes\C)$. Passing
from $c$ to $p$ discards the odd Chern classes, and every discarded term is
$2$-torsion by Lemma~\ref{lem:complexification-selfconj}; multiplying the
difference by $2$ therefore kills it.
\end{proof}

\begin{corollary}\label{cor:pontryagin-of-complex}
If $E$ is a complex bundle then $E\otimes\C\iso E\oplus\bar E$, hence
\[
1-p_1+p_2-\cdots+(-1)^np_n
=\big(1-c_1+c_2-\cdots\big)\big(1+c_1+c_2+\cdots\big),
\]
so for instance $p_1=c_1^2-2c_2$ and $p_2=c_2^2-2c_1c_3+2c_4$.
\end{corollary}

\begin{proof}
For a complex $E$ with structure $J$, the map
\[
E\otimes\C=E\oplus E\longrightarrow E\oplus\bar E,
\qquad
(x,y)\longmapsto\Big(\tfrac12(x+Jy),\ \tfrac12(x-Jy)\Big),
\]
is a complex isomorphism, as one checks by applying $J$ to both sides. Hence
$c(E\otimes\C)=c(E)\,c(\bar E)$, and
Proposition~\ref{prop:conjugate} turns $c(\bar E)$ into the alternating series.
Comparing coefficients of even degree gives the stated identity.
\end{proof}

\begin{example}
$T(S^n)\oplus\nu(S^n)$ is trivial with $\nu$ trivial, so by
Theorem~\ref{thm:pontryagin-whitney} and the torsion-freeness of $H^*(S^n;\Z)$,
\[
p(TS^n)=1 .
\]
Pontryagin classes do not distinguish spheres from parallelizable manifolds.
\end{example}

In the next lecture we compute $p(T\CP^n)$, which is far from trivial, extract
numerical invariants and put them to work on cobordism questions.

\lecture

\begin{guiding}
Time to harvest. First the model computation of the whole theory,
$c(T\CP^n)=(1+a)^{n+1}$ and $p(T\CP^n)=(1+u^2)^{n+1}$. Then we turn classes into
numbers and prove genuinely geometric statements: $\CP^{2n}$ admits no
orientation reversing diffeomorphism, and it is not an oriented boundary.
Finally we build the classes $S_I$, designed to detect products.
\end{guiding}

\section{The tangent bundle of complex projective space}

\begin{theorem}\label{thm:cTCPn}
Let $a=-c_1(\gamma_1\to\CP^n)$, a generator of $H^2(\CP^n;\Z)$. Then
\[
c\big(T\CP^n\big)=(1+a)^{n+1}.
\]
\end{theorem}

\begin{proof}
\begin{strategy}
Repeat the argument of Lemma~\ref{lem:TRPn} over $\C$: a tangent direction at a
line $L\subset\C^{n+1}$ is a complex linear map $L\to L^\perp$; adding the
trivial summand $\Hom_\C(\gamma_1,\gamma_1)$ produces a sum of $n+1$ copies of
the dual of $\gamma_1$, which is the conjugate bundle.
\end{strategy}

\step{1} \emph{The tangent bundle.} Exactly as in the real case,
\[
T\CP^n\ \iso\ \Hom_\C\big(\gamma_1,\gamma_1^{\perp}\big),
\]
where $\gamma_1^\perp$ is the complex orthogonal complement of $\gamma_1$ inside
the trivial bundle $\CP^n\times\C^{n+1}$, taken with respect to a Hermitian
metric.

\step{2} \emph{Adding a trivial line.} The bundle
$\Hom_\C(\gamma_1,\gamma_1)$ has the nowhere zero section $\mathrm{id}$, hence
is the trivial complex line bundle $\eps^1_\C$. Therefore
\[
T\CP^n\oplus\eps^1_\C
\iso\Hom_\C\big(\gamma_1,\ \gamma_1^\perp\oplus\gamma_1\big)
=\Hom_\C\big(\gamma_1,\ \C^{n+1}\big)
\iso\big(\gamma_1^{\,*}\big)^{\oplus(n+1)} .
\]

\step{3} \emph{The dual is the conjugate.} A Hermitian metric $h$ identifies
$\gamma_1^{\,*}=\Hom_\C(\gamma_1,\C)$ with $\overline{\gamma_1}$: the map
$v\mapsto h(\cdot,v)$ is conjugate linear in $v$, hence complex linear as a map
into the conjugate bundle. By Proposition~\ref{prop:conjugate},
$c_1(\overline{\gamma_1})=-c_1(\gamma_1)=a$.

\step{4} \emph{Conclusion.} Using Step 2, the Whitney formula and Step 3,
\[
c\big(T\CP^n\big)=c\big(T\CP^n\oplus\eps^1_\C\big)
=c\big(\overline{\gamma_1}\big)^{n+1}=(1+a)^{n+1}.
\]
\end{proof}

\begin{theorem}\label{thm:pTCPn}
With $u=a$ as above,
\[
p\big(T\CP^n\big)=(1+u^2)^{n+1},
\qquad\text{i.e.}\qquad
p_k\big(T\CP^n\big)=\binom{n+1}{k}u^{2k}.
\]
\end{theorem}

\begin{proof}
By Corollary~\ref{cor:pontryagin-of-complex} applied to $E=T\CP^n$,
\[
1-p_1+p_2-\cdots
=c\big(T\CP^n\big)\ c\big(\overline{T\CP^n}\big)
=(1+u)^{n+1}(1-u)^{n+1}=(1-u^2)^{n+1},
\]
where the middle equality uses Theorem~\ref{thm:cTCPn} and
Proposition~\ref{prop:conjugate}. Comparing the coefficients of $u^{2k}$ on both
sides, the sign $(-1)^k$ occurs on each, so
$p_k=\binom{n+1}{k}u^{2k}$.
\end{proof}

\begin{example}
$p(T\CP^5)=(1+u^2)^6=1+6u^2+15u^4$, the higher terms vanishing in
$H^*(\CP^5;\Z)$. In contrast with spheres, Pontryagin classes see projective
spaces very clearly.
\end{example}

\begin{remark}\label{rem:two-orientations}
For an oriented real bundle $V$ of rank $n$ there are two natural orientations
on $V\otimes\C\iso V\oplus V$: the complex one, given by
$v_1,Jv_1,\dots,v_n,Jv_n$, and the direct sum one, given by
$v_1,\dots,v_n,Jv_1,\dots,Jv_n$. Reordering one into the other requires
$n(n-1)/2$ transpositions, so the two differ by the sign $(-1)^{n(n-1)/2}$.
\end{remark}

\begin{lemma}\label{lem:pn-euler-square}
If $V$ is an oriented real bundle of rank $2n$ then $p_n(V)=e(V)^2$.
\end{lemma}

\begin{proof}
By definition $p_n(V)=(-1)^nc_{2n}(V\otimes\C)$, and
$c_{2n}=e$ for a complex bundle of rank $2n$ by
Theorem~\ref{thm:e-gamma-n}, computed with the complex orientation. By
Remark~\ref{rem:two-orientations} with $n$ replaced by $2n$, the complex and sum
orientations differ by $(-1)^{2n(2n-1)/2}=(-1)^n$, so
\[
p_n(V)=(-1)^n\cdot(-1)^n\,e\big(V\oplus V\big)=e(V)^2,
\]
using $e(V\oplus V)=e(V)\smile e(V)$ from
Theorem~\ref{thm:euler-properties}(5).
\end{proof}

\section{Chern numbers and Pontryagin numbers}

\begin{definition}
Let $X$ be a closed almost complex manifold of complex dimension $n$, so that
$TX$ carries a complex structure, and let $I=(i_1,\dots,i_r)$ be a partition of
$n$. The \emph{Chern number} is
\[
C_I(X)=\big\langle c_{i_1}(TX)\smile\cdots\smile c_{i_r}(TX),\ [X]\big\rangle
\in\Z,
\]
where $[X]$ is the fundamental class of the complex orientation.
\end{definition}

\begin{definition}
Let $M^{4n}$ be closed and oriented and $I=(i_1,\dots,i_r)$ a partition of $n$.
The \emph{Pontryagin number} is
\[
P_I(M)=\big\langle p_{i_1}(TM)\cdots p_{i_r}(TM),\ [M]\big\rangle\in\Z .
\]
\end{definition}

\begin{example}
By Theorem~\ref{thm:cTCPn}, $C_I(\CP^n)=\prod_j\binom{n+1}{i_j}>0$ for every
partition $I$ of $n$; and by Theorem~\ref{thm:pTCPn},
$P_I(\CP^{2n})=\prod_j\binom{2n+1}{i_j}\neq0$ for every partition $I$ of $n$.
\end{example}

\begin{remark}
For an oriented bundle $V$ and its reverse $\bar V$ one has $p_i(V)=p_i(\bar V)$,
since the complexification does not see the orientation. Reversing the
orientation of the \emph{manifold}, however, reverses $[M]$ and therefore
changes the sign of every Pontryagin number.
\end{remark}

\begin{lemma}\label{lem:no-orientation-reversal}
If some Pontryagin number of $M^{4n}$ is non-zero, then $M$ admits no
orientation reversing diffeomorphism.
\end{lemma}

\begin{proof}
Let $f\colon M\to M$ be a diffeomorphism. Its differential is a bundle
isomorphism $TM\iso f^*TM$, so $f^*p_i(TM)=p_i(f^*TM)=p_i(TM)$. If moreover $f$
reverses orientation then $f_*[M]=-[M]$, whence for every $I$
\[
P_I(M)=\big\langle p_I(TM),[M]\big\rangle
=\big\langle f^*p_I(TM),[M]\big\rangle
=\big\langle p_I(TM),f_*[M]\big\rangle
=-P_I(M),
\]
forcing all Pontryagin numbers to vanish.
\end{proof}

\begin{corollary}
$\CP^{2n}$ admits no orientation reversing diffeomorphism. By contrast
$\CP^{2n+1}$ does, since its Pontryagin numbers live in the wrong dimension to
obstruct anything.
\end{corollary}

\begin{lemma}\label{lem:pontryagin-boundary}
If $M^{4n}$ is the boundary of a compact oriented $(4n+1)$-manifold $B$, then
all Pontryagin numbers of $M$ vanish.
\end{lemma}

\begin{proof}
\begin{strategy}
The same argument as Theorem~\ref{thm:boundary-numbers-vanish}, now with
integral coefficients, which is why orientations are needed: the normal bundle
of the boundary is canonically oriented by the outward direction.
\end{strategy}

\step{1} Along $M=\partial B$ the outward pointing direction trivializes the
normal bundle, so $TB|_M\iso TM\oplus\eps^1$ as oriented bundles and
$i^*p_j(TB)=p_j(TM)$ for the inclusion $i\colon M\hookrightarrow B$.

\step{2} With $\Z$ coefficients, $\partial[B]=[M]$ for the fundamental class
$[B]\in H_{4n+1}(B,M;\Z)$ of the compact oriented manifold with boundary.

\step{3} For $v=p_{i_1}(TB)\cdots p_{i_r}(TB)\in H^{4n}(B;\Z)$,
\[
P_I(M)=\big\langle i^*v,[M]\big\rangle
=\big\langle i^*v,\partial[B]\big\rangle
=\big\langle\delta\,i^*v,[B]\big\rangle=0,
\]
since $\delta\circ i^*=0$ by exactness of the sequence of the pair $(B,M)$.
\end{proof}

\begin{example}
$\CP^{2n}$ is not an oriented boundary. On the other hand
$\CP^{2n}\sqcup(-\CP^{2n})=\partial\big(\CP^{2n}\times[0,1]\big)$, so in the
oriented cobordism group $[\CP^{2n}]+[-\CP^{2n}]=0$, while
$2[\CP^{2n}]\neq0$: by Lemma~\ref{lem:no-orientation-reversal} the two
orientations of $\CP^{2n}$ really are different.
\end{example}
\section[The symmetric function classes SI]{\texorpdfstring{The symmetric function classes $S_I$}{The symmetric function classes SI}}


\begin{guiding}
Chern numbers of a product mix the classes of the two factors in a complicated
way. Is there a family of characteristic classes adapted to sums and products,
one that vanishes on non-trivial products?
\end{guiding}

For partitions $J=(j_1,\dots,j_s)$ and $K=(k_1,\dots,k_t)$ write $JK$ for the
concatenated partition $(j_1,\dots,j_s,k_1,\dots,k_t)$.

\begin{definition}
For a partition $I=(i_1,\dots,i_r)$ let
\[
S_I(u_1,\dots,u_n)=\sum u_{\alpha_1}^{i_1}\cdots u_{\alpha_r}^{i_r}
\]
be the smallest symmetric polynomial containing the monomial
$u_1^{i_1}\cdots u_r^{i_r}$, each distinct monomial appearing exactly once.
Being symmetric, it is a polynomial in $\sigma_1,\dots,\sigma_n$, and for a
complex bundle $E$ we set
\[
S_I(E):=S_I\big(c_1(E),\dots,c_n(E)\big)\in H^{2|I|}(B).
\]
Equivalently $S_I(E)=f^*S_I(\tilde\sigma_1,\dots,\tilde\sigma_n)$ for a
classifying map $f$. For instance $S_{(1)}=\sigma_1$,
$S_{(1,1)}=\sigma_2$, and $S_{(2)}=\sigma_1^2-2\sigma_2$, so
$S_{(2)}(E)=c_1^2-2c_2$.
\end{definition}

\begin{theorem}\label{thm:SI-sum}
For complex bundles $E,E'$ over the same base,
\[
S_I\big(E\oplus E'\big)=\sum_{JK=I}S_J(E)\smile S_K(E').
\]
\end{theorem}

\begin{proof}
\begin{strategy}
Verify the identity for sums of line bundles, where the Chern roots of
$E\oplus E'$ are the union of those of $E$ and $E'$; the combinatorics is the
distribution of the exponents of $I$ between the two groups of roots. The
splitting principle then gives the general case.
\end{strategy}

\step{1} \emph{Reduction to roots.} By the splitting principle
(Remark~\ref{rem:complex-splitting}) there is a map $f$, injective on cohomology,
such that $f^*E$ and $f^*E'$ split into line bundles, with Chern roots
$u'_1,\dots,u'_n$ and $u''_1,\dots,u''_{n'}$. The Chern roots of the sum are all
of these together, so it suffices to prove the corresponding polynomial identity.

\step{2} \emph{Distributing the exponents.} By definition, $S_I$ evaluated on
the combined variables is the sum of all distinct monomials
$v_{\alpha_1}^{i_1}\cdots v_{\alpha_r}^{i_r}$, where each $v_{\alpha}$ is one of
the $u'$ or one of the $u''$ and the indices are distinct. Group the terms
according to which exponents are assigned to primed variables: this is exactly a
splitting of the partition $I$ into a sub-partition $J$ used on the $u'$ and the
complementary $K$ used on the $u''$. Summing over all such splittings gives
\[
S_I(u',u'')=\sum_{JK=I}S_J(u')\,S_K(u''),
\]
each monomial being counted once on both sides.

\step{3} \emph{Back to bundles.} Applying $f^*$ and using its injectivity gives
the identity for $E,E'$.
\end{proof}

\begin{corollary}\label{cor:Sn-additive}
For a partition with a single part, $S_{(m)}(E\oplus E')=S_{(m)}(E)+S_{(m)}(E')$,
since the only splittings of $(m)$ are $J=(m),K=\emptyset$ and the reverse.
\end{corollary}

\begin{definition}
For a closed complex manifold $X$ of complex dimension $m$ put
\[
S_m[X]:=\big\langle S_{(m)}\big(c(TX)\big),\ [X]\big\rangle\in\Z .
\]
\end{definition}

\begin{corollary}\label{cor:Sn-product-vanishes}
If $K^m$ and $L^{n}$ are closed complex manifolds of positive dimension then
$S_{m+n}\big[K\times L\big]=0$.
\end{corollary}

\begin{proof}
Since $T(K\times L)=p_1^*TK\oplus p_2^*TL$,
Corollary~\ref{cor:Sn-additive} gives
\[
S_{(m+n)}\big(T(K\times L)\big)
=p_1^*S_{(m+n)}(TK)+p_2^*S_{(m+n)}(TL).
\]
The first summand lies in the image of $H^{2(m+n)}(K)$, which vanishes because
$K$ has complex dimension $m<m+n$; likewise for the second. Hence the class
itself is zero.
\end{proof}

\begin{example}
$S_n[\CP^n]=n+1\neq0$, because $T\CP^n\oplus\eps^1$ has $n+1$ Chern roots all
equal to $a$, so $S_{(n)}=\sum a^n=(n+1)a^n$. By
Corollary~\ref{cor:Sn-product-vanishes}, $\CP^n$ is not a non-trivial product of
complex manifolds.
\end{example}

\begin{theorem}\label{thm:matrix-nondegenerate}
Let $X^1,\dots,X^n$ be closed complex manifolds with $\dim_\C X^k=k$ and
$S_k[X^k]\neq0$ for every $k$. Then the $p(n)\times p(n)$ matrix of Chern
numbers
\[
\Big(\,C_I\big(X^{j_1}\times\cdots\times X^{j_s}\big)\,\Big)_{I,\ J\vdash n}
\]
is non-degenerate, where $p(n)$ is the number of partitions of $n$ and
$J=(j_1,\dots,j_s)$ runs over them.
\end{theorem}

The proof is a triangularity argument with respect to a suitable order on
partitions; it opens the final lecture, where the same result in its Pontryagin
form feeds directly into the computation of the oriented cobordism ring.

\lecture

\begin{guiding}
Everything converges here. The $S_I$ matrix theorem shows that products of
projective spaces are independent in cobordism; the oriented cobordism ring
organizes the question; and Thom's theorem converts cobordism into homotopy
theory. Together they compute $\Omega_*\otimes\Q$.
\end{guiding}

\section{Proof of the matrix theorem}

\begin{proof}[Proof of Theorem~\ref{thm:matrix-nondegenerate}]
\begin{strategy}
Order partitions by their number of parts. The multiplicativity of $S_I$ shows
that $S_I$ can be non-zero on a product indexed by $J$ only when $I$ has at
least as many parts as $J$, and that on the diagonal $I=J$ it is a product of
non-zero numbers. That makes the matrix of the $S_I$ triangular with non-zero
diagonal; a change of basis converts it into the matrix of Chern numbers.
\end{strategy}

\step{1} \emph{Multiplicativity.} Iterating Theorem~\ref{thm:SI-sum} across a
product $X^{j_1}\times\dots\times X^{j_s}$, whose tangent bundle is the sum of
the pull-backs of the $TX^{j_t}$,
\[
S_I\big(X^{j_1}\times\cdots\times X^{j_s}\big)
=\sum_{I_1I_2\cdots I_s=I}
S_{I_1}\big(X^{j_1}\big)\cdots S_{I_s}\big(X^{j_s}\big),
\]
the sum being over all ways of splitting $I$ into $s$ consecutive
sub-partitions.

\step{2} \emph{Which terms survive evaluation.} Evaluating on the fundamental
class, a term vanishes unless each $I_t$ has weight $j_t$, that is, unless
$I_t$ is a partition of $j_t$: otherwise the corresponding factor lies in a
cohomology group above or below the top degree of $X^{j_t}$. In particular each
$I_t$ is non-empty, so
\[
S_I\big[X^{j_1}\times\cdots\times X^{j_s}\big]\neq0
\quad\Longrightarrow\quad
\#\{\text{parts of }I\}\ \ge\ \#\{\text{parts of }J\}.
\]

\step{3} \emph{The diagonal.} If $I=J$ as partitions, the only splittings that
contribute are those with $I_t=(j_t)$ a single part, so
\[
S_J\big[X^{j_1}\times\cdots\times X^{j_s}\big]
=\lambda_J\prod_{t=1}^{s}S_{j_t}\big[X^{j_t}\big],
\]
where $\lambda_J$ is a positive integer counting the number of such splittings.
By hypothesis every factor is non-zero, so the diagonal entry is non-zero.

\step{4} \emph{Triangularity.} Order the $p(n)$ partitions of $n$ so that the
number of parts is non-decreasing. By Step 2 the matrix
$\big(S_I[X^J]\big)_{I,J}$ vanishes below the diagonal blocks and by Step 3 has
non-zero entries on the diagonal, so its determinant is non-zero.

\step{5} \emph{From $S_I$ to Chern numbers.} Each $S_I$ is an integral linear
combination of the monomials $c_{i_1}\cdots c_{i_r}$ of the same weight, and
conversely, because both families are bases of the same lattice, as shown in the
next paragraph. Hence the matrix of Chern numbers is obtained from the matrix of
the $S_I$ by an invertible change of basis and is therefore also
non-degenerate.
\end{proof}

\medskip\noindent\emph{The change of basis.} Inside $\Z[u_1,\dots,u_n]$ the
symmetric polynomials form $\Z[\sigma_1,\dots,\sigma_n]$. Its part of degree $k$
has two natural bases indexed by the partitions $\lambda\vdash k$ into at most
$n$ parts:

\begin{itemize}
\item the monomials in the elementary symmetric functions,
$\sigma_\lambda=\sigma_{\lambda_1}\sigma_{\lambda_2}\cdots$, equivalently
$\sigma_1^{r_1}\cdots\sigma_n^{r_n}$ with $\sum i\,r_i=k$;
\item the symmetrized monomials $S_\lambda$, each distinct monomial of exponent
pattern $\lambda$ appearing once.
\end{itemize}

The first is a basis by the fundamental theorem of symmetric polynomials; the
second because every symmetric polynomial is by definition a sum of
symmetrized monomials. The transition matrix between two bases of the same free
module is invertible over $\Z$, which is what Step 5 used.

\begin{remark}
When applying Theorem~\ref{thm:SI-sum} it matters that each monomial is counted
once. Writing $S_I$ as a sum over all of $S_n$ would count each monomial
$|\mathrm{Iso}|$ times, where $\mathrm{Iso}$ is the isotropy group of the
exponent pattern; the multiplicativity formula holds for the normalized version.
\end{remark}

\medskip\noindent\emph{The Pontryagin version.} For a closed oriented $M^{4m}$
and $I\vdash m$ define $S_I(p(E))$ by the same universal polynomial in the
Pontryagin classes and put $S_I[M]=\langle S_I(p(TM)),[M]\rangle$. Since the
Pontryagin classes are multiplicative only modulo $2$-torsion
(Theorem~\ref{thm:pontryagin-whitney}),
\[
S_I\big(p(E\oplus E')\big)=\sum_{JK=I}S_J\big(p(E)\big)\,S_K\big(p(E')\big)
\qquad\text{modulo $2$-torsion,}
\]
which is enough for evaluation against fundamental classes, torsion classes
pairing to zero with an integral class after multiplication by $2$.

\begin{theorem}\label{thm:pontryagin-matrix}
Let $\{M^{4k}\}_{k\ge1}$ be closed oriented manifolds with $S_k[M^{4k}]\neq0$.
Then for every $n$ the $p(n)\times p(n)$ matrix of Pontryagin numbers
\[
\Big(P_I\big(M^{4j_1}\times\cdots\times M^{4j_s}\big)\Big)_{I,J\vdash n}
\]
is non-degenerate. The family $M^{4k}=\CP^{2k}$ qualifies, since
$p(T\CP^{2k})=(1+u^2)^{2k+1}$ gives $S_k[\CP^{2k}]=2k+1\neq0$.
\end{theorem}

\begin{proof}
The five steps of the proof above apply verbatim, with $c$ replaced by $p$ and
degrees multiplied by two; the only change is that the identities hold modulo
$2$-torsion, which does not affect the evaluations.
\end{proof}

\begin{remark}[quoted]
The classifying space for \emph{oriented} bundles is the oriented Grassmannian
$\tilde G_n$ of oriented $n$-planes, and
\[
H^*\big(\tilde G_{2m+1};\Q\big)=\Q[p_1,\dots,p_m],
\]
\[
H^*\big(\tilde G_{2m};\Q\big)=\Q[p_1,\dots,p_{m-1},e],
\qquad e^2=p_m .
\]
So rationally the Pontryagin classes, together with the Euler class in even
rank, are all the characteristic classes of oriented bundles.
\end{remark}

\section{The oriented cobordism ring}

\begin{definition}
Closed oriented manifolds $M,M'$ of dimension $n$ are \emph{oriented cobordant}
if there is a compact oriented $X$ together with an orientation preserving
diffeomorphism $M\sqcup(-M')\iso\partial X$, the boundary carrying its induced
orientation. Let $\Omega_n$ be the set of equivalence classes.
\end{definition}

\begin{proposition}
$\Omega_n$ is an abelian group under disjoint union, with zero the class of
boundaries and with $-[M]=[-M]$; and
$\Omega_*=\bigoplus_{n\ge0}\Omega_n$ is a graded commutative ring under
Cartesian product.
\end{proposition}

\begin{proof}
Reflexivity uses $M\sqcup(-M)=\partial(M\times[0,1])$; symmetry is reversal of
orientation of the cobordism; transitivity is gluing two cobordisms along a
common boundary component, which is again a smooth manifold after smoothing the
corner. Associativity and commutativity of $\sqcup$ and $\times$ are clear, and
the product of a boundary with a closed manifold is a boundary, so the product
descends to classes.
\end{proof}

\begin{lemma}
Oriented cobordant manifolds have equal Pontryagin numbers.
\end{lemma}

\begin{proof}
If $M\sqcup(-M')$ bounds, Lemma~\ref{lem:pontryagin-boundary} gives
$P_I(M)+P_I(-M')=0$, and $P_I(-M')=-P_I(M')$.
\end{proof}

\begin{corollary}\label{cor:rank-omega}
For every partition $I$ of $k$ the assignment $[M]\mapsto P_I(M)$ is a group
homomorphism $\Omega_{4k}\to\Z$. The products
$\CP^{2i_1}\times\dots\times\CP^{2i_r}$, as $I=(i_1,\dots,i_r)$ runs over the
partitions of $k$, are linearly independent in $\Omega_{4k}$, so
\[
\rk\,\Omega_{4k}\ \ge\ p(k).
\]
\end{corollary}

\begin{proof}
Assemble the homomorphisms into
$\Omega_{4k}\to\Z^{p(k)}$, $[M]\mapsto(P_I(M))_I$. By
Theorem~\ref{thm:pontryagin-matrix} the images of the $p(k)$ products
$\CP^{2i_1}\times\dots\times\CP^{2i_r}$ form a non-degenerate matrix, hence are
linearly independent over $\Q$; so the classes themselves are independent.
\end{proof}

\section{Thom spaces}

\begin{definition}
Let $\pi\colon E\to B$ be a real bundle of rank $k$ with a Euclidean metric, and
$E^{\ge1}=\{e\mid |e|\ge1\}$. The \emph{Thom space} is the pointed space
\[
T(E)=E\big/E^{\ge1},
\qquad
*_E=\big[E^{\ge1}\big],
\]
that is, the unit disc bundle with its boundary sphere bundle collapsed to a
point.
\end{definition}

\begin{center}
\begin{tikzpicture}[scale=0.9]
\draw[thick] (0,0) ellipse (1.1 and 0.34);
\draw[thick] (0,1.5) ellipse (1.1 and 0.34);
\draw[thick] (-1.1,0) -- (-1.1,1.5);
\draw[thick] (1.1,0) -- (1.1,1.5);
\node at (0,-0.85) {\small disc bundle over $B$};
\draw[-{Stealth},thick] (1.9,0.75) -- (3.1,0.75);
\fill (4.6,0.75) circle (1.7pt);
\draw[thick] (4.6,0.75) .. controls (4.05,1.75) and (5.8,1.75) .. (4.6,0.75);
\draw[thick] (4.6,0.75) .. controls (4.05,-0.25) and (5.8,-0.25) .. (4.6,0.75);
\node at (4.7,-0.85) {\small $T(E)$};
\end{tikzpicture}
\end{center}

\begin{proposition}\label{prop:thom-space}
$\tilde H^{i+k}\big(T(E)\big)\iso H^i(B)$ by the Thom isomorphism, for oriented
$E$ over $\Z$. If $B$ is a CW complex then $T(E)$ is a CW complex with one
$(d+k)$-cell for each $d$-cell of $B$, plus the basepoint; in particular
$T(E)$ is $(k-1)$-connected.
\end{proposition}

\begin{proof}
Since $E^{\ge1}$ is a deformation retract of $E_0$, we have
\[
\tilde H^*\big(T(E)\big)\iso H^*\big(E,E^{\ge1}\big)\iso H^*(E,E_0)
\iso H^{*-k}(B)
\]
by Theorem~\ref{thm:thom}. For the cell structure, over a $d$-cell $D^d$ the
bundle is trivial, so the part of the disc bundle above it is
$D^d\times D^k$, contributing one $(d+k)$-cell after the collapse. Since
$T(E)$ then has no cells in dimensions $1,\dots,k-1$ apart from the basepoint, a
cellular approximation argument shows every map $S^m\to T(E)$ with $m<k$ is
null-homotopic.
\end{proof}

\section{Thom's theorem}

\begin{definition}
Let $\mathcal C$ be the class of finite abelian groups. A homomorphism
$h\colon A\to B$ is a \emph{$\mathcal C$-isomorphism} if $\ker h$ and
$\mathrm{coker}\,h$ both lie in $\mathcal C$.
\end{definition}

\begin{theorem}[Hurewicz modulo $\mathcal C$, quoted]
If $X$ is a finite CW complex which is $(k-1)$-connected, then the Hurewicz map
$\pi_r(X)\to H_r(X;\Z)$ is a $\mathcal C$-isomorphism for $r<2k-1$.
\end{theorem}

\begin{corollary}\label{cor:hurewicz-thom}
For an oriented rank $k$ bundle $E$ over a finite CW complex $B$, and $n<k-1$,
the composite
\[
\pi_{n+k}\big(T(E)\big)\longrightarrow H_{n+k}\big(T(E);\Z\big)\iso H_n(B;\Z)
\]
is a $\mathcal C$-isomorphism.
\end{corollary}

\begin{theorem}[Thom]\label{thm:thom-cobordism}
For $k>n+1$ there is a canonical isomorphism
\[
\pi_{n+k}\big(T(\tilde\gamma^k)\big)\ \iso\ \Omega_n,
\]
where $\tilde\gamma^k$ is the universal oriented $k$-plane bundle over
$\tilde G_k$.
\end{theorem}

\begin{proof}[Proof of surjectivity]
\begin{strategy}
Given a closed oriented $M^n$, embed it in a sphere, and collapse everything
outside a tubular neighbourhood; what remains is a map from the sphere to the
Thom space of the normal bundle, which maps to the Thom space of the universal
bundle. The preimage of the zero section is $M$ by construction.
\end{strategy}

\step{1} \emph{Embedding.} By Whitney, embed $M^n\subset\R^{n+k}\subset S^{n+k}$
and let $E$ be its normal bundle, oriented by
$TM\oplus E=T\R^{n+k}|_M$.

\step{2} \emph{The collapse map.} Let $U\iso E$ be a tubular neighbourhood
(Theorem~\ref{thm:tnt}). Collapsing the complement of $U$ gives
\[
S^{n+k}\longrightarrow S^{n+k}\big/\big(S^{n+k}\smallsetminus U\big)\iso T(E),
\]
a continuous map sending everything outside the tube to the basepoint.

\step{3} \emph{To the universal Thom space.} A bundle map $E\to\tilde\gamma^k$,
which exists by the oriented version of Theorem~\ref{thm:universal}, induces a
pointed map $T(E)\to T(\tilde\gamma^k)$. Composing with Step 2 gives
$g\colon S^{n+k}\to T(\tilde\gamma^k)$.

\step{4} \emph{Recovering $M$.} By construction $g$ is smooth and transverse to
the zero section near the preimage of the zero section, and
$g^{-1}(\text{zero section})=M$ with normal bundle $g^*\tilde\gamma^k\iso E$.
Hence the class of $g$ maps to $[M]$.
\end{proof}

\begin{remark}[quoted]
For injectivity, and to define the inverse map, one uses transversality: every
continuous $g\colon S^{n+k}\to T(E)$ is homotopic to one which is smooth and
transverse to the zero section on the preimage of $E^{\le1}$, and then
$g^{-1}(\text{zero section})$ is a closed oriented $n$-manifold whose cobordism
class depends only on the homotopy class of $g$, a homotopy made transverse
providing the cobordism. We do not prove the transversality statements.
\end{remark}

\begin{corollary}
For $k$ large,
\[
\Omega_n\otimes\Q\ \iso\ \pi_{n+k}\big(T(\tilde\gamma^k)\big)\otimes\Q
\ \iso\ H_n\big(\tilde G_k;\Q\big)
=\begin{cases}
\Q^{\,p(m)} & n=4m,\\
0 & \text{otherwise,}
\end{cases}
\]
using Corollary~\ref{cor:hurewicz-thom}, which becomes an isomorphism after
tensoring with $\Q$ since finite groups die, together with the rational
cohomology of $\tilde G_k$ quoted above. Combining with
Corollary~\ref{cor:rank-omega}, where the lower bound $p(m)$ is realized by
products of even complex projective spaces,
\[
\Omega_*\otimes\Q\ =\ \Q\big[\,[\CP^2],\,[\CP^4],\,[\CP^6],\dots\,\big].
\]
\end{corollary}

\begin{remark}[where this leads]
Everything in the course feeds forward. Hirzebruch's signature theorem expresses
the signature of $M^{4k}$ as a universal polynomial in the Pontryagin classes,
precisely because both sides are cobordism invariants and agree on the
generators $\CP^{2i}$. Milnor's exotic spheres are detected by the failure of an
identity between Pontryagin numbers. And index theory reinterprets all of these
numbers analytically.
\end{remark}


\end{document}
